\newlabel{tab:front-notation}{{1}{ii}{记号（测试）\relax }{table.caption.3}{}}
\newlabel{tab:front-notation@cref}{{[table][1][0]1}{[i][2][]ii}}
\newlabel{ch:intro}{{1}{2}{引言：部分子图像、因子化的含义与严格性分级}{chapter.1}{}}
\newlabel{ch:intro@cref}{{[chapter][1][]1}{[1][2][]2}}
\newlabel{sec:intro-parton}{{1.1}{2}{第 1 节（测试）}{section.1.1}{}}
\newlabel{sec:intro-parton@cref}{{[section][1][1]1.1}{[1][2][]2}}
\newlabel{def:intro-hard-process}{{1.1}{2}{硬过程与幂次展开}{theorem.1.1}{}}
\newlabel{def:intro-hard-process@cref}{{[definition][1][1]1.1}{[1][2][]2}}
\newlabel{def:intro-factorization}{{1.2}{2}{共线因子化陈述}{theorem.1.2}{}}
\newlabel{def:intro-factorization@cref}{{[definition][2][1]1.2}{[1][2][]2}}
\newlabel{sec:intro-statement}{{1.2}{2}{第 2 节（测试）}{section.1.2}{}}
\newlabel{sec:intro-statement@cref}{{[section][2][1]1.2}{[1][2][]2}}
\newlabel{rem:intro-incoherence}{{1.3}{2}{部分子模型的非相干性论证（启发式）}{theorem.1.3}{}}
\newlabel{rem:intro-incoherence@cref}{{[remark][3][1]1.3}{[1][2][]2}}
\newlabel{ex:intro-handbag}{{1.4}{2}{树图手提包与 Callan–Gross 关系}{theorem.1.4}{}}
\newlabel{ex:intro-handbag@cref}{{[example][4][1]1.4}{[1][2][]2}}
\newlabel{sec:intro-proof}{{1.3}{2}{第 3 节（测试）}{section.1.3}{}}
\newlabel{sec:intro-proof@cref}{{[section][3][1]1.3}{[1][2][]2}}
\newlabel{rem:intro-sharpness}{{1.5}{2}{"每一阶、领头幂次"的精确含义与幂次的锐利性}{theorem.1.5}{}}
\newlabel{rem:intro-sharpness@cref}{{[remark][5][1]1.5}{[1][2][]2}}
\newlabel{def:intro-tags}{{1.6}{2}{严格性标签、限定语与等级传递}{theorem.1.6}{}}
\newlabel{def:intro-tags@cref}{{[definition][6][1]1.6}{[1][2][]2}}
\newlabel{sec:intro-tags}{{1.4}{2}{第 4 节（测试）}{section.1.4}{}}
\newlabel{sec:intro-tags@cref}{{[section][4][1]1.4}{[1][2][]2}}
\newlabel{def:intro-error-types}{{1.7}{2}{误差类型}{theorem.1.7}{}}
\newlabel{def:intro-error-types@cref}{{[definition][7][1]1.7}{[1][2][]2}}
\newlabel{rem:intro-grade-dictionary}{{1.8}{2}{与文献中另一种分级的对照}{theorem.1.8}{}}
\newlabel{rem:intro-grade-dictionary@cref}{{[remark][8][1]1.8}{[1][2][]2}}
\newlabel{sec:intro-kinds}{{1.5}{3}{第 5 节（测试）}{section.1.5}{}}
\newlabel{sec:intro-kinds@cref}{{[section][5][1]1.5}{[1][2][]3}}
\newlabel{asm:intro-wightman}{{1.9}{3}{QCD 的 Wightman 结构 \tagC }{theorem.1.9}{}}
\newlabel{asm:intro-wightman@cref}{{[assumption][9][1]1.9}{[1][2][]3}}
\newlabel{asm:intro-composite}{{1.10}{3}{规范不变复合算符与光线算符的存在与重整化群 \tagC }{theorem.1.10}{}}
\newlabel{asm:intro-composite@cref}{{[assumption][10][1]1.10}{[1][3][]3}}
\newlabel{sec:intro-overview}{{1.6}{3}{第 6 节（测试）}{section.1.6}{}}
\newlabel{sec:intro-overview@cref}{{[section][6][1]1.6}{[1][3][]3}}
\newlabel{asm:intro-np-ope}{{1.11}{3}{非微扰算符乘积展开 \tagC }{theorem.1.11}{}}
\newlabel{asm:intro-np-ope@cref}{{[assumption][11][1]1.11}{[1][3][]3}}
\newlabel{asm:intro-af-coeff}{{1.12}{3}{系数的渐近自由行为 \tagC }{theorem.1.12}{}}
\newlabel{asm:intro-af-coeff@cref}{{[assumption][12][1]1.12}{[1][3][]3}}
\newlabel{sec:intro-misconceptions}{{1.7}{4}{第 7 节（测试）}{section.1.7}{}}
\newlabel{sec:intro-misconceptions@cref}{{[section][7][1]1.7}{[1][4][]4}}
\newlabel{asm:intro-growth}{{1.13}{4}{前向振幅的增长条件 \tagC }{theorem.1.13}{}}
\newlabel{asm:intro-growth@cref}{{[assumption][13][1]1.13}{[1][4][]4}}
\newlabel{asm:intro-bridge}{{1.14}{4}{从部分子微扰论到强子的桥接假设 \tagC }{theorem.1.14}{}}
\newlabel{asm:intro-bridge@cref}{{[assumption][14][1]1.14}{[1][4][]4}}
\newlabel{sec:intro-organisation}{{1.8}{4}{第 8 节（测试）}{section.1.8}{}}
\newlabel{sec:intro-organisation@cref}{{[section][8][1]1.8}{[1][4][]4}}
\newlabel{def:intro-MA}{{1.15}{4}{元假设 MA0–MA4（总称）与编号假设的对应}{theorem.1.15}{}}
\newlabel{def:intro-MA@cref}{{[definition][15][1]1.15}{[1][4][]4}}
\newlabel{ch:qcd}{{2}{5}{QCD 基础——从规范不变性到渐近自由}{chapter.2}{}}
\newlabel{ch:qcd@cref}{{[chapter][2][]2}{[1][5][]5}}
\newlabel{sec:qcd-why}{{2.1}{5}{第 1 节（测试）}{section.2.1}{}}
\newlabel{sec:qcd-why@cref}{{[section][1][2]2.1}{[1][5][]5}}
\newlabel{def:qcd-gauge-group}{{2.1}{5}{规范群 SU($N$) 及其表示}{theorem.2.1}{}}
\newlabel{def:qcd-gauge-group@cref}{{[definition][1][2]2.1}{[1][5][]5}}
\newlabel{lem:qcd-colour-identities}{{2.2}{5}{色代数恒等式}{theorem.2.2}{}}
\newlabel{lem:qcd-colour-identities@cref}{{[lemma][2][2]2.2}{[1][5][]5}}
\newlabel{ex:qcd-colour-algebra}{{2.3}{5}{单圈图的色因子与双线记号}{theorem.2.3}{}}
\newlabel{ex:qcd-colour-algebra@cref}{{[example][3][2]2.3}{[1][5][]5}}
\newlabel{prop:qcd-covariance}{{2.4}{5}{协变导数与场强的规范协变性 \tagR }{theorem.2.4}{}}
\newlabel{prop:qcd-covariance@cref}{{[proposition][4][2]2.4}{[1][5][]5}}
\newlabel{sec:qcd-gauge}{{2.2}{5}{第 2 节（测试）}{section.2.2}{}}
\newlabel{sec:qcd-gauge@cref}{{[section][2][2]2.2}{[1][5][]5}}
\newlabel{def:prelim-qcd}{{2.5}{5}{QCD 拉氏量与规范变换}{theorem.2.5}{}}
\newlabel{def:prelim-qcd@cref}{{[definition][5][2]2.5}{[1][5][]5}}
\newlabel{thm:qcd-gauge-invariance}{{2.6}{5}{经典 QCD 拉氏量的规范不变性与唯一性 \tagR }{theorem.2.6}{}}
\newlabel{thm:qcd-gauge-invariance@cref}{{[theorem][6][2]2.6}{[1][5][]5}}
\newlabel{rem:qcd-theta}{{2.7}{5}{$\theta $ 项}{theorem.2.7}{}}
\newlabel{rem:qcd-theta@cref}{{[remark][7][2]2.7}{[1][5][]5}}
\newlabel{rem:qcd-fp-derivation}{{2.8}{5}{Faddeev–Popov 构造（形式推导）}{theorem.2.8}{}}
\newlabel{rem:qcd-fp-derivation@cref}{{[remark][8][2]2.8}{[1][5][]5}}
\newlabel{sec:prelim-qcd}{{2.3}{5}{第 3 节（测试）}{section.2.3}{}}
\newlabel{sec:prelim-qcd@cref}{{[section][3][2]2.3}{[1][5][]5}}
\newlabel{def:qcd-fp}{{2.9}{5}{协变规范、Faddeev–Popov 鬼场与微扰 QCD}{theorem.2.9}{}}
\newlabel{def:qcd-fp@cref}{{[definition][9][2]2.9}{[1][5][]5}}
\newlabel{ex:qcd-fp-lowest-order}{{2.10}{6}{由拉氏量导出传播子、鬼场顶点与闭合鬼圈因子}{theorem.2.10}{}}
\newlabel{ex:qcd-fp-lowest-order@cref}{{[example][10][2]2.10}{[1][5][]6}}
\newlabel{def:prelim-feynman-rules}{{2.11}{6}{Feynman 规则}{theorem.2.11}{}}
\newlabel{def:prelim-feynman-rules@cref}{{[definition][11][2]2.11}{[1][6][]6}}
\newlabel{ex:qcd-qqbar-gg}{{2.12}{6}{树图 $q\bar q\to gg$：为什么需要鬼场}{theorem.2.12}{}}
\newlabel{ex:qcd-qqbar-gg@cref}{{[example][12][2]2.12}{[1][6][]6}}
\newlabel{sec:qcd-quantization}{{2.4}{6}{第 4 节（测试）}{section.2.4}{}}
\newlabel{sec:qcd-quantization@cref}{{[section][4][2]2.4}{[1][6][]6}}
\newlabel{def:qcd-dimreg}{{2.13}{6}{维数正规化}{theorem.2.13}{}}
\newlabel{def:qcd-dimreg@cref}{{[definition][13][2]2.13}{[1][6][]6}}
\newlabel{prop:qcd-superficial-degree}{{2.14}{6}{紫外超表观发散度 \tagR }{theorem.2.14}{}}
\newlabel{prop:qcd-superficial-degree@cref}{{[proposition][14][2]2.14}{[1][6][]6}}
\newlabel{def:prelim-brst}{{2.15}{6}{BRST 变换}{theorem.2.15}{}}
\newlabel{def:prelim-brst@cref}{{[definition][15][2]2.15}{[1][6][]6}}
\newlabel{lem:prelim-brst-nilpotent}{{2.16}{6}{$s^2=0$ 与拉氏量的 BRST 不变性 \tagR }{theorem.2.16}{}}
\newlabel{lem:prelim-brst-nilpotent@cref}{{[lemma][16][2]2.16}{[1][6][]6}}
\newlabel{sec:qcd-regularization}{{2.5}{6}{第 5 节（测试）}{section.2.5}{}}
\newlabel{sec:qcd-regularization@cref}{{[section][5][2]2.5}{[1][6][]6}}
\newlabel{prop:qcd-brst-cohomology}{{2.17}{6}{自由渐近 Fock 空间上的 BRST 上同调（Kugo–Ojima 四重态） \tagR }{theorem.2.17}{}}
\newlabel{prop:qcd-brst-cohomology@cref}{{[proposition][17][2]2.17}{[1][6][]6}}
\newlabel{ex:qcd-brst-quartet}{{2.18}{6}{单一动量模式的四重态}{theorem.2.18}{}}
\newlabel{ex:qcd-brst-quartet@cref}{{[example][18][2]2.18}{[1][6][]6}}
\newlabel{lem:qcd-st-regularized}{{2.19}{6}{维数正规化理论的 Slavnov–Taylor 恒等式 \tagR }{theorem.2.19}{}}
\newlabel{lem:qcd-st-regularized@cref}{{[lemma][19][2]2.19}{[1][6][]6}}
\newlabel{def:qcd-renormalization}{{2.20}{6}{重整化常数与抵消项}{theorem.2.20}{}}
\newlabel{def:qcd-renormalization@cref}{{[definition][20][2]2.20}{[1][6][]6}}
\newlabel{sec:prelim-brst}{{2.6}{6}{第 6 节（测试）}{section.2.6}{}}
\newlabel{sec:prelim-brst@cref}{{[section][6][2]2.6}{[1][6][]6}}
\newlabel{def:prelim-msbar}{{2.21}{6}{$\overline {\rm MS}$ 方案}{theorem.2.21}{}}
\newlabel{def:prelim-msbar@cref}{{[definition][21][2]2.21}{[1][6][]6}}
\newlabel{thm:qcd-renormalizability}{{2.22}{7}{QCD 的可重整性 \tagR }{theorem.2.22}{}}
\newlabel{thm:qcd-renormalizability@cref}{{[theorem][22][2]2.22}{[1][6][]7}}
\newlabel{prop:qcd-st-relations}{{2.23}{7}{重整化常数间的 Slavnov–Taylor 关系 \tagR }{theorem.2.23}{}}
\newlabel{prop:qcd-st-relations@cref}{{[proposition][23][2]2.23}{[1][7][]7}}
\newlabel{thm:prelim-slavnov-taylor}{{2.24}{7}{Slavnov–Taylor 恒等式（因子化所用形式） \tagR }{theorem.2.24}{}}
\newlabel{thm:prelim-slavnov-taylor@cref}{{[theorem][24][2]2.24}{[1][7][]7}}
\newlabel{sec:prelim-renorm}{{2.7}{7}{第 7 节（测试）}{section.2.7}{}}
\newlabel{sec:prelim-renorm@cref}{{[section][7][2]2.7}{[1][7][]7}}
\newlabel{prop:qcd-xi-independence}{{2.25}{7}{物理量的规范参数无关性 \tagR }{theorem.2.25}{}}
\newlabel{prop:qcd-xi-independence@cref}{{[proposition][25][2]2.25}{[1][7][]7}}
\newlabel{thm:prelim-composite-renorm}{{2.26}{7}{规范不变复合算符的重整化（Joglekar–Lee） \tagR }{theorem.2.26}{}}
\newlabel{thm:prelim-composite-renorm@cref}{{[theorem][26][2]2.26}{[1][7][]7}}
\newlabel{rem:qcd-bphz-brst}{{2.27}{7}{BPHZ 减除与 Slavnov–Taylor 恒等式；减除度的顶点可加性}{theorem.2.27}{}}
\newlabel{rem:qcd-bphz-brst@cref}{{[remark][27][2]2.27}{[1][7][]7}}
\newlabel{ex:qcd-gluon-self-energy}{{2.28}{7}{单圈胶子自能与 $Z_3$}{theorem.2.28}{}}
\newlabel{ex:qcd-gluon-self-energy@cref}{{[example][28][2]2.28}{[1][7][]7}}
\newlabel{sec:qcd-renorm}{{2.8}{7}{第 8 节（测试）}{section.2.8}{}}
\newlabel{sec:qcd-renorm@cref}{{[section][8][2]2.8}{[1][7][]7}}
\newlabel{ex:qcd-quark-self-energy-vertex}{{2.29}{7}{夸克自能与夸克—胶子顶点}{theorem.2.29}{}}
\newlabel{ex:qcd-quark-self-energy-vertex@cref}{{[example][29][2]2.29}{[1][7][]7}}
\newlabel{ex:qcd-beta0-four-ways}{{2.30}{7}{用四种 Slavnov–Taylor 等价途径计算 $\beta _0$}{theorem.2.30}{}}
\newlabel{ex:qcd-beta0-four-ways@cref}{{[example][30][2]2.30}{[1][7][]7}}
\newlabel{prop:prelim-beta}{{2.31}{7}{$\beta $ 函数 \tagR }{theorem.2.31}{}}
\newlabel{prop:prelim-beta@cref}{{[proposition][31][2]2.31}{[1][7][]7}}
\newlabel{prop:qcd-rge}{{2.32}{7}{重整化群方程 \tagR }{theorem.2.32}{}}
\newlabel{prop:qcd-rge@cref}{{[proposition][32][2]2.32}{[1][7][]7}}
\newlabel{sec:qcd-running}{{2.9}{8}{第 9 节（测试）}{section.2.9}{}}
\newlabel{sec:qcd-running@cref}{{[section][9][2]2.9}{[1][7][]8}}
\newlabel{prop:prelim-rg}{{2.33}{8}{重整化群改进 \tagR }{theorem.2.33}{}}
\newlabel{prop:prelim-rg@cref}{{[proposition][33][2]2.33}{[1][7][]8}}
\newlabel{def:qcd-lambda}{{2.34}{8}{$\Lambda _{\overline {\rm MS}}$ 与跑动耦合}{theorem.2.34}{}}
\newlabel{def:qcd-lambda@cref}{{[definition][34][2]2.34}{[1][8][]8}}
\newlabel{rem:qcd-decoupling}{{2.35}{8}{重夸克退耦与有效味数}{theorem.2.35}{}}
\newlabel{rem:qcd-decoupling@cref}{{[remark][35][2]2.35}{[1][8][]8}}
\newlabel{ex:qcd-running-numbers}{{2.36}{8}{从 $\alpha _s(M_Z)$ 出发的数值跑动}{theorem.2.36}{}}
\newlabel{ex:qcd-running-numbers@cref}{{[example][36][2]2.36}{[1][8][]8}}
\newlabel{sec:qcd-af}{{2.10}{8}{第 10 节（测试）}{section.2.10}{}}
\newlabel{sec:qcd-af@cref}{{[section][10][2]2.10}{[1][8][]8}}
\newlabel{thm:qcd-asymptotic-freedom}{{2.37}{8}{渐近自由 \tagR }{theorem.2.37}{}}
\newlabel{thm:qcd-asymptotic-freedom@cref}{{[theorem][37][2]2.37}{[1][8][]8}}
\newlabel{rem:qcd-antiscreening}{{2.38}{8}{屏蔽与反屏蔽的直观图像}{theorem.2.38}{}}
\newlabel{rem:qcd-antiscreening@cref}{{[remark][38][2]2.38}{[1][8][]8}}
\newlabel{rem:qcd-af-meaning}{{2.39}{8}{渐近自由意味着什么、不意味着什么}{theorem.2.39}{}}
\newlabel{rem:qcd-af-meaning@cref}{{[remark][39][2]2.39}{[1][8][]8}}
\newlabel{prop:qcd-soft-collinear-limits}{{2.40}{8}{树图层次的软极限与共线极限 \tagR }{theorem.2.40}{}}
\newlabel{prop:qcd-soft-collinear-limits@cref}{{[proposition][40][2]2.40}{[1][8][]8}}
\newlabel{sec:qcd-ir}{{2.11}{8}{第 11 节（测试）}{section.2.11}{}}
\newlabel{sec:qcd-ir@cref}{{[section][11][2]2.11}{[1][8][]8}}
\newlabel{def:qcd-ir-safety-informal}{{2.41}{8}{红外安全（非形式）}{theorem.2.41}{}}
\newlabel{def:qcd-ir-safety-informal@cref}{{[definition][41][2]2.41}{[1][8][]8}}
\newlabel{ex:qcd-ee-real-virtual}{{2.42}{8}{$e^+e^-\to q\bar q(g)$：实辐射与虚修正的相消}{theorem.2.42}{}}
\newlabel{ex:qcd-ee-real-virtual@cref}{{[example][42][2]2.42}{[1][8][]8}}
\newlabel{ex:qcd-sterman-weinberg}{{2.43}{8}{第一个红外安全的排他量：Sterman–Weinberg 两喷注截面}{theorem.2.43}{}}
\newlabel{ex:qcd-sterman-weinberg@cref}{{[example][43][2]2.43}{[1][8][]8}}
\newlabel{rem:qcd-ir-theorems-preview}{{2.44}{8}{红外定理预览}{theorem.2.44}{}}
\newlabel{rem:qcd-ir-theorems-preview@cref}{{[remark][44][2]2.44}{[1][8][]8}}
\newlabel{ch:prelim}{{3}{9}{因子化的数学与场论工具}{chapter.3}{}}
\newlabel{ch:prelim@cref}{{[chapter][3][]3}{[1][9][]9}}
\newlabel{sec:prelim-lightfront}{{3.1}{9}{第 1 节（测试）}{section.3.1}{}}
\newlabel{sec:prelim-lightfront@cref}{{[section][1][3]3.1}{[1][9][]9}}
\newlabel{def:prelim-lightfront}{{3.1}{9}{光锥坐标与 Sudakov 分解}{theorem.3.1}{}}
\newlabel{def:prelim-lightfront@cref}{{[definition][1][3]3.1}{[1][9][]9}}
\newlabel{lem:prelim-boost}{{3.2}{9}{推动下的标度 \tagR }{theorem.3.2}{}}
\newlabel{lem:prelim-boost@cref}{{[lemma][2][3]3.2}{[1][9][]9}}
\newlabel{ex:prelim-boost-scaling}{{3.3}{9}{共线动量与好/坏分量的标度}{theorem.3.3}{}}
\newlabel{ex:prelim-boost-scaling@cref}{{[example][3][3]3.3}{[1][9][]9}}
\newlabel{def:prelim-wilson}{{3.4}{9}{Wilson 线}{theorem.3.4}{}}
\newlabel{def:prelim-wilson@cref}{{[definition][4][3]3.4}{[1][9][]9}}
\newlabel{prop:prelim-wilson-covariance}{{3.5}{9}{Wilson 线的规范协变性 \tagR }{theorem.3.5}{}}
\newlabel{prop:prelim-wilson-covariance@cref}{{[proposition][5][3]3.5}{[1][9][]9}}
\newlabel{sec:prelim-wilson}{{3.2}{9}{第 2 节（测试）}{section.3.2}{}}
\newlabel{sec:prelim-wilson@cref}{{[section][2][3]3.2}{[1][9][]9}}
\newlabel{prop:prelim-wilson-rules}{{3.6}{9}{半无限直 Wilson 线的动量空间 Feynman 规则 \tagR }{theorem.3.6}{}}
\newlabel{prop:prelim-wilson-rules@cref}{{[proposition][6][3]3.6}{[1][9][]9}}
\newlabel{lem:prelim-eikonal-identity}{{3.7}{9}{程函恒等式}{theorem.3.7}{}}
\newlabel{lem:prelim-eikonal-identity@cref}{{[lemma][7][3]3.7}{[1][9][]9}}
\newlabel{ex:prelim-eikonal-two}{{3.8}{9}{$m=2$ 的程函恒等式与入射夸克传播子的程函极限}{theorem.3.8}{}}
\newlabel{ex:prelim-eikonal-two@cref}{{[example][8][3]3.8}{[1][9][]9}}
\newlabel{ex:prelim-wilson-segment}{{3.9}{9}{有限类光线段的单胶子项}{theorem.3.9}{}}
\newlabel{ex:prelim-wilson-segment@cref}{{[example][9][3]3.9}{[1][9][]9}}
\newlabel{def:prelim-dr-conventions}{{3.10}{9}{维数正规化的技术约定}{theorem.3.10}{}}
\newlabel{def:prelim-dr-conventions@cref}{{[definition][10][3]3.10}{[1][9][]9}}
\newlabel{sec:prelim-dimreg}{{3.3}{10}{第 3 节（测试）}{section.3.3}{}}
\newlabel{sec:prelim-dimreg@cref}{{[section][3][3]3.3}{[1][9][]10}}
\newlabel{lem:prelim-scaleless}{{3.11}{10}{无标度积分为零}{theorem.3.11}{}}
\newlabel{lem:prelim-scaleless@cref}{{[lemma][11][3]3.11}{[1][9][]10}}
\newlabel{ex:prelim-scaleless}{{3.12}{10}{无标度积分的紫外/红外记账}{theorem.3.12}{}}
\newlabel{ex:prelim-scaleless@cref}{{[example][12][3]3.12}{[1][10][]10}}
\newlabel{def:prelim-gamma5}{{3.13}{10}{$\gamma _5$ 与 $\varepsilon $ 张量}{theorem.3.13}{}}
\newlabel{def:prelim-gamma5@cref}{{[definition][13][3]3.13}{[1][10][]10}}
\newlabel{def:prelim-plus}{{3.14}{10}{加号分布}{theorem.3.14}{}}
\newlabel{def:prelim-plus@cref}{{[definition][14][3]3.14}{[1][10][]10}}
\newlabel{lem:prelim-plus-expansion}{{3.15}{10}{$(1-z)^{-1-a\varepsilon }$ 的分布展开}{theorem.3.15}{}}
\newlabel{lem:prelim-plus-expansion@cref}{{[lemma][15][3]3.15}{[1][10][]10}}
\newlabel{sec:prelim-distributions}{{3.4}{10}{第 4 节（测试）}{section.3.4}{}}
\newlabel{sec:prelim-distributions@cref}{{[section][4][3]3.4}{[1][10][]10}}
\newlabel{ex:prelim-plus}{{3.16}{10}{单圈实辐射的典型积分}{theorem.3.16}{}}
\newlabel{ex:prelim-plus@cref}{{[example][16][3]3.16}{[1][10][]10}}
\newlabel{def:prelim-mellin}{{3.17}{10}{Mellin 变换与卷积}{theorem.3.17}{}}
\newlabel{def:prelim-mellin@cref}{{[definition][17][3]3.17}{[1][10][]10}}
\newlabel{lem:prelim-mellin}{{3.18}{10}{Mellin 卷积定理与加号分布的矩}{theorem.3.18}{}}
\newlabel{lem:prelim-mellin@cref}{{[lemma][18][3]3.18}{[1][10][]10}}
\newlabel{rem:prelim-moment-tools}{{3.19}{10}{矩问题工具箱}{theorem.3.19}{}}
\newlabel{rem:prelim-moment-tools@cref}{{[remark][19][3]3.19}{[1][10][]10}}
\newlabel{def:prelim-cut}{{3.20}{10}{割图}{theorem.3.20}{}}
\newlabel{def:prelim-cut@cref}{{[definition][20][3]3.20}{[1][10][]10}}
\newlabel{sec:prelim-moments}{{3.5}{10}{第 5 节（测试）}{section.3.5}{}}
\newlabel{sec:prelim-moments@cref}{{[section][5][3]3.5}{[1][10][]10}}
\newlabel{thm:prelim-cutting-rules}{{3.21}{10}{切割规则 \tagR }{theorem.3.21}{}}
\newlabel{thm:prelim-cutting-rules@cref}{{[theorem][21][3]3.21}{[1][10][]10}}
\newlabel{prop:prelim-unitarity}{{3.22}{10}{物理幺正性与鬼场割线 \tagR }{theorem.3.22}{}}
\newlabel{prop:prelim-unitarity@cref}{{[proposition][22][3]3.22}{[1][10][]10}}
\newlabel{rem:prelim-cutting-dr}{{3.23}{10}{无质量维数正规化中的割图规则 \tagA }{theorem.3.23}{}}
\newlabel{rem:prelim-cutting-dr@cref}{{[remark][23][3]3.23}{[1][10][]10}}
\newlabel{prop:prelim-optical}{{3.24}{11}{流关联函数的光学定理 \tagR }{theorem.3.24}{}}
\newlabel{prop:prelim-optical@cref}{{[proposition][24][3]3.24}{[1][11][]11}}
\newlabel{ex:prelim-vacpol}{{3.25}{11}{维数正规化中的单圈无质量流关联函数}{theorem.3.25}{}}
\newlabel{ex:prelim-vacpol@cref}{{[example][25][3]3.25}{[1][11][]11}}
\newlabel{ch:ir}{{4}{12}{费曼积分的红外结构：钳夹、幂次计数与红外安全}{chapter.4}{}}
\newlabel{ch:ir@cref}{{[chapter][4][]4}{[1][12][]12}}
\newlabel{sec:ir-landau}{{4.1}{12}{第 1 节（测试）}{section.4.1}{}}
\newlabel{sec:ir-landau@cref}{{[section][1][4]4.1}{[1][12][]12}}
\newlabel{def:ir-landau-setup}{{4.1}{12}{Feynman 参数表示与 Symanzik 多项式}{theorem.4.1}{}}
\newlabel{def:ir-landau-setup@cref}{{[definition][1][4]4.1}{[1][12][]12}}
\newlabel{lem:ir-envelope}{{4.2}{12}{包络恒等式}{theorem.4.2}{}}
\newlabel{lem:ir-envelope@cref}{{[lemma][2][4]4.2}{[1][12][]12}}
\newlabel{thm:ir-landau-alpha}{{4.3}{12}{Landau 条件（参数形式）：解析性与红外有限性判据 \tagR }{theorem.4.3}{}}
\newlabel{thm:ir-landau-alpha@cref}{{[theorem][3][4]4.3}{[1][12][]12}}
\newlabel{thm:ir-landau-momentum}{{4.4}{12}{动量空间的钳夹必要条件 \tagR }{theorem.4.4}{}}
\newlabel{thm:ir-landau-momentum@cref}{{[theorem][4][4]4.4}{[1][12][]12}}
\newlabel{rem:ir-landau-status}{{4.5}{12}{Landau 条件的充分性与整体形变 \tagA }{theorem.4.5}{}}
\newlabel{rem:ir-landau-status@cref}{{[remark][5][4]4.5}{[1][12][]12}}
\newlabel{sec:ir-coleman-norton}{{4.2}{12}{第 2 节（测试）}{section.4.2}{}}
\newlabel{sec:ir-coleman-norton@cref}{{[section][2][4]4.2}{[1][12][]12}}
\newlabel{rem:ir-euclid-corners}{{4.6}{12}{欧氏区的 Landau 解与红外有限性}{theorem.4.6}{}}
\newlabel{rem:ir-euclid-corners@cref}{{[remark][6][4]4.6}{[1][12][]12}}
\newlabel{ex:ir-euclid-triangle}{{4.7}{12}{离壳三角图的 Symanzik 多项式}{theorem.4.7}{}}
\newlabel{ex:ir-euclid-triangle@cref}{{[example][7][4]4.7}{[1][12][]12}}
\newlabel{def:ir-pinch-surface}{{4.8}{13}{钳夹面、内禀坐标与法坐标}{theorem.4.8}{}}
\newlabel{def:ir-pinch-surface@cref}{{[definition][8][4]4.8}{[1][12][]13}}
\newlabel{prop:ir-coleman-norton}{{4.9}{13}{Coleman–Norton 定理（代数形式） \tagR }{theorem.4.9}{}}
\newlabel{prop:ir-coleman-norton@cref}{{[proposition][9][4]4.9}{[1][12][]13}}
\newlabel{prop:ir-pinch-classification}{{4.10}{13}{钳夹面与约化图 \tagA }{theorem.4.10}{}}
\newlabel{prop:ir-pinch-classification@cref}{{[proposition][10][4]4.10}{[1][13][]13}}
\newlabel{sec:ir-pinch}{{4.3}{13}{第 3 节（测试）}{section.4.3}{}}
\newlabel{sec:ir-pinch@cref}{{[section][3][4]4.3}{[1][13][]13}}
\newlabel{def:ir-scalings}{{4.11}{13}{区域标度}{theorem.4.11}{}}
\newlabel{def:ir-scalings@cref}{{[definition][11][4]4.11}{[1][13][]13}}
\newlabel{prop:ir-power-counting}{{4.12}{13}{超表观红外发散度 \tagR }{theorem.4.12}{}}
\newlabel{prop:ir-power-counting@cref}{{[proposition][12][4]4.12}{[1][13][]13}}
\newlabel{prop:ir-super-leading}{{4.13}{13}{超领头构型及其相消 \tagA }{theorem.4.13}{}}
\newlabel{prop:ir-super-leading@cref}{{[proposition][13][4]4.13}{[1][13][]13}}
\newlabel{rem:ir-scaling-vs-bound}{{4.14}{13}{标度陈述与积分界}{theorem.4.14}{}}
\newlabel{rem:ir-scaling-vs-bound@cref}{{[remark][14][4]4.14}{[1][13][]13}}
\newlabel{thm:ir-leading-regions}{{4.15}{13}{领头区域 \tagA }{theorem.4.15}{}}
\newlabel{thm:ir-leading-regions@cref}{{[theorem][15][4]4.15}{[1][13][]13}}
\newlabel{sec:ir-power}{{4.4}{14}{第 4 节（测试）}{section.4.4}{}}
\newlabel{sec:ir-power@cref}{{[section][4][4]4.4}{[1][13][]14}}
\newlabel{ex:ir-vertex}{{4.16}{14}{单圈无质量顶点}{theorem.4.16}{}}
\newlabel{ex:ir-vertex@cref}{{[example][16][4]4.16}{[1][13][]14}}
\newlabel{ex:ir-k-dominance}{{4.17}{14}{Feynman 规范下共线区域的标量极化主导}{theorem.4.17}{}}
\newlabel{ex:ir-k-dominance@cref}{{[example][17][4]4.17}{[1][13][]14}}
\newlabel{ex:ir-real-emission}{{4.18}{14}{$e^+e^-\to q\bar qg$ 的共线区}{theorem.4.18}{}}
\newlabel{ex:ir-real-emission@cref}{{[example][18][4]4.18}{[1][13][]14}}
\newlabel{def:ir-glauber}{{4.19}{14}{Glauber 区}{theorem.4.19}{}}
\newlabel{def:ir-glauber@cref}{{[definition][19][4]4.19}{[1][14][]14}}
\newlabel{prop:ir-glauber-trapping}{{4.20}{14}{单圈 Glauber 区的钳夹判据 \tagR }{theorem.4.20}{}}
\newlabel{prop:ir-glauber-trapping@cref}{{[proposition][20][4]4.20}{[1][14][]14}}
\newlabel{sec:ir-glauber}{{4.5}{14}{第 5 节（测试）}{section.4.5}{}}
\newlabel{sec:ir-glauber@cref}{{[section][5][4]4.5}{[1][14][]14}}
\newlabel{thm:ir-bn-qed}{{4.21}{14}{QED 中的 Bloch–Nordsieck 定理 \tagR }{theorem.4.21}{}}
\newlabel{thm:ir-bn-qed@cref}{{[theorem][21][4]4.21}{[1][14][]14}}
\newlabel{thm:ir-kln}{{4.22}{14}{KLN 定理 \tagA }{theorem.4.22}{}}
\newlabel{thm:ir-kln@cref}{{[theorem][22][4]4.22}{[1][14][]14}}
\newlabel{prop:ir-bn-failure}{{4.23}{14}{有色初态时非阿贝尔 Bloch–Nordsieck 抵消的失效 \tagV }{theorem.4.23}{}}
\newlabel{prop:ir-bn-failure@cref}{{[proposition][23][4]4.23}{[1][14][]14}}
\newlabel{def:ir-safe}{{4.24}{14}{红外安全可观测量}{theorem.4.24}{}}
\newlabel{def:ir-safe@cref}{{[definition][24][4]4.24}{[1][14][]14}}
\newlabel{thm:ir-safe-finite}{{4.25}{14}{红外安全可观测量的有限性 \tagA }{theorem.4.25}{}}
\newlabel{thm:ir-safe-finite@cref}{{[theorem][25][4]4.25}{[1][14][]14}}
\newlabel{sec:ir-cancellation}{{4.6}{15}{第 6 节（测试）}{section.4.6}{}}
\newlabel{sec:ir-cancellation@cref}{{[section][6][4]4.6}{[1][14][]15}}
\newlabel{def:ir-regions}{{4.26}{15}{区域方法}{theorem.4.26}{}}
\newlabel{def:ir-regions@cref}{{[definition][26][4]4.26}{[1][14][]15}}
\newlabel{thm:ir-regions-euclid}{{4.27}{15}{欧氏极限的渐近展开 \tagR }{theorem.4.27}{}}
\newlabel{thm:ir-regions-euclid@cref}{{[theorem][27][4]4.27}{[1][15][]15}}
\newlabel{prop:ir-regions-minkowski}{{4.28}{15}{Minkowski 极限下区域方法的地位 \tagP }{theorem.4.28}{}}
\newlabel{prop:ir-regions-minkowski@cref}{{[proposition][28][4]4.28}{[1][15][]15}}
\newlabel{rem:ir-regions-values}{{4.29}{15}{发散区域积分的值依赖于表示}{theorem.4.29}{}}
\newlabel{rem:ir-regions-values@cref}{{[remark][29][4]4.29}{[1][15][]15}}
\newlabel{ex:ir-regions-toy}{{4.30}{15}{一维玩具积分与零区减除}{theorem.4.30}{}}
\newlabel{ex:ir-regions-toy@cref}{{[example][30][4]4.30}{[1][15][]15}}
\newlabel{ch:sudakov}{{5}{16}{Sudakov 形状因子与振幅的红外结构}{chapter.5}{}}
\newlabel{ch:sudakov@cref}{{[chapter][5][]5}{[1][16][]16}}
\newlabel{sec:sudakov-oneloop}{{5.1}{16}{第 1 节（测试）}{section.5.1}{}}
\newlabel{sec:sudakov-oneloop@cref}{{[section][1][5]5.1}{[1][16][]16}}
\newlabel{def:sudakov-ff}{{5.1}{16}{夸克形状因子}{theorem.5.1}{}}
\newlabel{def:sudakov-ff@cref}{{[definition][1][5]5.1}{[1][16][]16}}
\newlabel{prop:sudakov-oneloop}{{5.2}{16}{一圈形状因子（精确到 $\varepsilon $ 的所有阶） \tagR }{theorem.5.2}{}}
\newlabel{prop:sudakov-oneloop@cref}{{[proposition][2][5]5.2}{[1][16][]16}}
\newlabel{ex:sudakov-onshell-hard}{{5.3}{16}{在壳无质量形状因子 = 硬区}{theorem.5.3}{}}
\newlabel{ex:sudakov-onshell-hard@cref}{{[example][3][5]5.3}{[1][16][]16}}
\newlabel{sec:sudakov-regions}{{5.2}{16}{第 2 节（测试）}{section.5.2}{}}
\newlabel{sec:sudakov-regions@cref}{{[section][2][5]5.2}{[1][16][]16}}
\newlabel{ex:sudakov-offshell}{{5.4}{16}{离壳 Sudakov 三角形的区域展开}{theorem.5.4}{}}
\newlabel{ex:sudakov-offshell@cref}{{[example][4][5]5.4}{[1][16][]16}}
\newlabel{ex:sudakov-massive}{{5.5}{16}{有质量交换：需要解析正规化}{theorem.5.5}{}}
\newlabel{ex:sudakov-massive@cref}{{[example][5][5]5.5}{[1][16][]16}}
\newlabel{ex:sudakov-ll}{{5.6}{16}{双对数系数}{theorem.5.6}{}}
\newlabel{ex:sudakov-ll@cref}{{[example][6][5]5.6}{[1][16][]16}}
\newlabel{sec:sudakov-factorization}{{5.3}{16}{第 3 节（测试）}{section.5.3}{}}
\newlabel{sec:sudakov-factorization@cref}{{[section][3][5]5.3}{[1][16][]16}}
\newlabel{rem:sudakov-consistency}{{5.7}{16}{一致性论证：单标度结构、$\mu $ 无关性与极点相消}{theorem.5.7}{}}
\newlabel{rem:sudakov-consistency@cref}{{[remark][7][5]5.7}{[1][16][]16}}
\newlabel{thm:sudakov-factorization}{{5.8}{16}{Sudakov 形状因子的全阶因子化 \tagA }{theorem.5.8}{}}
\newlabel{thm:sudakov-factorization@cref}{{[theorem][8][5]5.8}{[1][16][]16}}
\newlabel{def:sudakov-cusp}{{5.9}{17}{cusp 反常维度}{theorem.5.9}{}}
\newlabel{def:sudakov-cusp@cref}{{[definition][9][5]5.9}{[1][17][]17}}
\newlabel{sec:sudakov-evolution}{{5.4}{17}{第 4 节（测试）}{section.5.4}{}}
\newlabel{sec:sudakov-evolution@cref}{{[section][4][5]5.4}{[1][17][]17}}
\newlabel{prop:sudakov-kg}{{5.10}{17}{$K+G$ 演化方程 \tagA }{theorem.5.10}{}}
\newlabel{prop:sudakov-kg@cref}{{[proposition][10][5]5.10}{[1][17][]17}}
\newlabel{cor:sudakov-exponentiation}{{5.11}{17}{指数化与解析延拓 \tagA }{theorem.5.11}{}}
\newlabel{cor:sudakov-exponentiation@cref}{{[corollary][11][5]5.11}{[1][17][]17}}
\newlabel{prop:sudakov-webs}{{5.12}{17}{非阿贝尔指数化（webs） \tagA }{theorem.5.12}{}}
\newlabel{prop:sudakov-webs@cref}{{[proposition][12][5]5.12}{[1][17][]17}}
\newlabel{sec:sudakov-multileg}{{5.5}{17}{第 5 节（测试）}{section.5.5}{}}
\newlabel{sec:sudakov-multileg@cref}{{[section][5][5]5.5}{[1][17][]17}}
\newlabel{thm:sudakov-ir-poles}{{5.13}{17}{无质量多腿振幅红外极点的因子化 \tagA }{theorem.5.13}{}}
\newlabel{thm:sudakov-ir-poles@cref}{{[theorem][13][5]5.13}{[1][17][]17}}
\newlabel{prop:sudakov-dipole-violation}{{5.14}{17}{偶极子公式作为全阶精确公式的失效 \tagV }{theorem.5.14}{}}
\newlabel{prop:sudakov-dipole-violation@cref}{{[proposition][14][5]5.14}{[1][17][]17}}
\newlabel{prop:sudakov-coulomb}{{5.15}{17}{Coulomb 相位 \tagR }{theorem.5.15}{}}
\newlabel{prop:sudakov-coulomb@cref}{{[proposition][15][5]5.15}{[1][17][]17}}
\newlabel{sec:sudakov-collinear}{{5.6}{18}{第 6 节（测试）}{section.5.6}{}}
\newlabel{sec:sudakov-collinear@cref}{{[section][6][5]5.6}{[1][17][]18}}
\newlabel{prop:sudakov-amplitude-factorization}{{5.16}{18}{振幅的硬—软—共线因子化（出射或阈值运动学） \tagA }{theorem.5.16}{}}
\newlabel{prop:sudakov-amplitude-factorization@cref}{{[proposition][16][5]5.16}{[1][17][]18}}
\newlabel{rem:sudakov-spacelike}{{5.17}{18}{类空共线极限预览 \tagV }{theorem.5.17}{}}
\newlabel{rem:sudakov-spacelike@cref}{{[remark][17][5]5.17}{[1][18][]18}}
\newlabel{ch:ee}{{6}{20}{严格核心 I：$e^+e^-$ 湮灭总截面}{chapter.6}{}}
\newlabel{ch:ee@cref}{{[chapter][6][]6}{[1][20][]20}}
\newlabel{sec:ee-R}{{6.1}{20}{第 1 节（测试）}{section.6.1}{}}
\newlabel{sec:ee-R@cref}{{[section][1][6]6.1}{[1][20][]20}}
\newlabel{def:ee-R}{{6.1}{20}{$R$ 比}{theorem.6.1}{}}
\newlabel{def:ee-R@cref}{{[definition][1][6]6.1}{[1][20][]20}}
\newlabel{prop:ee-R-ImPi}{{6.2}{20}{$R=12\pi \,{\rm Im}\,\Pi $ \tagR }{theorem.6.2}{}}
\newlabel{prop:ee-R-ImPi@cref}{{[proposition][2][6]6.2}{[1][20][]20}}
\newlabel{lem:ee-single-scale}{{6.3}{20}{无质量传播子型图的单标度结构 \tagR }{theorem.6.3}{}}
\newlabel{lem:ee-single-scale@cref}{{[lemma][3][6]6.3}{[1][20][]20}}
\newlabel{sec:ee-analytic}{{6.2}{20}{第 2 节（测试）}{section.6.2}{}}
\newlabel{sec:ee-analytic@cref}{{[section][2][6]6.2}{[1][20][]20}}
\newlabel{prop:ee-analytic}{{6.4}{20}{解析性与正定性 \tagR }{theorem.6.4}{}}
\newlabel{prop:ee-analytic@cref}{{[proposition][4][6]6.4}{[1][20][]20}}
\newlabel{def:ee-adler}{{6.5}{20}{Adler 函数}{theorem.6.5}{}}
\newlabel{def:ee-adler@cref}{{[definition][5][6]6.5}{[1][20][]20}}
\newlabel{prop:ee-adler-dispersion}{{6.6}{20}{Adler 函数的色散表示 \tagR }{theorem.6.6}{}}
\newlabel{prop:ee-adler-dispersion@cref}{{[proposition][6][6]6.6}{[1][20][]20}}
\newlabel{sec:ee-euclid}{{6.3}{20}{第 3 节（测试）}{section.6.3}{}}
\newlabel{sec:ee-euclid@cref}{{[section][3][6]6.3}{[1][20][]20}}
\newlabel{thm:ee-euclid-ir}{{6.7}{20}{欧氏红外有限性 \tagR }{theorem.6.7}{}}
\newlabel{thm:ee-euclid-ir@cref}{{[theorem][7][6]6.7}{[1][20][]20}}
\newlabel{prop:ee-adler-rg}{{6.8}{21}{Adler 函数的重整化群 \tagR }{theorem.6.8}{}}
\newlabel{prop:ee-adler-rg@cref}{{[proposition][8][6]6.8}{[1][20][]21}}
\newlabel{thm:ee-R-finite}{{6.9}{21}{$R(s)$ 在微扰论每一阶的红外有限性 \tagR }{theorem.6.9}{}}
\newlabel{thm:ee-R-finite@cref}{{[theorem][9][6]6.9}{[1][21][]21}}
\newlabel{sec:ee-theorem}{{6.4}{21}{第 4 节（测试）}{section.6.4}{}}
\newlabel{sec:ee-theorem@cref}{{[section][4][6]6.4}{[1][21][]21}}
\newlabel{prop:ee-oneloop}{{6.10}{21}{$O(\alpha _s)$ 的实辐射与虚修正 \tagR }{theorem.6.10}{}}
\newlabel{prop:ee-oneloop@cref}{{[proposition][10][6]6.10}{[1][21][]21}}
\newlabel{ex:ee-quark-spectrum}{{6.11}{21}{KLN 的作用与局限}{theorem.6.11}{}}
\newlabel{ex:ee-quark-spectrum@cref}{{[example][11][6]6.11}{[1][21][]21}}
\newlabel{prop:ee-twoloop}{{6.12}{21}{欧氏两圈关联函数 \tagR }{theorem.6.12}{}}
\newlabel{prop:ee-twoloop@cref}{{[proposition][12][6]6.12}{[1][21][]21}}
\newlabel{sec:ee-oneloop}{{6.5}{21}{第 5 节（测试）}{section.6.5}{}}
\newlabel{sec:ee-oneloop@cref}{{[section][5][6]6.5}{[1][21][]21}}
\newlabel{ex:ee-twoloop-route}{{6.13}{21}{欧氏两圈路线}{theorem.6.13}{}}
\newlabel{ex:ee-twoloop-route@cref}{{[example][13][6]6.13}{[1][21][]21}}
\newlabel{prop:ee-cutting-exact}{{6.14}{21}{$O(g^2)$ 处精确到 $\varepsilon $ 的切割恒等式 \tagR }{theorem.6.14}{}}
\newlabel{prop:ee-cutting-exact@cref}{{[proposition][14][6]6.14}{[1][21][]21}}
\newlabel{cor:ee-cuts}{{6.15}{21}{与部分子末态求和的等同；KLN \tagA }{theorem.6.15}{}}
\newlabel{cor:ee-cuts@cref}{{[corollary][15][6]6.15}{[1][21][]21}}
\newlabel{sec:ee-higher}{{6.6}{21}{第 6 节（测试）}{section.6.6}{}}
\newlabel{sec:ee-higher@cref}{{[section][6][6]6.6}{[1][21][]21}}
\newlabel{prop:ee-contour}{{6.16}{21}{围道公式 \tagR }{theorem.6.16}{}}
\newlabel{prop:ee-contour@cref}{{[proposition][16][6]6.16}{[1][21][]21}}
\newlabel{cor:ee-R-alphas}{{6.17}{22}{$R(s)$ 的微扰展开 \tagR }{theorem.6.17}{}}
\newlabel{cor:ee-R-alphas@cref}{{[corollary][17][6]6.17}{[1][21][]22}}
\newlabel{ex:ee-pi2}{{6.18}{22}{$O(a^3)$ 的 $\pi ^2$ 项}{theorem.6.18}{}}
\newlabel{ex:ee-pi2@cref}{{[example][18][6]6.18}{[1][22][]22}}
\newlabel{sec:ee-ope}{{6.7}{22}{第 7 节（测试）}{section.6.7}{}}
\newlabel{sec:ee-ope@cref}{{[section][7][6]6.7}{[1][22][]22}}
\newlabel{rem:ee-status-layout}{{6.19}{22}{四个逻辑层次}{theorem.6.19}{}}
\newlabel{rem:ee-status-layout@cref}{{[remark][19][6]6.19}{[1][22][]22}}
\newlabel{thm:ee-ope}{{6.20}{22}{流关联函数的欧氏 OPE \tagA }{theorem.6.20}{}}
\newlabel{thm:ee-ope@cref}{{[theorem][20][6]6.20}{[1][22][]22}}
\newlabel{conj:ee-np-ope}{{6.21}{22}{非微扰 OPE 与凝聚 \tagC }{theorem.6.21}{}}
\newlabel{conj:ee-np-ope@cref}{{[conjecture][21][6]6.21}{[1][22][]22}}
\newlabel{sec:ee-duality}{{6.8}{22}{第 8 节（测试）}{section.6.8}{}}
\newlabel{sec:ee-duality@cref}{{[section][8][6]6.8}{[1][22][]22}}
\newlabel{rem:ee-complex-ope}{{6.22}{22}{复 $q^2$ 处的 OPE 是额外假设}{theorem.6.22}{}}
\newlabel{rem:ee-complex-ope@cref}{{[remark][22][6]6.22}{[1][22][]22}}
\newlabel{prop:ee-smearing}{{6.23}{22}{Poggio–Quinn–Weinberg 抹平恒等式 \tagR }{theorem.6.23}{}}
\newlabel{prop:ee-smearing@cref}{{[proposition][23][6]6.23}{[1][22][]22}}
\newlabel{conj:ee-duality}{{6.24}{22}{夸克—强子对偶 \tagC }{theorem.6.24}{}}
\newlabel{conj:ee-duality@cref}{{[conjecture][24][6]6.24}{[1][22][]22}}
\newlabel{ch:dis}{{7}{23}{严格核心 II：深度非弹散射与算符乘积展开}{chapter.7}{}}
\newlabel{ch:dis@cref}{{[chapter][7][]7}{[1][23][]23}}
\newlabel{sec:dis-kinematics}{{7.1}{23}{第 1 节（测试）}{section.7.1}{}}
\newlabel{sec:dis-kinematics@cref}{{[section][1][7]7.1}{[1][23][]23}}
\newlabel{def:dis-kinematics}{{7.1}{23}{DIS 运动学与强子张量}{theorem.7.1}{}}
\newlabel{def:dis-kinematics@cref}{{[definition][1][7]7.1}{[1][23][]23}}
\newlabel{prop:dis-W-structure}{{7.2}{23}{强子张量的一般结构 \tagR }{theorem.7.2}{}}
\newlabel{prop:dis-W-structure@cref}{{[proposition][2][7]7.2}{[1][23][]23}}
\newlabel{def:dis-compton}{{7.3}{23}{前向 Compton 振幅}{theorem.7.3}{}}
\newlabel{def:dis-compton@cref}{{[definition][3][7]7.3}{[1][23][]23}}
\newlabel{prop:dis-compton-analytic}{{7.4}{23}{解析性与交叉（部分子靶） \tagR }{theorem.7.4}{}}
\newlabel{prop:dis-compton-analytic@cref}{{[proposition][4][7]7.4}{[1][23][]23}}
\newlabel{sec:dis-compton}{{7.2}{23}{第 2 节（测试）}{section.7.2}{}}
\newlabel{sec:dis-compton@cref}{{[section][2][7]7.2}{[1][23][]23}}
\newlabel{prop:dis-dispersion}{{7.5}{23}{色散关系：Taylor 系数 = 矩（部分子靶） \tagR }{theorem.7.5}{}}
\newlabel{prop:dis-dispersion@cref}{{[proposition][5][7]7.5}{[1][23][]23}}
\newlabel{ex:dis-free-quark}{{7.6}{23}{自由夸克的前向振幅}{theorem.7.6}{}}
\newlabel{ex:dis-free-quark@cref}{{[example][6][7]7.6}{[1][23][]23}}
\newlabel{def:dis-twist}{{7.7}{23}{扭度与扭度 2 算符}{theorem.7.7}{}}
\newlabel{def:dis-twist@cref}{{[definition][7][7]7.7}{[1][23][]23}}
\newlabel{lem:dis-twist-classification}{{7.8}{23}{扭度 2 规范不变算符的分类 \tagR }{theorem.7.8}{}}
\newlabel{lem:dis-twist-classification@cref}{{[lemma][8][7]7.8}{[1][23][]23}}
\newlabel{sec:dis-ope}{{7.3}{24}{第 3 节（测试）}{section.7.3}{}}
\newlabel{sec:dis-ope@cref}{{[section][3][7]7.3}{[1][23][]24}}
\newlabel{lem:dis-twist-chirality}{{7.9}{24}{扭度奇偶性与手征 \tagR }{theorem.7.9}{}}
\newlabel{lem:dis-twist-chirality@cref}{{[lemma][9][7]7.9}{[1][23][]24}}
\newlabel{lem:dis-tree-ope}{{7.10}{24}{树图 OPE \tagR }{theorem.7.10}{}}
\newlabel{lem:dis-tree-ope@cref}{{[lemma][10][7]7.10}{[1][24][]24}}
\newlabel{prop:dis-ope-euclid}{{7.11}{24}{欧氏非例外动量下 Green 函数的逐图渐近展开 \tagA }{theorem.7.11}{}}
\newlabel{prop:dis-ope-euclid@cref}{{[proposition][11][7]7.11}{[1][24][]24}}
\newlabel{prop:dis-ope-gi}{{7.12}{24}{物理矩阵元只需规范不变算符 \tagA }{theorem.7.12}{}}
\newlabel{prop:dis-ope-gi@cref}{{[proposition][12][7]7.12}{[1][24][]24}}
\newlabel{sec:dis-np}{{7.4}{24}{第 4 节（测试）}{section.7.4}{}}
\newlabel{sec:dis-np@cref}{{[section][4][7]7.4}{[1][24][]24}}
\newlabel{prop:dis-ope-forward}{{7.13}{24}{前向在壳部分子矩阵元 \tagA }{theorem.7.13}{}}
\newlabel{prop:dis-ope-forward@cref}{{[proposition][13][7]7.13}{[1][24][]24}}
\newlabel{thm:dis-ope}{{7.14}{24}{重整化微扰论中的欧氏 OPE \tagA }{theorem.7.14}{}}
\newlabel{thm:dis-ope@cref}{{[theorem][14][7]7.14}{[1][24][]24}}
\newlabel{prop:dis-chiral-suppression}{{7.15}{24}{奇扭度系数的手征压低 \tagR }{theorem.7.15}{}}
\newlabel{prop:dis-chiral-suppression@cref}{{[proposition][15][7]7.15}{[1][24][]24}}
\newlabel{prop:dis-wilson-partons}{{7.16}{25}{由部分子结构函数提取 Wilson 系数 \tagR }{theorem.7.16}{}}
\newlabel{prop:dis-wilson-partons@cref}{{[proposition][16][7]7.16}{[1][24][]25}}
\newlabel{sec:dis-moments}{{7.5}{25}{第 5 节（测试）}{section.7.5}{}}
\newlabel{sec:dis-moments@cref}{{[section][5][7]7.5}{[1][25][]25}}
\newlabel{prop:dis-ope-rg}{{7.17}{25}{Wilson 系数与约化矩阵元的重整化群 \tagR }{theorem.7.17}{}}
\newlabel{prop:dis-ope-rg@cref}{{[proposition][17][7]7.17}{[1][25][]25}}
\newlabel{lem:dis-euclid-segment}{{7.18}{25}{欧氏点在 $\omega $ 平面上的位置 \tagR }{theorem.7.18}{}}
\newlabel{lem:dis-euclid-segment@cref}{{[lemma][18][7]7.18}{[1][25][]25}}
\newlabel{prop:dis-dispersion-np}{{7.19}{25}{由 Wightman 公理得到的有限减除色散关系 \tagR }{theorem.7.19}{}}
\newlabel{prop:dis-dispersion-np@cref}{{[proposition][19][7]7.19}{[1][25][]25}}
\newlabel{prop:dis-low-moments}{{7.20}{25}{低阶矩与增长条件的等价 \tagR }{theorem.7.20}{}}
\newlabel{prop:dis-low-moments@cref}{{[proposition][20][7]7.20}{[1][25][]25}}
\newlabel{sec:dis-rg}{{7.6}{25}{第 6 节（测试）}{section.7.6}{}}
\newlabel{sec:dis-rg@cref}{{[section][6][7]7.6}{[1][25][]25}}
\newlabel{lem:dis-ope-segment}{{7.21}{25}{OPE 在线段上的多项式逼近 \tagR }{theorem.7.21}{}}
\newlabel{lem:dis-ope-segment@cref}{{[lemma][21][7]7.21}{[1][25][]25}}
\newlabel{thm:dis-fixed-n}{{7.22}{25}{由定域 OPE 得到固定 $n$ 的 Taylor 系数 \tagR }{theorem.7.22}{}}
\newlabel{thm:dis-fixed-n@cref}{{[theorem][22][7]7.22}{[1][25][]25}}
\newlabel{thm:dis-moments-np}{{7.23}{25}{DIS 矩因子化的非微扰条件定理 \tagR }{theorem.7.23}{}}
\newlabel{thm:dis-moments-np@cref}{{[theorem][23][7]7.23}{[1][25][]25}}
\newlabel{cor:dis-weak-xspace}{{7.24}{26}{$x$ 空间的弱收敛 \tagR }{theorem.7.24}{}}
\newlabel{cor:dis-weak-xspace@cref}{{[corollary][24][7]7.24}{[1][25][]26}}
\newlabel{sec:dis-wilson}{{7.7}{26}{第 7 节（测试）}{section.7.7}{}}
\newlabel{sec:dis-wilson@cref}{{[section][7][7]7.7}{[1][26][]26}}
\newlabel{thm:dis-moments}{{7.25}{26}{DIS 结构函数矩的因子化 \tagR }{theorem.7.25}{}}
\newlabel{thm:dis-moments@cref}{{[theorem][25][7]7.25}{[1][26][]26}}
\newlabel{cor:dis-scaling}{{7.26}{26}{标度破坏 \tagR }{theorem.7.26}{}}
\newlabel{cor:dis-scaling@cref}{{[corollary][26][7]7.26}{[1][26][]26}}
\newlabel{ex:dis-n2}{{7.27}{26}{$n=2$：能动张量}{theorem.7.27}{}}
\newlabel{ex:dis-n2@cref}{{[example][27][7]7.27}{[1][26][]26}}
\newlabel{prop:dis-c-oneloop}{{7.28}{26}{单圈 Wilson 系数 \tagR }{theorem.7.28}{}}
\newlabel{prop:dis-c-oneloop@cref}{{[proposition][28][7]7.28}{[1][26][]26}}
\newlabel{sec:dis-xspace}{{7.8}{26}{第 8 节（测试）}{section.7.8}{}}
\newlabel{sec:dis-xspace@cref}{{[section][8][7]7.8}{[1][26][]26}}
\newlabel{ex:dis-oneloop-ope}{{7.29}{26}{单圈 OPE 检验}{theorem.7.29}{}}
\newlabel{ex:dis-oneloop-ope@cref}{{[example][29][7]7.29}{[1][26][]26}}
\newlabel{rem:dis-uniqueness}{{7.30}{26}{偶数矩决定 C 偶组合（数学）}{theorem.7.30}{}}
\newlabel{rem:dis-uniqueness@cref}{{[remark][30][7]7.30}{[1][26][]26}}
\newlabel{rem:dis-nonuniform}{{7.31}{26}{矩的 $O(1/Q^2)$ 不蕴含逐点 $O(1/Q^2)$}{theorem.7.31}{}}
\newlabel{rem:dis-nonuniform@cref}{{[remark][31][7]7.31}{[1][26][]26}}
\newlabel{prop:dis-light-cone}{{7.32}{26}{光锥展开 \tagA }{theorem.7.32}{}}
\newlabel{prop:dis-light-cone@cref}{{[proposition][32][7]7.32}{[1][26][]26}}
\newlabel{sec:dis-higher-twist}{{7.9}{26}{第 9 节（测试）}{section.7.9}{}}
\newlabel{sec:dis-higher-twist@cref}{{[section][9][7]7.9}{[1][26][]26}}
\newlabel{prop:dis-twist4}{{7.33}{26}{扭度 4 修正的算符结构 \tagR }{theorem.7.33}{}}
\newlabel{prop:dis-twist4@cref}{{[proposition][33][7]7.33}{[1][26][]26}}
\newlabel{prop:dis-tmc}{{7.34}{27}{靶质量修正 \tagR }{theorem.7.34}{}}
\newlabel{prop:dis-tmc@cref}{{[proposition][34][7]7.34}{[1][26][]27}}
\newlabel{ex:dis-tmc-n2}{{7.35}{27}{第二阶矩的靶质量修正}{theorem.7.35}{}}
\newlabel{ex:dis-tmc-n2@cref}{{[example][35][7]7.35}{[1][27][]27}}
\newlabel{thm:dis-bjorken}{{7.36}{27}{Bjorken 求和规则 \tagR }{theorem.7.36}{}}
\newlabel{thm:dis-bjorken@cref}{{[theorem][36][7]7.36}{[1][27][]27}}
\newlabel{sec:dis-polarized}{{7.10}{27}{第 10 节（测试）}{section.7.10}{}}
\newlabel{sec:dis-polarized@cref}{{[section][10][7]7.10}{[1][27][]27}}
\newlabel{rem:dis-crewther}{{7.37}{27}{Crewther 关系}{theorem.7.37}{}}
\newlabel{rem:dis-crewther@cref}{{[remark][37][7]7.37}{[1][27][]27}}
\newlabel{ch:pdf}{{8}{28}{严格核心 III：部分子分布与碎裂函数的算符定义}{chapter.8}{}}
\newlabel{ch:pdf@cref}{{[chapter][8][]8}{[1][28][]28}}
\newlabel{sec:pdf-definitions}{{8.1}{28}{第 1 节（测试）}{section.8.1}{}}
\newlabel{sec:pdf-definitions@cref}{{[section][1][8]8.1}{[1][28][]28}}
\newlabel{def:pdf-quark}{{8.1}{28}{夸克分布}{theorem.8.1}{}}
\newlabel{def:pdf-quark@cref}{{[definition][1][8]8.1}{[1][28][]28}}
\newlabel{def:pdf-antiquark}{{8.2}{28}{反夸克分布}{theorem.8.2}{}}
\newlabel{def:pdf-antiquark@cref}{{[definition][2][8]8.2}{[1][28][]28}}
\newlabel{def:pdf-gluon}{{8.3}{28}{胶子分布}{theorem.8.3}{}}
\newlabel{def:pdf-gluon@cref}{{[definition][3][8]8.3}{[1][28][]28}}
\newlabel{prop:pdf-invariance}{{8.4}{28}{规范不变性、推动不变性与自由场值 \tagR }{theorem.8.4}{}}
\newlabel{prop:pdf-invariance@cref}{{[proposition][4][8]8.4}{[1][28][]28}}
\newlabel{sec:pdf-properties}{{8.2}{28}{第 2 节（测试）}{section.8.2}{}}
\newlabel{sec:pdf-properties@cref}{{[section][2][8]8.2}{[1][28][]28}}
\newlabel{ex:pdf-free-field}{{8.5}{28}{自由场值}{theorem.8.5}{}}
\newlabel{ex:pdf-free-field@cref}{{[example][5][8]8.5}{[1][28][]28}}
\newlabel{thm:pdf-support-pt}{{8.6}{28}{支撑（微扰论） \tagR }{theorem.8.6}{}}
\newlabel{thm:pdf-support-pt@cref}{{[theorem][6][8]8.6}{[1][28][]28}}
\newlabel{prop:pdf-support-hadron}{{8.7}{28}{支撑（强子态） \tagA }{theorem.8.7}{}}
\newlabel{prop:pdf-support-hadron@cref}{{[proposition][7][8]8.7}{[1][28][]28}}
\newlabel{rem:pdf-support-jaffe}{{8.8}{29}{支撑的第二种论证：$k^-$ 极点}{theorem.8.8}{}}
\newlabel{rem:pdf-support-jaffe@cref}{{[remark][8][8]8.8}{[1][28][]29}}
\newlabel{sec:pdf-interpretation}{{8.3}{29}{第 3 节（测试）}{section.8.3}{}}
\newlabel{sec:pdf-interpretation@cref}{{[section][3][8]8.3}{[1][29][]29}}
\newlabel{rem:pdf-number-density}{{8.9}{29}{数密度诠释（启发式）}{theorem.8.9}{}}
\newlabel{rem:pdf-number-density@cref}{{[remark][9][8]8.9}{[1][29][]29}}
\newlabel{rem:pdf-gauge-residual}{{8.10}{29}{光锥规范的剩余规范自由与 $1/k^+$ 处方}{theorem.8.10}{}}
\newlabel{rem:pdf-gauge-residual@cref}{{[remark][10][8]8.10}{[1][29][]29}}
\newlabel{def:pdf-lightray}{{8.11}{29}{光线算符}{theorem.8.11}{}}
\newlabel{def:pdf-lightray@cref}{{[definition][11][8]8.11}{[1][29][]29}}
\newlabel{lem:pdf-wilson-moments}{{8.12}{29}{Wilson 线顶点的矩 = 定域算符顶点 \tagR }{theorem.8.12}{}}
\newlabel{lem:pdf-wilson-moments@cref}{{[lemma][12][8]8.12}{[1][29][]29}}
\newlabel{sec:pdf-moments}{{8.4}{29}{第 4 节（测试）}{section.8.4}{}}
\newlabel{sec:pdf-moments@cref}{{[section][4][8]8.4}{[1][29][]29}}
\newlabel{thm:pdf-moments-local}{{8.13}{29}{PDF 的矩 = 扭度 2 定域算符矩阵元 \tagR }{theorem.8.13}{}}
\newlabel{thm:pdf-moments-local@cref}{{[theorem][13][8]8.13}{[1][29][]29}}
\newlabel{thm:pdf-dglap}{{8.14}{29}{由算符重整化得到 DGLAP 方程 \tagR }{theorem.8.14}{}}
\newlabel{thm:pdf-dglap@cref}{{[theorem][14][8]8.14}{[1][29][]29}}
\newlabel{prop:pdf-lo-kernels}{{8.15}{29}{领头阶分裂函数 \tagR }{theorem.8.15}{}}
\newlabel{prop:pdf-lo-kernels@cref}{{[proposition][15][8]8.15}{[1][29][]29}}
\newlabel{prop:pdf-pqq-operator}{{8.16}{29}{Feynman 规范下由算符定义计算 $P_{qq}$ \tagR }{theorem.8.16}{}}
\newlabel{prop:pdf-pqq-operator@cref}{{[proposition][16][8]8.16}{[1][29][]29}}
\newlabel{sec:pdf-dglap}{{8.5}{29}{第 5 节（测试）}{section.8.5}{}}
\newlabel{sec:pdf-dglap@cref}{{[section][5][8]8.5}{[1][29][]29}}
\newlabel{ex:pdf-oneloop}{{8.17}{29}{单圈夸克 PDF（逐图）}{theorem.8.17}{}}
\newlabel{ex:pdf-oneloop@cref}{{[example][17][8]8.17}{[1][29][]29}}
\newlabel{prop:pdf-partonic-dr}{{8.18}{29}{纯维数正规化中的部分子分布 \tagR }{theorem.8.18}{}}
\newlabel{prop:pdf-partonic-dr@cref}{{[proposition][18][8]8.18}{[1][29][]29}}
\newlabel{prop:pdf-local-anomalous}{{8.19}{30}{定域算符的单圈反常维度（含混合） \tagR }{theorem.8.19}{}}
\newlabel{prop:pdf-local-anomalous@cref}{{[proposition][19][8]8.19}{[1][29][]30}}
\newlabel{ex:pdf-moments-exact}{{8.20}{30}{精确到 $\varepsilon $ 的单圈矩}{theorem.8.20}{}}
\newlabel{ex:pdf-moments-exact@cref}{{[example][20][8]8.20}{[1][30][]30}}
\newlabel{sec:pdf-oneloop}{{8.6}{30}{第 6 节（测试）}{section.8.6}{}}
\newlabel{sec:pdf-oneloop@cref}{{[section][6][8]8.6}{[1][30][]30}}
\newlabel{ex:pdf-quark-in-gluon}{{8.21}{30}{胶子中的夸克}{theorem.8.21}{}}
\newlabel{ex:pdf-quark-in-gluon@cref}{{[example][21][8]8.21}{[1][30][]30}}
\newlabel{thm:pdf-dis-xspace}{{8.22}{30}{DIS 的 $x$ 空间因子化（扭度 2 部分） \tagR }{theorem.8.22}{}}
\newlabel{thm:pdf-dis-xspace@cref}{{[theorem][22][8]8.22}{[1][30][]30}}
\newlabel{prop:pdf-sum-rules}{{8.23}{30}{求和规则 \tagR }{theorem.8.23}{}}
\newlabel{prop:pdf-sum-rules@cref}{{[proposition][23][8]8.23}{[1][30][]30}}
\newlabel{rem:pdf-ns-plusminus}{{8.24}{30}{解析延拓与 $P^\pm _{\rm NS}$}{theorem.8.24}{}}
\newlabel{rem:pdf-ns-plusminus@cref}{{[remark][24][8]8.24}{[1][30][]30}}
\newlabel{sec:pdf-dis}{{8.7}{30}{第 7 节（测试）}{section.8.7}{}}
\newlabel{sec:pdf-dis@cref}{{[section][7][8]8.7}{[1][30][]30}}
\newlabel{rem:pdf-schemes}{{8.25}{30}{方案变换}{theorem.8.25}{}}
\newlabel{rem:pdf-schemes@cref}{{[remark][25][8]8.25}{[1][30][]30}}
\newlabel{conj:pdf-positivity}{{8.26}{30}{$\overline {\rm MS}$ PDF 的正定性 \tagC }{theorem.8.26}{}}
\newlabel{conj:pdf-positivity@cref}{{[conjecture][26][8]8.26}{[1][30][]30}}
\newlabel{def:pdf-ff}{{8.27}{30}{碎裂函数}{theorem.8.27}{}}
\newlabel{def:pdf-ff@cref}{{[definition][27][8]8.27}{[1][30][]30}}
\newlabel{rem:pdf-ff-multren}{{8.28}{30}{碎裂函数的乘法可重整性是独立输入}{theorem.8.28}{}}
\newlabel{rem:pdf-ff-multren@cref}{{[remark][28][8]8.28}{[1][30][]30}}
\newlabel{sec:pdf-sumrules}{{8.8}{30}{第 8 节（测试）}{section.8.8}{}}
\newlabel{sec:pdf-sumrules@cref}{{[section][8][8]8.8}{[1][30][]30}}
\newlabel{prop:pdf-timelike-dglap}{{8.29}{30}{类时 DGLAP 演化 \tagA }{theorem.8.29}{}}
\newlabel{prop:pdf-timelike-dglap@cref}{{[proposition][29][8]8.29}{[1][30][]30}}
\newlabel{def:pdf-polarized}{{8.30}{31}{螺旋度与横度分布}{theorem.8.30}{}}
\newlabel{def:pdf-polarized@cref}{{[definition][30][8]8.30}{[1][31][]31}}
\newlabel{ch:method}{{9}{33}{全阶因子化证明的一般方法}{chapter.9}{}}
\newlabel{ch:method@cref}{{[chapter][9][]9}{[1][33][]33}}
\newlabel{sec:method-strategy}{{9.1}{33}{第 1 节（测试）}{section.9.1}{}}
\newlabel{sec:method-strategy@cref}{{[section][1][9]9.1}{[1][33][]33}}
\newlabel{def:method-regions}{{9.1}{33}{区域：逐线模式指派与偏序}{theorem.9.1}{}}
\newlabel{def:method-regions@cref}{{[definition][1][9]9.1}{[1][33][]33}}
\newlabel{lem:method-meets}{{9.2}{33}{逐线交与钳夹集的交}{theorem.9.2}{}}
\newlabel{lem:method-meets@cref}{{[lemma][2][9]9.2}{[1][33][]33}}
\newlabel{def:method-power-suppressed}{{9.3}{33}{积分范数意义下的幂次压低}{theorem.9.3}{}}
\newlabel{def:method-power-suppressed@cref}{{[definition][3][9]9.3}{[1][33][]33}}
\newlabel{def:method-neighbourhoods}{{9.4}{33}{依赖于 $\lambda $ 的邻域}{theorem.9.4}{}}
\newlabel{def:method-neighbourhoods@cref}{{[definition][4][9]9.4}{[1][33][]33}}
\newlabel{lem:method-covering}{{9.5}{33}{邻域覆盖}{theorem.9.5}{}}
\newlabel{lem:method-covering@cref}{{[lemma][5][9]9.5}{[1][33][]33}}
\newlabel{def:method-collinear-approx}{{9.6}{33}{共线近似子（幂等投影形式）}{theorem.9.6}{}}
\newlabel{def:method-collinear-approx@cref}{{[definition][6][9]9.6}{[1][33][]33}}
\newlabel{sec:method-regions}{{9.2}{33}{第 2 节（测试）}{section.9.2}{}}
\newlabel{sec:method-regions@cref}{{[section][2][9]9.2}{[1][33][]33}}
\newlabel{def:method-soft-approx}{{9.7}{33}{软近似子（幂等投影形式）}{theorem.9.7}{}}
\newlabel{def:method-soft-approx@cref}{{[definition][7][9]9.7}{[1][33][]33}}
\newlabel{lem:method-i0-compat}{{9.8}{33}{软近似分母的 $i0$ 相容性与变形幅度 \tagR }{theorem.9.8}{}}
\newlabel{lem:method-i0-compat@cref}{{[lemma][8][9]9.8}{[1][33][]33}}
\newlabel{lem:method-telescoping}{{9.9}{33}{乘积的伸缩恒等式}{theorem.9.9}{}}
\newlabel{lem:method-telescoping@cref}{{[lemma][9][9]9.9}{[1][33][]33}}
\newlabel{lem:method-collinear-accuracy}{{9.10}{33}{共线近似的精度 \tagA }{theorem.9.10}{}}
\newlabel{lem:method-collinear-accuracy@cref}{{[lemma][10][9]9.10}{[1][33][]33}}
\newlabel{lem:method-soft-accuracy}{{9.11}{34}{软近似的精度 \tagA }{theorem.9.11}{}}
\newlabel{lem:method-soft-accuracy@cref}{{[lemma][11][9]9.11}{[1][33][]34}}
\newlabel{lem:method-local-accuracy}{{9.12}{34}{局部精度（积分范数） \tagR }{theorem.9.12}{}}
\newlabel{lem:method-local-accuracy@cref}{{[lemma][12][9]9.12}{[1][34][]34}}
\newlabel{sec:method-gy}{{9.3}{34}{第 3 节（测试）}{section.9.3}{}}
\newlabel{sec:method-gy@cref}{{[section][3][9]9.3}{[1][34][]34}}
\newlabel{rem:method-abs-failure}{{9.13}{34}{吸收性质 (Abs) 何时失效}{theorem.9.13}{}}
\newlabel{rem:method-abs-failure@cref}{{[remark][13][9]9.13}{[1][34][]34}}
\newlabel{lem:method-gy}{{9.14}{34}{Grammer–Yennie 分解}{theorem.9.14}{}}
\newlabel{lem:method-gy@cref}{{[lemma][14][9]9.14}{[1][34][]34}}
\newlabel{prop:method-gy-suppression}{{9.15}{34}{G 部分的幂次压低 \tagA }{theorem.9.15}{}}
\newlabel{prop:method-gy-suppression@cref}{{[proposition][15][9]9.15}{[1][34][]34}}
\newlabel{lem:method-ward-elementary}{{9.16}{34}{基本 Ward 恒等式 \tagR }{theorem.9.16}{}}
\newlabel{lem:method-ward-elementary@cref}{{[lemma][16][9]9.16}{[1][34][]34}}
\newlabel{lem:method-ward-abelian}{{9.17}{34}{阿贝尔 Ward 恒等式 \tagR }{theorem.9.17}{}}
\newlabel{lem:method-ward-abelian@cref}{{[lemma][17][9]9.17}{[1][34][]34}}
\newlabel{prop:method-ward-one}{{9.18}{34}{单个纵向胶子的非阿贝尔 Ward 恒等式 \tagR }{theorem.9.18}{}}
\newlabel{prop:method-ward-one@cref}{{[proposition][18][9]9.18}{[1][34][]34}}
\newlabel{sec:method-ward}{{9.4}{34}{第 4 节（测试）}{section.9.4}{}}
\newlabel{sec:method-ward@cref}{{[section][4][9]9.4}{[1][34][]34}}
\newlabel{prop:method-ward-tree}{{9.19}{34}{树图层次的非阿贝尔恒等式与 PDF Wilson 线的种子 \tagR }{theorem.9.19}{}}
\newlabel{prop:method-ward-tree@cref}{{[proposition][19][9]9.19}{[1][34][]34}}
\newlabel{ex:method-ward-tree}{{9.20}{35}{$q\bar q\to g^*g^*$ 的两个纵向胶子}{theorem.9.20}{}}
\newlabel{ex:method-ward-tree@cref}{{[example][20][9]9.20}{[1][34][]35}}
\newlabel{lem:method-ward-multi}{{9.21}{35}{多个 K 胶子的 Ward 恒等式与 Wilson 线 \tagA }{theorem.9.21}{}}
\newlabel{lem:method-ward-multi@cref}{{[lemma][21][9]9.21}{[1][34][]35}}
\newlabel{prop:method-collinear-wilson}{{9.22}{35}{共线 Wilson 线与 PDF \tagA }{theorem.9.22}{}}
\newlabel{prop:method-collinear-wilson@cref}{{[proposition][22][9]9.22}{[1][35][]35}}
\newlabel{prop:method-soft-wilson}{{9.23}{35}{软 Wilson 线 \tagA }{theorem.9.23}{}}
\newlabel{prop:method-soft-wilson@cref}{{[proposition][23][9]9.23}{[1][35][]35}}
\newlabel{rem:method-pure-gauge}{{9.24}{35}{纯规范论证的地位}{theorem.9.24}{}}
\newlabel{rem:method-pure-gauge@cref}{{[remark][24][9]9.24}{[1][35][]35}}
\newlabel{sec:method-subtraction}{{9.5}{35}{第 5 节（测试）}{section.9.5}{}}
\newlabel{sec:method-subtraction@cref}{{[section][5][9]9.5}{[1][35][]35}}
\newlabel{def:method-subtraction}{{9.25}{35}{减除项}{theorem.9.25}{}}
\newlabel{def:method-subtraction@cref}{{[definition][25][9]9.25}{[1][35][]35}}
\newlabel{ex:method-naive-induction}{{9.26}{35}{朴素归纳为何失败：离壳单圈三角图 \tagR }{theorem.9.26}{}}
\newlabel{ex:method-naive-induction@cref}{{[example][26][9]9.26}{[1][35][]35}}
\newlabel{def:method-collins-conditions}{{9.27}{35}{减除条件 (CS1)–(CS3)}{theorem.9.27}{}}
\newlabel{def:method-collins-conditions@cref}{{[definition][27][9]9.27}{[1][35][]35}}
\newlabel{prop:method-chains}{{9.28}{35}{减除的链展开}{theorem.9.28}{}}
\newlabel{prop:method-chains@cref}{{[proposition][28][9]9.28}{[1][35][]35}}
\newlabel{lem:method-pairing}{{9.29}{35}{配对分解}{theorem.9.29}{}}
\newlabel{lem:method-pairing@cref}{{[lemma][29][9]9.29}{[1][35][]35}}
\newlabel{prop:method-product}{{9.30}{35}{乘积公式}{theorem.9.30}{}}
\newlabel{prop:method-product@cref}{{[proposition][30][9]9.30}{[1][35][]35}}
\newlabel{sec:method-unitarity}{{9.6}{36}{第 6 节（测试）}{section.9.6}{}}
\newlabel{sec:method-unitarity@cref}{{[section][6][9]9.6}{[1][35][]36}}
\newlabel{lem:method-no-pole}{{9.31}{36}{有界的亚纯函数无极点}{theorem.9.31}{}}
\newlabel{lem:method-no-pole@cref}{{[lemma][31][9]9.31}{[1][35][]36}}
\newlabel{cor:method-removing-eps}{{9.32}{36}{去掉维数正规化 \tagR }{theorem.9.32}{}}
\newlabel{cor:method-removing-eps@cref}{{[corollary][32][9]9.32}{[1][36][]36}}
\newlabel{thm:method-subtraction}{{9.33}{36}{减除定理 \tagA }{theorem.9.33}{}}
\newlabel{thm:method-subtraction@cref}{{[theorem][33][9]9.33}{[1][36][]36}}
\newlabel{ex:method-sudakov-subtraction}{{9.34}{36}{完整的单圈例子：离壳三角图的 Collins 减除}{theorem.9.34}{}}
\newlabel{ex:method-sudakov-subtraction@cref}{{[example][34][9]9.34}{[1][36][]36}}
\newlabel{ex:method-dis-oneloop}{{9.35}{36}{减除语言中的单圈 DIS}{theorem.9.35}{}}
\newlabel{ex:method-dis-oneloop@cref}{{[example][35][9]9.35}{[1][36][]36}}
\newlabel{lem:method-soft-unitarity}{{9.36}{36}{软因子的幺正性 \tagR }{theorem.9.36}{}}
\newlabel{lem:method-soft-unitarity@cref}{{[lemma][36][9]9.36}{[1][36][]36}}
\newlabel{sec:method-glauber}{{9.7}{36}{第 7 节（测试）}{section.9.7}{}}
\newlabel{sec:method-glauber@cref}{{[section][7][9]9.7}{[1][36][]36}}
\newlabel{prop:method-fsi}{{9.37}{36}{末态相互作用的幺正性相消 \tagA }{theorem.9.37}{}}
\newlabel{prop:method-fsi@cref}{{[proposition][37][9]9.37}{[1][36][]36}}
\newlabel{prop:method-glauber-oneloop}{{9.38}{36}{两条程函线之间的单胶子交换 \tagR }{theorem.9.38}{}}
\newlabel{prop:method-glauber-oneloop@cref}{{[proposition][38][9]9.38}{[1][36][]36}}
\newlabel{ex:method-glauber-oneloop}{{9.39}{36}{单胶子交换的极点图}{theorem.9.39}{}}
\newlabel{ex:method-glauber-oneloop@cref}{{[example][39][9]9.39}{[1][36][]36}}
\newlabel{prop:method-glauber-deformation}{{9.40}{36}{Glauber 区的围道变形判据 \tagA }{theorem.9.40}{}}
\newlabel{prop:method-glauber-deformation@cref}{{[proposition][40][9]9.40}{[1][36][]36}}
\newlabel{rem:method-two-component}{{9.41}{36}{一个分量的形变为何不够}{theorem.9.41}{}}
\newlabel{rem:method-two-component@cref}{{[remark][41][9]9.41}{[1][36][]36}}
\newlabel{thm:method-template}{{9.42}{37}{因子化模板定理 \tagA }{theorem.9.42}{}}
\newlabel{thm:method-template@cref}{{[theorem][42][9]9.42}{[1][36][]37}}
\newlabel{sec:method-template}{{9.8}{37}{第 8 节（测试）}{section.9.8}{}}
\newlabel{sec:method-template@cref}{{[section][8][9]9.8}{[1][37][]37}}
\newlabel{prop:method-scale-scheme}{{9.43}{37}{因子化标度与方案 \tagR }{theorem.9.43}{}}
\newlabel{prop:method-scale-scheme@cref}{{[proposition][43][9]9.43}{[1][37][]37}}
\newlabel{ch:disall}{{10}{38}{DIS 的全阶动量空间证明、重夸克与核靶}{chapter.10}{}}
\newlabel{ch:disall@cref}{{[chapter][10][]10}{[1][38][]38}}
\newlabel{sec:disall-regions}{{10.1}{38}{第 1 节（测试）}{section.10.1}{}}
\newlabel{sec:disall-regions@cref}{{[section][1][10]10.1}{[1][38][]38}}
\newlabel{def:disall-breit}{{10.1}{38}{Breit 系与投影算符}{theorem.10.1}{}}
\newlabel{def:disall-breit@cref}{{[definition][1][10]10.1}{[1][38][]38}}
\newlabel{prop:disall-two-routes}{{10.2}{38}{从非物理区前向振幅到矩与 $x$ 空间的两条途径 \tagR }{theorem.10.2}{}}
\newlabel{prop:disall-two-routes@cref}{{[proposition][2][10]10.2}{[1][38][]38}}
\newlabel{lem:disall-fsi}{{10.3}{38}{末态喷注的非钳夹性 \tagA }{theorem.10.3}{}}
\newlabel{lem:disall-fsi@cref}{{[lemma][3][10]10.3}{[1][38][]38}}
\newlabel{sec:disall-fsi}{{10.2}{38}{第 2 节（测试）}{section.10.2}{}}
\newlabel{sec:disall-fsi@cref}{{[section][2][10]10.2}{[1][38][]38}}
\newlabel{prop:disall-regions}{{10.4}{38}{DIS 的领头区域 \tagA }{theorem.10.4}{}}
\newlabel{prop:disall-regions@cref}{{[proposition][4][10]10.4}{[1][38][]38}}
\newlabel{prop:disall-factorized-subtraction}{{10.5}{38}{DIS 减除项的因子化结构 \tagR }{theorem.10.5}{}}
\newlabel{prop:disall-factorized-subtraction@cref}{{[proposition][5][10]10.5}{[1][38][]38}}
\newlabel{prop:disall-h2}{{10.6}{39}{(H2) 的验证 \tagA }{theorem.10.6}{}}
\newlabel{prop:disall-h2@cref}{{[proposition][6][10]10.6}{[1][38][]39}}
\newlabel{sec:disall-subtraction}{{10.3}{39}{第 3 节（测试）}{section.10.3}{}}
\newlabel{sec:disall-subtraction@cref}{{[section][3][10]10.3}{[1][39][]39}}
\newlabel{prop:disall-h3h4}{{10.7}{39}{(H3)(H4) 平凡成立 \tagA }{theorem.10.7}{}}
\newlabel{prop:disall-h3h4@cref}{{[proposition][7][10]10.7}{[1][39][]39}}
\newlabel{thm:disall-renorm-compat}{{10.8}{39}{减除与 $\overline {\rm MS}$ 重整化的相容性（矩空间） \tagR }{theorem.10.8}{}}
\newlabel{thm:disall-renorm-compat@cref}{{[theorem][8][10]10.8}{[1][39][]39}}
\newlabel{prop:disall-h5}{{10.9}{39}{(H5) 的验证 \tagA }{theorem.10.9}{}}
\newlabel{prop:disall-h5@cref}{{[proposition][9][10]10.9}{[1][39][]39}}
\newlabel{sec:disall-theorem}{{10.4}{39}{第 4 节（测试）}{section.10.4}{}}
\newlabel{sec:disall-theorem@cref}{{[section][4][10]10.4}{[1][39][]39}}
\newlabel{ex:disall-ladder}{{10.10}{39}{两级梯形图：三个区域与分组余项}{theorem.10.10}{}}
\newlabel{ex:disall-ladder@cref}{{[example][10][10]10.10}{[1][39][]39}}
\newlabel{thm:disall-factorization}{{10.11}{39}{DIS 共线因子化（动量空间，$x$ 空间逐点） \tagA }{theorem.10.11}{}}
\newlabel{thm:disall-factorization@cref}{{[theorem][11][10]10.11}{[1][39][]39}}
\newlabel{cor:disall-consistency}{{10.12}{40}{与 OPE 系数的一致性 \tagR }{theorem.10.12}{}}
\newlabel{cor:disall-consistency@cref}{{[corollary][12][10]10.12}{[1][39][]40}}
\newlabel{sec:disall-consistency}{{10.5}{40}{第 5 节（测试）}{section.10.5}{}}
\newlabel{sec:disall-consistency@cref}{{[section][5][10]10.5}{[1][40][]40}}
\newlabel{prop:disall-nlo}{{10.13}{40}{$O(\alpha _s)$ 的 $\overline {\rm MS}$ 系数函数 \tagR }{theorem.10.13}{}}
\newlabel{prop:disall-nlo@cref}{{[proposition][13][10]10.13}{[1][40][]40}}
\newlabel{ex:disall-nlo}{{10.14}{40}{$O(\alpha _s)$ 系数函数的完整推导}{theorem.10.14}{}}
\newlabel{ex:disall-nlo@cref}{{[example][14][10]10.14}{[1][40][]40}}
\newlabel{ex:disall-gluon}{{10.15}{40}{光子—胶子融合}{theorem.10.15}{}}
\newlabel{ex:disall-gluon@cref}{{[example][15][10]10.15}{[1][40][]40}}
\newlabel{sec:disall-nlo}{{10.6}{40}{第 6 节（测试）}{section.10.6}{}}
\newlabel{sec:disall-nlo@cref}{{[section][6][10]10.6}{[1][40][]40}}
\newlabel{prop:disall-scale}{{10.16}{40}{标度无关性 \tagR }{theorem.10.16}{}}
\newlabel{prop:disall-scale@cref}{{[proposition][16][10]10.16}{[1][40][]40}}
\newlabel{prop:disall-dis-scheme}{{10.17}{40}{DIS 方案 \tagR }{theorem.10.17}{}}
\newlabel{prop:disall-dis-scheme@cref}{{[proposition][17][10]10.17}{[1][40][]40}}
\newlabel{thm:disall-heavy}{{10.18}{40}{含重夸克质量的因子化（Collins 1998） \tagA }{theorem.10.18}{}}
\newlabel{thm:disall-heavy@cref}{{[theorem][18][10]10.18}{[1][40][]40}}
\newlabel{sec:disall-scale}{{10.7}{40}{第 7 节（测试）}{section.10.7}{}}
\newlabel{sec:disall-scale@cref}{{[section][7][10]10.7}{[1][40][]40}}
\newlabel{ex:disall-heavy-lo}{{10.19}{40}{领头阶 $\gamma ^*g\to Q\bar Q$}{theorem.10.19}{}}
\newlabel{ex:disall-heavy-lo@cref}{{[example][19][10]10.19}{[1][40][]40}}
\newlabel{prop:disall-nuclear}{{10.20}{41}{核靶上的 DIS \tagA }{theorem.10.20}{}}
\newlabel{prop:disall-nuclear@cref}{{[proposition][20][10]10.20}{[1][40][]41}}
\newlabel{ch:coll}{{11}{42}{其他共线过程：单举湮灭、SIDIS、光生、极化与横度}{chapter.11}{}}
\newlabel{ch:coll@cref}{{[chapter][11][]11}{[1][42][]42}}
\newlabel{sec:coll-sia}{{11.1}{42}{第 1 节（测试）}{section.11.1}{}}
\newlabel{sec:coll-sia@cref}{{[section][1][11]11.1}{[1][42][]42}}
\newlabel{def:coll-sia}{{11.1}{42}{单举湮灭的运动学}{theorem.11.1}{}}
\newlabel{def:coll-sia@cref}{{[definition][1][11]11.1}{[1][42][]42}}
\newlabel{lem:coll-final-poles}{{11.2}{42}{全部为末态的共线子图的极点在同一半平面 \tagR }{theorem.11.2}{}}
\newlabel{lem:coll-final-poles@cref}{{[lemma][2][11]11.2}{[1][42][]42}}
\newlabel{lem:coll-unobserved-jet}{{11.3}{42}{未观测喷注与软辐射的幺正性相消 \tagA }{theorem.11.3}{}}
\newlabel{lem:coll-unobserved-jet@cref}{{[lemma][3][11]11.3}{[1][42][]42}}
\newlabel{prop:coll-sia-regions}{{11.4}{42}{SIA 的领头区域 \tagA }{theorem.11.4}{}}
\newlabel{prop:coll-sia-regions@cref}{{[proposition][4][11]11.4}{[1][42][]42}}
\newlabel{thm:coll-sia}{{11.5}{42}{$e^+e^-\to hX$ 的因子化 \tagA }{theorem.11.5}{}}
\newlabel{thm:coll-sia@cref}{{[theorem][5][11]11.5}{[1][42][]42}}
\newlabel{sec:coll-sidis}{{11.2}{43}{第 2 节（测试）}{section.11.2}{}}
\newlabel{sec:coll-sidis@cref}{{[section][2][11]11.2}{[1][42][]43}}
\newlabel{rem:coll-sia-nlo}{{11.6}{43}{系数函数与检验}{theorem.11.6}{}}
\newlabel{rem:coll-sia-nlo@cref}{{[remark][6][11]11.6}{[1][42][]43}}
\newlabel{ex:coll-sia-lo}{{11.7}{43}{领头阶 $e^+e^-\to hX$}{theorem.11.7}{}}
\newlabel{ex:coll-sia-lo@cref}{{[example][7][11]11.7}{[1][42][]43}}
\newlabel{thm:coll-sidis}{{11.8}{43}{共线 SIDIS 的因子化 \tagA }{theorem.11.8}{}}
\newlabel{thm:coll-sidis@cref}{{[theorem][8][11]11.8}{[1][43][]43}}
\newlabel{thm:coll-direct-photo}{{11.9}{43}{直接光生与光子结构函数 \tagA }{theorem.11.9}{}}
\newlabel{thm:coll-direct-photo@cref}{{[theorem][9][11]11.9}{[1][43][]43}}
\newlabel{prop:coll-photon-sf}{{11.10}{43}{光子结构函数的点状部分 \tagR }{theorem.11.10}{}}
\newlabel{prop:coll-photon-sf@cref}{{[proposition][10][11]11.10}{[1][43][]43}}
\newlabel{sec:coll-photo}{{11.3}{43}{第 3 节（测试）}{section.11.3}{}}
\newlabel{sec:coll-photo@cref}{{[section][3][11]11.3}{[1][43][]43}}
\newlabel{conj:coll-resolved-photo}{{11.11}{43}{分辨光生 \tagC }{theorem.11.11}{}}
\newlabel{conj:coll-resolved-photo@cref}{{[conjecture][11][11]11.11}{[1][43][]43}}
\newlabel{prop:coll-diffractive-photo}{{11.12}{43}{衍射光生中的因子化破坏 \tagV }{theorem.11.12}{}}
\newlabel{prop:coll-diffractive-photo@cref}{{[proposition][12][11]11.12}{[1][43][]43}}
\newlabel{prop:coll-luxqed}{{11.13}{43}{质子的光子分布 \tagA }{theorem.11.13}{}}
\newlabel{prop:coll-luxqed@cref}{{[proposition][13][11]11.13}{[1][43][]43}}
\newlabel{def:coll-g1}{{11.14}{44}{极化结构函数}{theorem.11.14}{}}
\newlabel{def:coll-g1@cref}{{[definition][14][11]11.14}{[1][44][]44}}
\newlabel{thm:coll-g1-moments}{{11.15}{44}{$g_1$ 的奇数矩 \tagR }{theorem.11.15}{}}
\newlabel{thm:coll-g1-moments@cref}{{[theorem][15][11]11.15}{[1][44][]44}}
\newlabel{sec:coll-polarized}{{11.4}{44}{第 4 节（测试）}{section.11.4}{}}
\newlabel{sec:coll-polarized@cref}{{[section][4][11]11.4}{[1][44][]44}}
\newlabel{prop:coll-anomaly-scheme}{{11.16}{44}{一阶矩的方案依赖 \tagR }{theorem.11.16}{}}
\newlabel{prop:coll-anomaly-scheme@cref}{{[proposition][16][11]11.16}{[1][44][]44}}
\newlabel{prop:coll-ellis-jaffe}{{11.17}{44}{Ellis–Jaffe 求和规则作为预言 \tagV }{theorem.11.17}{}}
\newlabel{prop:coll-ellis-jaffe@cref}{{[proposition][17][11]11.17}{[1][44][]44}}
\newlabel{conj:coll-bc}{{11.18}{44}{Burkhardt–Cottingham 求和规则 \tagC }{theorem.11.18}{}}
\newlabel{conj:coll-bc@cref}{{[conjecture][18][11]11.18}{[1][44][]44}}
\newlabel{prop:coll-ww}{{11.19}{44}{Wandzura–Wilczek 关系的一般失效 \tagV }{theorem.11.19}{}}
\newlabel{prop:coll-ww@cref}{{[proposition][19][11]11.19}{[1][44][]44}}
\newlabel{def:coll-transversity}{{11.20}{44}{横度分布}{theorem.11.20}{}}
\newlabel{def:coll-transversity@cref}{{[definition][20][11]11.20}{[1][44][]44}}
\newlabel{sec:coll-transversity}{{11.5}{44}{第 5 节（测试）}{section.11.5}{}}
\newlabel{sec:coll-transversity@cref}{{[section][5][11]11.5}{[1][44][]44}}
\newlabel{prop:coll-transversity-evolution}{{11.21}{44}{横度的演化 \tagR }{theorem.11.21}{}}
\newlabel{prop:coll-transversity-evolution@cref}{{[proposition][21][11]11.21}{[1][44][]44}}
\newlabel{prop:coll-soffer}{{11.22}{44}{Soffer 不等式（LO 演化的保持性） \tagR }{theorem.11.22}{}}
\newlabel{prop:coll-soffer@cref}{{[proposition][22][11]11.22}{[1][44][]44}}
\newlabel{rem:coll-christ-lee}{{11.23}{44}{单举 DIS 中无单靶自旋横向不对称}{theorem.11.23}{}}
\newlabel{rem:coll-christ-lee@cref}{{[remark][23][11]11.23}{[1][44][]44}}
\newlabel{ch:dy}{{12}{45}{Drell–Yan 过程与强子—强子碰撞}{chapter.12}{}}
\newlabel{ch:dy@cref}{{[chapter][12][]12}{[1][45][]45}}
\newlabel{sec:dy-born}{{12.1}{45}{第 1 节（测试）}{section.12.1}{}}
\newlabel{sec:dy-born@cref}{{[section][1][12]12.1}{[1][45][]45}}
\newlabel{def:dy-kin}{{12.1}{45}{DY 运动学}{theorem.12.1}{}}
\newlabel{def:dy-kin@cref}{{[definition][1][12]12.1}{[1][45][]45}}
\newlabel{prop:dy-born}{{12.2}{45}{Born 截面 \tagR }{theorem.12.2}{}}
\newlabel{prop:dy-born@cref}{{[proposition][2][12]12.2}{[1][45][]45}}
\newlabel{prop:dy-regions}{{12.3}{45}{DY 的领头区域 \tagA }{theorem.12.3}{}}
\newlabel{prop:dy-regions@cref}{{[proposition][3][12]12.3}{[1][45][]45}}
\newlabel{prop:dy-glauber-pinch}{{12.4}{45}{单个割图中的 Glauber 钳夹 \tagR }{theorem.12.4}{}}
\newlabel{prop:dy-glauber-pinch@cref}{{[proposition][4][12]12.4}{[1][45][]45}}
\newlabel{sec:dy-regions}{{12.2}{45}{第 2 节（测试）}{section.12.2}{}}
\newlabel{sec:dy-regions@cref}{{[section][2][12]12.2}{[1][45][]45}}
\newlabel{lem:dy-poles-orientation}{{12.5}{45}{光锥序中间态的取向与 $k^-$ 极点 \tagR }{theorem.12.5}{}}
\newlabel{lem:dy-poles-orientation@cref}{{[lemma][5][12]12.5}{[1][45][]45}}
\newlabel{rem:dy-one-time-per-jet}{{12.6}{45}{每个喷注一个光锥时间}{theorem.12.6}{}}
\newlabel{rem:dy-one-time-per-jet@cref}{{[remark][6][12]12.6}{[1][45][]45}}
\newlabel{lem:dy-fsi}{{12.7}{45}{旁观者末态相互作用的幺正性相消 \tagA }{theorem.12.7}{}}
\newlabel{lem:dy-fsi@cref}{{[lemma][7][12]12.7}{[1][45][]45}}
\newlabel{ex:dy-future-vs-past}{{12.8}{46}{单个割图给出未来指向的 Wilson 线 \tagR }{theorem.12.8}{}}
\newlabel{ex:dy-future-vs-past@cref}{{[example][8][12]12.8}{[1][45][]46}}
\newlabel{sec:dy-orientation}{{12.3}{46}{第 3 节（测试）}{section.12.3}{}}
\newlabel{sec:dy-orientation@cref}{{[section][3][12]12.3}{[1][46][]46}}
\newlabel{prop:dy-deformation}{{12.9}{46}{移除末态极点后的围道变形 \tagA }{theorem.12.9}{}}
\newlabel{prop:dy-deformation@cref}{{[proposition][9][12]12.9}{[1][46][]46}}
\newlabel{lem:dy-soft}{{12.10}{46}{软因子为 1 \tagA }{theorem.12.10}{}}
\newlabel{lem:dy-soft@cref}{{[lemma][10][12]12.10}{[1][46][]46}}
\newlabel{prop:dy-universality}{{12.11}{46}{DY 与 DIS 的共线 PDF 是同一对象 \tagR }{theorem.12.11}{}}
\newlabel{prop:dy-universality@cref}{{[proposition][11][12]12.11}{[1][46][]46}}
\newlabel{prop:dy-abelian-model}{{12.12}{46}{阿贝尔模型中旁观者 Glauber 相位的相消 \tagP }{theorem.12.12}{}}
\newlabel{prop:dy-abelian-model@cref}{{[proposition][12][12]12.12}{[1][46][]46}}
\newlabel{sec:dy-glauber}{{12.4}{46}{第 4 节（测试）}{section.12.4}{}}
\newlabel{sec:dy-glauber@cref}{{[section][4][12]12.4}{[1][46][]46}}
\newlabel{ex:dy-spectator-lcopt}{{12.13}{46}{旁观者模型的光锥序表示}{theorem.12.13}{}}
\newlabel{ex:dy-spectator-lcopt@cref}{{[example][13][12]12.13}{[1][46][]46}}
\newlabel{ex:dy-glauber-model}{{12.14}{46}{阿贝尔 Glauber 模型的数值演示}{theorem.12.14}{}}
\newlabel{ex:dy-glauber-model@cref}{{[example][14][12]12.14}{[1][46][]46}}
\newlabel{prop:dy-eikonal-twogluon}{{12.15}{46}{两条程函线之间的双胶子交换：$O(g^4)$ 幺正性与 Regge 相位 \tagR }{theorem.12.15}{}}
\newlabel{prop:dy-eikonal-twogluon@cref}{{[proposition][15][12]12.15}{[1][46][]46}}
\newlabel{rem:dy-common-ordering}{{12.16}{46}{共同时序是物理结果}{theorem.12.16}{}}
\newlabel{rem:dy-common-ordering@cref}{{[remark][16][12]12.16}{[1][46][]46}}
\newlabel{sec:dy-model}{{12.5}{47}{第 5 节（测试）}{section.12.5}{}}
\newlabel{sec:dy-model@cref}{{[section][5][12]12.5}{[1][46][]47}}
\newlabel{thm:dy-factorization}{{12.17}{47}{单举 DY 的共线因子化 \tagA }{theorem.12.17}{}}
\newlabel{thm:dy-factorization@cref}{{[theorem][17][12]12.17}{[1][46][]47}}
\newlabel{rem:dy-nnlo-not-glauber}{{12.18}{47}{固定阶部分子截面不检验旁观者 Glauber 相消}{theorem.12.18}{}}
\newlabel{rem:dy-nnlo-not-glauber@cref}{{[remark][18][12]12.18}{[1][47][]47}}
\newlabel{rem:dy-dft-premises}{{12.19}{47}{Doria–Frenkel–Taylor 非相消的前提}{theorem.12.19}{}}
\newlabel{rem:dy-dft-premises@cref}{{[remark][19][12]12.19}{[1][47][]47}}
\newlabel{prop:dy-nlo}{{12.20}{47}{$O(\alpha _s)$ 的 $\overline {\rm MS}$ 系数函数 \tagR }{theorem.12.20}{}}
\newlabel{prop:dy-nlo@cref}{{[proposition][20][12]12.20}{[1][47][]47}}
\newlabel{sec:dy-theorem}{{12.6}{47}{第 6 节（测试）}{section.12.6}{}}
\newlabel{sec:dy-theorem@cref}{{[section][6][12]12.6}{[1][47][]47}}
\newlabel{ex:dy-nlo}{{12.21}{47}{$O(\alpha _s)$ 系数函数的完整推导}{theorem.12.21}{}}
\newlabel{ex:dy-nlo@cref}{{[example][21][12]12.21}{[1][47][]47}}
\newlabel{prop:dy-scale}{{12.22}{47}{$\mu _F$ 无关性 \tagR }{theorem.12.22}{}}
\newlabel{prop:dy-scale@cref}{{[proposition][22][12]12.22}{[1][47][]47}}
\newlabel{rem:dy-dis-scheme}{{12.23}{47}{DIS 方案中的 DY 与 PDF 的普适性}{theorem.12.23}{}}
\newlabel{rem:dy-dis-scheme@cref}{{[remark][23][12]12.23}{[1][47][]47}}
\newlabel{ex:dy-threshold}{{12.24}{47}{大 $N$ 极限}{theorem.12.24}{}}
\newlabel{ex:dy-threshold@cref}{{[example][24][12]12.24}{[1][47][]47}}
\newlabel{sec:dy-nlo}{{12.7}{47}{第 7 节（测试）}{section.12.7}{}}
\newlabel{sec:dy-nlo@cref}{{[section][7][12]12.7}{[1][47][]47}}
\newlabel{thm:dy-colour-singlet}{{12.25}{47}{色单态末态 \tagA }{theorem.12.25}{}}
\newlabel{thm:dy-colour-singlet@cref}{{[theorem][25][12]12.25}{[1][47][]47}}
\newlabel{thm:dy-polarized}{{12.26}{48}{极化 Drell–Yan \tagA }{theorem.12.26}{}}
\newlabel{thm:dy-polarized@cref}{{[theorem][26][12]12.26}{[1][47][]48}}
\newlabel{conj:dy-coloured-final}{{12.27}{48}{强子对撞中有色末态的共线因子化 \tagC }{theorem.12.27}{}}
\newlabel{conj:dy-coloured-final@cref}{{[conjecture][27][12]12.27}{[1][48][]48}}
\newlabel{conj:dy-polarized-jets}{{12.28}{48}{极化强子对撞中的 $A_{LL}$（喷注、$\pi ^0$） \tagC }{theorem.12.28}{}}
\newlabel{conj:dy-polarized-jets@cref}{{[conjecture][28][12]12.28}{[1][48][]48}}
\newlabel{sec:dy-other}{{12.8}{48}{第 8 节（测试）}{section.12.8}{}}
\newlabel{sec:dy-other@cref}{{[section][8][12]12.8}{[1][48][]48}}
\newlabel{thm:dy-higher-twist}{{12.29}{48}{$1/Q^2$ 处的广义因子化 \tagA }{theorem.12.29}{}}
\newlabel{thm:dy-higher-twist@cref}{{[theorem][29][12]12.29}{[1][48][]48}}
\newlabel{prop:dy-twist6-failure}{{12.30}{48}{更高幂次的因子化失效 \tagV }{theorem.12.30}{}}
\newlabel{prop:dy-twist6-failure@cref}{{[proposition][30][12]12.30}{[1][48][]48}}
\newlabel{ch:thr}{{13}{49}{阈值重求和与重新因子化}{chapter.13}{}}
\newlabel{ch:thr@cref}{{[chapter][13][]13}{[1][49][]49}}
\newlabel{sec:thr-limit}{{13.1}{49}{第 1 节（测试）}{section.13.1}{}}
\newlabel{sec:thr-limit@cref}{{[section][1][13]13.1}{[1][49][]49}}
\newlabel{def:thr-limit}{{13.1}{49}{部分子阈值与强子阈值}{theorem.13.1}{}}
\newlabel{def:thr-limit@cref}{{[definition][1][13]13.1}{[1][49][]49}}
\newlabel{thm:thr-refactorization}{{13.2}{49}{DY 阈值重新因子化与指数化 \tagA }{theorem.13.2}{}}
\newlabel{thm:thr-refactorization@cref}{{[theorem][2][13]13.2}{[1][49][]49}}
\newlabel{sec:thr-refactorization}{{13.2}{49}{第 2 节（测试）}{section.13.2}{}}
\newlabel{sec:thr-refactorization@cref}{{[section][2][13]13.2}{[1][49][]49}}
\newlabel{prop:thr-oneloop}{{13.3}{49}{单圈检验 \tagR }{theorem.13.3}{}}
\newlabel{prop:thr-oneloop@cref}{{[proposition][3][13]13.3}{[1][49][]49}}
\newlabel{rem:thr-prescriptions}{{13.4}{49}{Landau 极点与处方}{theorem.13.4}{}}
\newlabel{rem:thr-prescriptions@cref}{{[remark][4][13]13.4}{[1][49][]49}}
\newlabel{sec:thr-check}{{13.3}{49}{第 3 节（测试）}{section.13.3}{}}
\newlabel{sec:thr-check@cref}{{[section][3][13]13.3}{[1][49][]49}}
\newlabel{prop:thr-dis}{{13.5}{49}{DIS 的大 $x$ 重新因子化 \tagA }{theorem.13.5}{}}
\newlabel{prop:thr-dis@cref}{{[proposition][5][13]13.5}{[1][49][]49}}
\newlabel{prop:thr-nlp}{{13.6}{49}{次领头幂次的阈值对数 \tagP }{theorem.13.6}{}}
\newlabel{prop:thr-nlp@cref}{{[proposition][6][13]13.6}{[1][49][]49}}
\newlabel{sec:thr-prescriptions}{{13.4}{50}{第 4 节（测试）}{section.13.4}{}}
\newlabel{sec:thr-prescriptions@cref}{{[section][4][13]13.4}{[1][49][]50}}
\newlabel{ex:thr-dy}{{13.7}{50}{指数的 $O(\alpha _s)$ 展开与 prop:dy-nlo 的 Mellin 变换}{theorem.13.7}{}}
\newlabel{ex:thr-dy@cref}{{[example][7][13]13.7}{[1][49][]50}}
\newlabel{ch:tmd}{{14}{52}{横动量依赖（TMD）因子化}{chapter.14}{}}
\newlabel{ch:tmd@cref}{{[chapter][14][]14}{[1][52][]52}}
\newlabel{sec:tmd-kinematics}{{14.1}{52}{第 1 节（测试）}{section.14.1}{}}
\newlabel{sec:tmd-kinematics@cref}{{[section][1][14]14.1}{[1][52][]52}}
\newlabel{def:tmd-kin}{{14.1}{52}{小 $q_T$ 区域}{theorem.14.1}{}}
\newlabel{def:tmd-kin@cref}{{[definition][1][14]14.1}{[1][52][]52}}
\newlabel{prop:tmd-regions}{{14.2}{52}{小 $q_T$ 的领头区域 \tagA }{theorem.14.2}{}}
\newlabel{prop:tmd-regions@cref}{{[proposition][2][14]14.2}{[1][52][]52}}
\newlabel{def:tmd-unsub}{{14.3}{52}{未减除的 TMD 关联函数}{theorem.14.3}{}}
\newlabel{def:tmd-unsub@cref}{{[definition][3][14]14.3}{[1][52][]52}}
\newlabel{def:tmd-soft}{{14.4}{52}{软因子}{theorem.14.4}{}}
\newlabel{def:tmd-soft@cref}{{[definition][4][14]14.4}{[1][52][]52}}
\newlabel{sec:tmd-regions}{{14.2}{52}{第 2 节（测试）}{section.14.2}{}}
\newlabel{sec:tmd-regions@cref}{{[section][2][14]14.2}{[1][52][]52}}
\newlabel{lem:tmd-soft-oneloop}{{14.5}{52}{单圈软因子 \tagR }{theorem.14.5}{}}
\newlabel{lem:tmd-soft-oneloop@cref}{{[lemma][5][14]14.5}{[1][52][]52}}
\newlabel{lem:tmd-sqrt-identity}{{14.6}{52}{平方根构造是恒等式 \tagR }{theorem.14.6}{}}
\newlabel{lem:tmd-sqrt-identity@cref}{{[lemma][6][14]14.6}{[1][52][]52}}
\newlabel{def:tmd-collins}{{14.7}{52}{Collins 的 TMD 定义}{theorem.14.7}{}}
\newlabel{def:tmd-collins@cref}{{[definition][7][14]14.7}{[1][52][]52}}
\newlabel{rem:tmd-schemes}{{14.8}{52}{其他方案}{theorem.14.8}{}}
\newlabel{rem:tmd-schemes@cref}{{[remark][8][14]14.8}{[1][52][]52}}
\newlabel{sec:tmd-definitions}{{14.3}{53}{第 3 节（测试）}{section.14.3}{}}
\newlabel{sec:tmd-definitions@cref}{{[section][3][14]14.3}{[1][52][]53}}
\newlabel{rem:tmd-scheme-constants}{{14.9}{53}{方案常数（$S_\varepsilon $ 约定）}{theorem.14.9}{}}
\newlabel{rem:tmd-scheme-constants@cref}{{[remark][9][14]14.9}{[1][52][]53}}
\newlabel{rem:tmd-order-of-limits}{{14.10}{53}{极限次序}{theorem.14.10}{}}
\newlabel{rem:tmd-order-of-limits@cref}{{[remark][10][14]14.10}{[1][52][]53}}
\newlabel{prop:tmd-zero-bin-soft}{{14.11}{53}{零箱减除即软因子 \tagR }{theorem.14.11}{}}
\newlabel{prop:tmd-zero-bin-soft@cref}{{[proposition][11][14]14.11}{[1][53][]53}}
\newlabel{rem:tmd-rapidity-window}{{14.12}{53}{快度窗口}{theorem.14.12}{}}
\newlabel{rem:tmd-rapidity-window@cref}{{[remark][12][14]14.12}{[1][53][]53}}
\newlabel{sec:tmd-subtraction}{{14.4}{53}{第 4 节（测试）}{section.14.4}{}}
\newlabel{sec:tmd-subtraction@cref}{{[section][4][14]14.4}{[1][53][]53}}
\newlabel{lem:tmd-glauber}{{14.13}{53}{TMD DY、SIDIS、$e^+e^-$ 中的 Glauber 相消 \tagA }{theorem.14.13}{}}
\newlabel{lem:tmd-glauber@cref}{{[lemma][13][14]14.13}{[1][53][]53}}
\newlabel{prop:tmd-h1h5}{{14.14}{53}{小 $q_T$ DY 对模板假设的验证 \tagA }{theorem.14.14}{}}
\newlabel{prop:tmd-h1h5@cref}{{[proposition][14][14]14.14}{[1][53][]53}}
\newlabel{thm:tmd-factorization}{{14.15}{53}{Drell–Yan 的 TMD 因子化 \tagA }{theorem.14.15}{}}
\newlabel{thm:tmd-factorization@cref}{{[theorem][15][14]14.15}{[1][53][]53}}
\newlabel{prop:tmd-evolution}{{14.16}{53}{Collins–Soper 方程与重整化群方程 \tagA }{theorem.14.16}{}}
\newlabel{prop:tmd-evolution@cref}{{[proposition][16][14]14.16}{[1][53][]53}}
\newlabel{sec:tmd-theorem}{{14.5}{54}{第 5 节（测试）}{section.14.5}{}}
\newlabel{sec:tmd-theorem@cref}{{[section][5][14]14.5}{[1][53][]54}}
\newlabel{prop:tmd-evolution-oneloop}{{14.17}{54}{单圈演化系数与硬因子的 RG 相容性 \tagR }{theorem.14.17}{}}
\newlabel{prop:tmd-evolution-oneloop@cref}{{[proposition][17][14]14.17}{[1][53][]54}}
\newlabel{prop:tmd-small-b}{{14.18}{54}{小 $b$ 展开 \tagA }{theorem.14.18}{}}
\newlabel{prop:tmd-small-b@cref}{{[proposition][18][14]14.18}{[1][54][]54}}
\newlabel{prop:tmd-css}{{14.19}{54}{CSS 解与非微扰输入 \tagA }{theorem.14.19}{}}
\newlabel{prop:tmd-css@cref}{{[proposition][19][14]14.19}{[1][54][]54}}
\newlabel{thm:tmd-sivers-sign}{{14.20}{54}{Sivers 函数的符号反转（算符层面） \tagR }{theorem.14.20}{}}
\newlabel{thm:tmd-sivers-sign@cref}{{[theorem][20][14]14.20}{[1][54][]54}}
\newlabel{sec:tmd-evolution}{{14.6}{54}{第 6 节（测试）}{section.14.6}{}}
\newlabel{sec:tmd-evolution@cref}{{[section][6][14]14.6}{[1][54][]54}}
\newlabel{cor:tmd-sivers-observable}{{14.21}{54}{可观测不对称的符号反转 \tagA }{theorem.14.21}{}}
\newlabel{cor:tmd-sivers-observable@cref}{{[corollary][21][14]14.21}{[1][54][]54}}
\newlabel{prop:tmd-collins-universality}{{14.22}{54}{T-奇碎裂函数的普适性 \tagA }{theorem.14.22}{}}
\newlabel{prop:tmd-collins-universality@cref}{{[proposition][22][14]14.22}{[1][54][]54}}
\newlabel{def:tmd-etqs}{{14.23}{54}{Qiu–Sterman (ETQS) 关联函数}{theorem.14.23}{}}
\newlabel{def:tmd-etqs@cref}{{[definition][23][14]14.23}{[1][54][]54}}
\newlabel{prop:tmd-twist3-matching}{{14.24}{54}{与扭度 3 共线因子化的匹配 \tagP }{theorem.14.24}{}}
\newlabel{prop:tmd-twist3-matching@cref}{{[proposition][24][14]14.24}{[1][54][]54}}
\newlabel{sec:tmd-small-b}{{14.7}{55}{第 7 节（测试）}{section.14.7}{}}
\newlabel{sec:tmd-small-b@cref}{{[section][7][14]14.7}{[1][55][]55}}
\newlabel{rem:tmd-kt-integral}{{14.25}{55}{$\int d^2k_T\,f\ne f$}{theorem.14.25}{}}
\newlabel{rem:tmd-kt-integral@cref}{{[remark][25][14]14.25}{[1][55][]55}}
\newlabel{ex:tmd-oneloop}{{14.26}{55}{单圈全过程}{theorem.14.26}{}}
\newlabel{ex:tmd-oneloop@cref}{{[example][26][14]14.26}{[1][55][]55}}
\newlabel{ex:tmd-qt-asymptotics}{{14.27}{55}{小 $q_T$ 渐近式的重现}{theorem.14.27}{}}
\newlabel{ex:tmd-qt-asymptotics@cref}{{[example][27][14]14.27}{[1][55][]55}}
\newlabel{ex:tmd-sivers-model}{{14.28}{55}{标量双夸克模型中的 Sivers 函数}{theorem.14.28}{}}
\newlabel{ex:tmd-sivers-model@cref}{{[example][28][14]14.28}{[1][55][]55}}
\newlabel{ch:scet}{{15}{56}{有效场论方法：SCET}{chapter.15}{}}
\newlabel{ch:scet@cref}{{[chapter][15][]15}{[1][56][]56}}
\newlabel{sec:scet-modes}{{15.1}{56}{第 1 节（测试）}{section.15.1}{}}
\newlabel{sec:scet-modes@cref}{{[section][1][15]15.1}{[1][56][]56}}
\newlabel{def:scet-modes}{{15.1}{56}{模式}{theorem.15.1}{}}
\newlabel{def:scet-modes@cref}{{[definition][1][15]15.1}{[1][56][]56}}
\newlabel{prop:scet-scalings}{{15.2}{56}{场的标度 \tagR }{theorem.15.2}{}}
\newlabel{prop:scet-scalings@cref}{{[proposition][2][15]15.2}{[1][56][]56}}
\newlabel{prop:scet-lagrangian}{{15.3}{56}{领头阶共线夸克拉氏量 \tagR }{theorem.15.3}{}}
\newlabel{prop:scet-lagrangian@cref}{{[proposition][3][15]15.3}{[1][56][]56}}
\newlabel{sec:scet-lagrangian}{{15.2}{56}{第 2 节（测试）}{section.15.2}{}}
\newlabel{sec:scet-lagrangian@cref}{{[section][2][15]15.2}{[1][56][]56}}
\newlabel{prop:scet-wilson}{{15.4}{56}{共线 Wilson 线 \tagR }{theorem.15.4}{}}
\newlabel{prop:scet-wilson@cref}{{[proposition][4][15]15.4}{[1][56][]56}}
\newlabel{ex:scet-wn-tree}{{15.5}{56}{$W_n$ 的树图匹配}{theorem.15.5}{}}
\newlabel{ex:scet-wn-tree@cref}{{[example][5][15]15.5}{[1][56][]56}}
\newlabel{prop:scet-gauge}{{15.6}{56}{共线与超软规范对称性 \tagR }{theorem.15.6}{}}
\newlabel{prop:scet-gauge@cref}{{[proposition][6][15]15.6}{[1][56][]56}}
\newlabel{sec:scet-decoupling}{{15.3}{56}{第 3 节（测试）}{section.15.3}{}}
\newlabel{sec:scet-decoupling@cref}{{[section][3][15]15.3}{[1][56][]56}}
\newlabel{thm:scet-decoupling}{{15.7}{56}{超软退耦（BPS 场重定义） \tagR }{theorem.15.7}{}}
\newlabel{thm:scet-decoupling@cref}{{[theorem][7][15]15.7}{[1][56][]56}}
\newlabel{ex:scet-decoupling-tree}{{15.8}{56}{BPS 场重定义的树图演示}{theorem.15.8}{}}
\newlabel{ex:scet-decoupling-tree@cref}{{[example][8][15]15.8}{[1][56][]56}}
\newlabel{rem:scet-y-direction}{{15.9}{57}{$Y_n$ 的方向与 $i0$}{theorem.15.9}{}}
\newlabel{rem:scet-y-direction@cref}{{[remark][9][15]15.9}{[1][56][]57}}
\newlabel{sec:scet-scet2}{{15.4}{57}{第 4 节（测试）}{section.15.4}{}}
\newlabel{sec:scet-scet2@cref}{{[section][4][15]15.4}{[1][57][]57}}
\newlabel{def:scet-zero-bin}{{15.10}{57}{零箱减除}{theorem.15.10}{}}
\newlabel{def:scet-zero-bin@cref}{{[definition][10][15]15.10}{[1][57][]57}}
\newlabel{prop:scet-scet2}{{15.11}{57}{SCET$_{\rm II}$ 与软 Wilson 线 \tagR }{theorem.15.11}{}}
\newlabel{prop:scet-scet2@cref}{{[proposition][11][15]15.11}{[1][57][]57}}
\newlabel{prop:scet-rapidity-rg}{{15.12}{57}{快度发散的处理与等价性 \tagA }{theorem.15.12}{}}
\newlabel{prop:scet-rapidity-rg@cref}{{[proposition][12][15]15.12}{[1][57][]57}}
\newlabel{sec:scet-glauber}{{15.5}{57}{第 5 节（测试）}{section.15.5}{}}
\newlabel{sec:scet-glauber@cref}{{[section][5][15]15.5}{[1][57][]57}}
\newlabel{prop:scet-glauber-operators}{{15.13}{57}{Glauber 算符 \tagA }{theorem.15.13}{}}
\newlabel{prop:scet-glauber-operators@cref}{{[proposition][13][15]15.13}{[1][57][]57}}
\newlabel{rem:scet-glauber-dictionary}{{15.14}{57}{Glauber SCET 与光锥序证明的对照}{theorem.15.14}{}}
\newlabel{rem:scet-glauber-dictionary@cref}{{[remark][14][15]15.14}{[1][57][]57}}
\newlabel{prop:scet-hard-matching}{{15.15}{57}{硬匹配系数 \tagR }{theorem.15.15}{}}
\newlabel{prop:scet-hard-matching@cref}{{[proposition][15][15]15.15}{[1][57][]57}}
\newlabel{sec:scet-thrust}{{15.6}{57}{第 6 节（测试）}{section.15.6}{}}
\newlabel{sec:scet-thrust@cref}{{[section][6][15]15.6}{[1][57][]57}}
\newlabel{thm:scet-thrust}{{15.16}{57}{推力分布的因子化 \tagA }{theorem.15.16}{}}
\newlabel{thm:scet-thrust@cref}{{[theorem][16][15]15.16}{[1][57][]57}}
\newlabel{prop:scet-thrust-oneloop}{{15.17}{58}{单圈因子与 RG 一致性 \tagR }{theorem.15.17}{}}
\newlabel{prop:scet-thrust-oneloop@cref}{{[proposition][17][15]15.17}{[1][57][]58}}
\newlabel{ex:scet-thrust-oneloop}{{15.18}{58}{单圈 $H$、$J$、$S_T$ 合成 QCD 奇异项}{theorem.15.18}{}}
\newlabel{ex:scet-thrust-oneloop@cref}{{[example][18][15]15.18}{[1][57][]58}}
\newlabel{sec:scet-rg}{{15.7}{58}{第 7 节（测试）}{section.15.7}{}}
\newlabel{sec:scet-rg@cref}{{[section][7][15]15.7}{[1][58][]58}}
\newlabel{prop:scet-ll}{{15.19}{58}{领头对数指数 \tagA }{theorem.15.19}{}}
\newlabel{prop:scet-ll@cref}{{[proposition][19][15]15.19}{[1][58][]58}}
\newlabel{rem:scet-assumptions}{{15.20}{58}{SCET 证明的前提}{theorem.15.20}{}}
\newlabel{rem:scet-assumptions@cref}{{[remark][20][15]15.20}{[1][58][]58}}
\newlabel{ch:jets}{{16}{59}{喷注、事件形状、非全局对数与超领头对数}{chapter.16}{}}
\newlabel{ch:jets@cref}{{[chapter][16][]16}{[1][59][]59}}
\newlabel{sec:jets-irsafe}{{16.1}{59}{第 1 节（测试）}{section.16.1}{}}
\newlabel{sec:jets-irsafe@cref}{{[section][1][16]16.1}{[1][59][]59}}
\newlabel{def:jets-algorithms}{{16.1}{59}{红外安全的喷注定义与事件形状}{theorem.16.1}{}}
\newlabel{def:jets-algorithms@cref}{{[definition][1][16]16.1}{[1][59][]59}}
\newlabel{prop:jets-sw}{{16.2}{59}{Sterman–Weinberg 两喷注截面 \tagR }{theorem.16.2}{}}
\newlabel{prop:jets-sw@cref}{{[proposition][2][16]16.2}{[1][59][]59}}
\newlabel{ex:jets-sw}{{16.3}{59}{$O(\alpha _s)$ 的 Sterman–Weinberg 截面}{theorem.16.3}{}}
\newlabel{ex:jets-sw@cref}{{[example][3][16]16.3}{[1][59][]59}}
\newlabel{sec:jets-hadronization}{{16.2}{59}{第 2 节（测试）}{section.16.2}{}}
\newlabel{sec:jets-hadronization@cref}{{[section][2][16]16.2}{[1][59][]59}}
\newlabel{cor:jets-ir-finite}{{16.4}{59}{红外安全喷注截面的有限性 \tagA }{theorem.16.4}{}}
\newlabel{cor:jets-ir-finite@cref}{{[corollary][4][16]16.4}{[1][59][]59}}
\newlabel{prop:jets-thrust-lo}{{16.5}{59}{领头阶推力分布 \tagR }{theorem.16.5}{}}
\newlabel{prop:jets-thrust-lo@cref}{{[proposition][5][16]16.5}{[1][59][]59}}
\newlabel{conj:jets-hadronization}{{16.6}{59}{强子化修正的普适性 \tagC }{theorem.16.6}{}}
\newlabel{conj:jets-hadronization@cref}{{[conjecture][6][16]16.6}{[1][59][]59}}
\newlabel{sec:jets-ngl}{{16.3}{60}{第 3 节（测试）}{section.16.3}{}}
\newlabel{sec:jets-ngl@cref}{{[section][3][16]16.3}{[1][59][]60}}
\newlabel{prop:jets-ngl}{{16.7}{60}{非全局对数的存在 \tagR }{theorem.16.7}{}}
\newlabel{prop:jets-ngl@cref}{{[proposition][7][16]16.7}{[1][59][]60}}
\newlabel{thm:jets-bms}{{16.8}{60}{BMS 方程 \tagA }{theorem.16.8}{}}
\newlabel{thm:jets-bms@cref}{{[theorem][8][16]16.8}{[1][60][]60}}
\newlabel{prop:jets-ngl-beyond}{{16.9}{60}{超出领头对数与有限 $N_c$ \tagP }{theorem.16.9}{}}
\newlabel{prop:jets-ngl-beyond@cref}{{[proposition][9][16]16.9}{[1][60][]60}}
\newlabel{sec:jets-sll}{{16.4}{60}{第 4 节（测试）}{section.16.4}{}}
\newlabel{sec:jets-sll@cref}{{[section][4][16]16.4}{[1][60][]60}}
\newlabel{thm:jets-ngl-factorization}{{16.10}{60}{喷注过程的有效场论因子化 \tagA }{theorem.16.10}{}}
\newlabel{thm:jets-ngl-factorization@cref}{{[theorem][10][16]16.10}{[1][60][]60}}
\newlabel{prop:jets-sll}{{16.11}{60}{超领头对数 \tagP }{theorem.16.11}{}}
\newlabel{prop:jets-sll@cref}{{[proposition][11][16]16.11}{[1][60][]60}}
\newlabel{ch:smallx}{{17}{61}{小 $x$ 与高能因子化}{chapter.17}{}}
\newlabel{ch:smallx@cref}{{[chapter][17][]17}{[1][61][]61}}
\newlabel{sec:smallx-regge}{{17.1}{61}{第 1 节（测试）}{section.17.1}{}}
\newlabel{sec:smallx-regge@cref}{{[section][1][17]17.1}{[1][61][]61}}
\newlabel{def:smallx-regge}{{17.1}{61}{高能极限}{theorem.17.1}{}}
\newlabel{def:smallx-regge@cref}{{[definition][1][17]17.1}{[1][61][]61}}
\newlabel{thm:smallx-bfkl}{{17.2}{61}{BFKL 方程 \tagA }{theorem.17.2}{}}
\newlabel{thm:smallx-bfkl@cref}{{[theorem][2][17]17.2}{[1][61][]61}}
\newlabel{sec:smallx-bfkl}{{17.2}{61}{第 2 节（测试）}{section.17.2}{}}
\newlabel{sec:smallx-bfkl@cref}{{[section][2][17]17.2}{[1][61][]61}}
\newlabel{prop:smallx-kt}{{17.3}{61}{$k_T$（高能）因子化 \tagP }{theorem.17.3}{}}
\newlabel{prop:smallx-kt@cref}{{[proposition][3][17]17.3}{[1][61][]61}}
\newlabel{thm:smallx-bk}{{17.4}{61}{BK 方程与 JIMWLK \tagA }{theorem.17.4}{}}
\newlabel{thm:smallx-bk@cref}{{[theorem][4][17]17.4}{[1][61][]61}}
\newlabel{sec:smallx-kt}{{17.3}{61}{第 3 节（测试）}{section.17.3}{}}
\newlabel{sec:smallx-kt@cref}{{[section][3][17]17.3}{[1][61][]61}}
\newlabel{conj:smallx-saturation}{{17.5}{61}{饱和与色玻璃凝聚 \tagC }{theorem.17.5}{}}
\newlabel{conj:smallx-saturation@cref}{{[conjecture][5][17]17.5}{[1][61][]61}}
\newlabel{rem:smallx-nonuniform}{{17.6}{62}{共线因子化在小 $x$ 的非一致性}{theorem.17.6}{}}
\newlabel{rem:smallx-nonuniform@cref}{{[remark][6][17]17.6}{[1][61][]62}}
\newlabel{sec:smallx-dipole}{{17.4}{62}{第 4 节（测试）}{section.17.4}{}}
\newlabel{sec:smallx-dipole@cref}{{[section][4][17]17.4}{[1][62][]62}}
\newlabel{ex:smallx-dla}{{17.7}{62}{DGLAP 与 BFKL 的双领头对数重叠}{theorem.17.7}{}}
\newlabel{ex:smallx-dla@cref}{{[example][7][17]17.7}{[1][62][]62}}
\newlabel{ch:break}{{18}{63}{因子化的破坏与非标准情形}{chapter.18}{}}
\newlabel{ch:break@cref}{{[chapter][18][]18}{[1][63][]63}}
\newlabel{sec:break-overview}{{18.1}{63}{第 1 节（测试）}{section.18.1}{}}
\newlabel{sec:break-overview@cref}{{[section][1][18]18.1}{[1][63][]63}}
\newlabel{def:break-generalized-tmd}{{18.1}{63}{标准与广义 TMD 因子化}{theorem.18.1}{}}
\newlabel{def:break-generalized-tmd@cref}{{[definition][1][18]18.1}{[1][63][]63}}
\newlabel{prop:break-measurement}{{18.2}{63}{测量相容性判据 \tagR }{theorem.18.2}{}}
\newlabel{prop:break-measurement@cref}{{[proposition][2][18]18.2}{[1][63][]63}}
\newlabel{rem:break-measurement-failures}{{18.3}{63}{失效的位置}{theorem.18.3}{}}
\newlabel{rem:break-measurement-failures@cref}{{[remark][3][18]18.3}{[1][63][]63}}
\newlabel{sec:break-tmd}{{18.2}{63}{第 2 节（测试）}{section.18.2}{}}
\newlabel{sec:break-tmd@cref}{{[section][2][18]18.2}{[1][63][]63}}
\newlabel{prop:break-colour}{{18.4}{63}{单个额外胶子的色因子（胶子极强度） \tagR }{theorem.18.4}{}}
\newlabel{prop:break-colour@cref}{{[proposition][4][18]18.4}{[1][63][]63}}
\newlabel{ex:break-colour}{{18.5}{63}{由投影算符得到 $C_G$ 表；阿贝尔极限；纠缠检验}{theorem.18.5}{}}
\newlabel{ex:break-colour@cref}{{[example][5][18]18.5}{[1][63][]63}}
\newlabel{thm:break-tmd}{{18.6}{63}{强子—强子末态强子的 TMD 因子化不成立 \tagV }{theorem.18.6}{}}
\newlabel{thm:break-tmd@cref}{{[theorem][6][18]18.6}{[1][63][]63}}
\newlabel{sec:break-spacelike}{{18.3}{63}{第 3 节（测试）}{section.18.3}{}}
\newlabel{sec:break-spacelike@cref}{{[section][3][18]18.3}{[1][63][]63}}
\newlabel{ex:break-c07}{{18.7}{63}{Collins 2007 的双胶子反常（阿贝尔标量模型）}{theorem.18.7}{}}
\newlabel{ex:break-c07@cref}{{[example][7][18]18.7}{[1][63][]63}}
\newlabel{prop:break-spacelike}{{18.8}{63}{类空共线极限中严格共线因子化的失效 \tagV }{theorem.18.8}{}}
\newlabel{prop:break-spacelike@cref}{{[proposition][8][18]18.8}{[1][63][]63}}
\newlabel{ex:break-spacelike}{{18.9}{64}{CdFR 单圈类空劈裂的色结构}{theorem.18.9}{}}
\newlabel{ex:break-spacelike@cref}{{[example][9][18]18.9}{[1][64][]64}}
\newlabel{sec:break-mpi}{{18.4}{64}{第 4 节（测试）}{section.18.4}{}}
\newlabel{sec:break-mpi@cref}{{[section][4][18]18.4}{[1][64][]64}}
\newlabel{prop:break-mass-singularity}{{18.10}{64}{超出 NNLO 的质量奇异性因子化 \tagC }{theorem.18.10}{}}
\newlabel{prop:break-mass-singularity@cref}{{[proposition][10][18]18.10}{[1][64][]64}}
\newlabel{prop:break-fss}{{18.11}{64}{Coulomb 相位与软交换不对易导致的共线因子化破坏 \tagV }{theorem.18.11}{}}
\newlabel{prop:break-fss@cref}{{[proposition][11][18]18.11}{[1][64][]64}}
\newlabel{prop:break-jet-veto}{{18.12}{64}{对多重部分子相互作用敏感的观测量 \tagV }{theorem.18.12}{}}
\newlabel{prop:break-jet-veto@cref}{{[proposition][12][18]18.12}{[1][64][]64}}
\newlabel{sec:break-diffraction}{{18.5}{64}{第 5 节（测试）}{section.18.5}{}}
\newlabel{sec:break-diffraction@cref}{{[section][5][18]18.5}{[1][64][]64}}
\newlabel{def:break-diffraction}{{18.13}{64}{衍射 DIS 与衍射部分子分布}{theorem.18.13}{}}
\newlabel{def:break-diffraction@cref}{{[definition][13][18]18.13}{[1][64][]64}}
\newlabel{thm:break-ddis}{{18.14}{64}{衍射 DIS 的因子化 \tagA }{theorem.18.14}{}}
\newlabel{thm:break-ddis@cref}{{[theorem][14][18]18.14}{[1][64][]64}}
\newlabel{prop:break-hh-diffraction}{{18.15}{64}{强子—强子硬衍射中的因子化破坏 \tagV }{theorem.18.15}{}}
\newlabel{prop:break-hh-diffraction@cref}{{[proposition][15][18]18.15}{[1][64][]64}}
\newlabel{sec:break-dps}{{18.6}{65}{第 6 节（测试）}{section.18.6}{}}
\newlabel{sec:break-dps@cref}{{[section][6][18]18.6}{[1][64][]65}}
\newlabel{def:break-dpdf}{{18.16}{65}{双部分子分布与双 Drell–Yan}{theorem.18.16}{}}
\newlabel{def:break-dpdf@cref}{{[definition][16][18]18.16}{[1][64][]65}}
\newlabel{thm:break-double-dy}{{18.17}{65}{双 Drell–Yan 的因子化 \tagA }{theorem.18.17}{}}
\newlabel{thm:break-double-dy@cref}{{[theorem][17][18]18.17}{[1][64][]65}}
\newlabel{conj:break-dps-coloured}{{18.18}{65}{有色末态的双部分子散射 \tagC }{theorem.18.18}{}}
\newlabel{conj:break-dps-coloured@cref}{{[conjecture][18][18]18.18}{[1][65][]65}}
\newlabel{sec:break-summary}{{18.7}{65}{第 7 节（测试）}{section.18.7}{}}
\newlabel{sec:break-summary@cref}{{[section][7][18]18.7}{[1][65][]65}}
\newlabel{conj:break-pocket}{{18.19}{65}{"口袋公式" $\sigma _{\rm DPS}=\sigma _A\sigma _B/\sigma _{\rm eff}$ \tagC }{theorem.18.19}{}}
\newlabel{conj:break-pocket@cref}{{[conjecture][19][18]18.19}{[1][65][]65}}
\newlabel{ch:excl}{{19}{67}{硬遍举过程}{chapter.19}{}}
\newlabel{ch:excl@cref}{{[chapter][19][]19}{[1][67][]67}}
\newlabel{sec:excl-da}{{19.1}{67}{第 1 节（测试）}{section.19.1}{}}
\newlabel{sec:excl-da@cref}{{[section][1][19]19.1}{[1][67][]67}}
\newlabel{def:excl-da}{{19.1}{67}{介子分布振幅}{theorem.19.1}{}}
\newlabel{def:excl-da@cref}{{[definition][1][19]19.1}{[1][67][]67}}
\newlabel{prop:excl-erbl}{{19.2}{67}{单圈 ERBL 演化 \tagR }{theorem.19.2}{}}
\newlabel{prop:excl-erbl@cref}{{[proposition][2][19]19.2}{[1][67][]67}}
\newlabel{prop:excl-soft-cancel}{{19.3}{67}{色单态介子的软胶子在被积函数层次相消 \tagR }{theorem.19.3}{}}
\newlabel{prop:excl-soft-cancel@cref}{{[proposition][3][19]19.3}{[1][67][]67}}
\newlabel{sec:excl-soft}{{19.2}{67}{第 2 节（测试）}{section.19.2}{}}
\newlabel{sec:excl-soft@cref}{{[section][2][19]19.2}{[1][67][]67}}
\newlabel{thm:excl-pion-ff}{{19.4}{67}{介子电磁形状因子的因子化 \tagA }{theorem.19.4}{}}
\newlabel{thm:excl-pion-ff@cref}{{[theorem][4][19]19.4}{[1][67][]67}}
\newlabel{prop:excl-pion-lo}{{19.5}{67}{领头阶硬核 \tagR }{theorem.19.5}{}}
\newlabel{prop:excl-pion-lo@cref}{{[proposition][5][19]19.5}{[1][67][]67}}
\newlabel{ex:excl-pion-ff-lo}{{19.6}{67}{领头阶 $\pi $ 形状因子}{theorem.19.6}{}}
\newlabel{ex:excl-pion-ff-lo@cref}{{[example][6][19]19.6}{[1][67][]67}}
\newlabel{sec:excl-ff}{{19.3}{68}{第 3 节（测试）}{section.19.3}{}}
\newlabel{sec:excl-ff@cref}{{[section][3][19]19.3}{[1][67][]68}}
\newlabel{thm:excl-transition}{{19.7}{68}{$\gamma ^*\gamma \to \pi ^0$ 形状因子 \tagA }{theorem.19.7}{}}
\newlabel{thm:excl-transition@cref}{{[theorem][7][19]19.7}{[1][67][]68}}
\newlabel{prop:excl-endpoint}{{19.8}{68}{端点区域与软重叠 \tagP }{theorem.19.8}{}}
\newlabel{prop:excl-endpoint@cref}{{[proposition][8][19]19.8}{[1][68][]68}}
\newlabel{def:excl-gpd}{{19.9}{68}{广义部分子分布}{theorem.19.9}{}}
\newlabel{def:excl-gpd@cref}{{[definition][9][19]19.9}{[1][68][]68}}
\newlabel{sec:excl-endpoint}{{19.4}{68}{第 4 节（测试）}{section.19.4}{}}
\newlabel{sec:excl-endpoint@cref}{{[section][4][19]19.4}{[1][68][]68}}
\newlabel{prop:excl-polynomiality}{{19.10}{68}{多项式性与 Ji 求和规则 \tagR }{theorem.19.10}{}}
\newlabel{prop:excl-polynomiality@cref}{{[proposition][10][19]19.10}{[1][68][]68}}
\newlabel{thm:excl-dvcs}{{19.11}{68}{DVCS 的因子化 \tagA }{theorem.19.11}{}}
\newlabel{thm:excl-dvcs@cref}{{[theorem][11][19]19.11}{[1][68][]68}}
\newlabel{prop:excl-dvcs-lo}{{19.12}{68}{DVCS 领头阶系数函数 \tagR }{theorem.19.12}{}}
\newlabel{prop:excl-dvcs-lo@cref}{{[proposition][12][19]19.12}{[1][68][]68}}
\newlabel{sec:excl-gpd}{{19.5}{68}{第 5 节（测试）}{section.19.5}{}}
\newlabel{sec:excl-gpd@cref}{{[section][5][19]19.5}{[1][68][]68}}
\newlabel{thm:excl-meson-L}{{19.13}{68}{纵向光子介子产生的因子化 \tagA }{theorem.19.13}{}}
\newlabel{thm:excl-meson-L@cref}{{[theorem][13][19]19.13}{[1][68][]68}}
\newlabel{prop:excl-meson-T}{{19.14}{69}{横向光子：标准共线因子化失效 \tagV }{theorem.19.14}{}}
\newlabel{prop:excl-meson-T@cref}{{[proposition][14][19]19.14}{[1][69][]69}}
\newlabel{rem:excl-hadron-scattering}{{19.15}{69}{宽角强子散射}{theorem.19.15}{}}
\newlabel{rem:excl-hadron-scattering@cref}{{[remark][15][19]19.15}{[1][69][]69}}
\newlabel{sec:excl-dvcs}{{19.6}{69}{第 6 节（测试）}{section.19.6}{}}
\newlabel{sec:excl-dvcs@cref}{{[section][6][19]19.6}{[1][69][]69}}
\newlabel{prop:excl-near-threshold}{{19.16}{69}{近阈重夸克偶素光生（重夸克极限） \tagP }{theorem.19.16}{}}
\newlabel{prop:excl-near-threshold@cref}{{[proposition][16][19]19.16}{[1][69][]69}}
\newlabel{ch:heavy}{{20}{70}{重夸克物理：$B$ 介子衰变与重夸克偶素}{chapter.20}{}}
\newlabel{ch:heavy@cref}{{[chapter][20][]20}{[1][70][]70}}
\newlabel{sec:heavy-bdpi}{{20.1}{70}{第 1 节（测试）}{section.20.1}{}}
\newlabel{sec:heavy-bdpi@cref}{{[section][1][20]20.1}{[1][70][]70}}
\newlabel{thm:heavy-bdpi}{{20.1}{70}{$B\to D\pi $ 的全阶因子化 \tagA }{theorem.20.1}{}}
\newlabel{thm:heavy-bdpi@cref}{{[theorem][1][20]20.1}{[1][70][]70}}
\newlabel{prop:heavy-btopi-ff}{{20.2}{70}{重到轻形状因子：软重叠与硬旁观者 \tagA }{theorem.20.2}{}}
\newlabel{prop:heavy-btopi-ff@cref}{{[proposition][2][20]20.2}{[1][70][]70}}
\newlabel{prop:heavy-qcdf-lp}{{20.3}{70}{无粲两体衰变的领头幂次 QCD 因子化 \tagP }{theorem.20.3}{}}
\newlabel{prop:heavy-qcdf-lp@cref}{{[proposition][3][20]20.3}{[1][70][]70}}
\newlabel{sec:heavy-qcdf}{{20.2}{70}{第 2 节（测试）}{section.20.2}{}}
\newlabel{sec:heavy-qcdf@cref}{{[section][2][20]20.2}{[1][70][]70}}
\newlabel{prop:heavy-qcdf-endpoint}{{20.4}{70}{幂次修正中的端点发散 \tagV }{theorem.20.4}{}}
\newlabel{prop:heavy-qcdf-endpoint@cref}{{[proposition][4][20]20.4}{[1][70][]70}}
\newlabel{thm:heavy-inclusive-b}{{20.5}{70}{单举 $B$ 衰变的重夸克展开 \tagA }{theorem.20.5}{}}
\newlabel{thm:heavy-inclusive-b@cref}{{[theorem][5][20]20.5}{[1][70][]70}}
\newlabel{def:heavy-nrqcd}{{20.6}{71}{NRQCD 及其幂次计数}{theorem.20.6}{}}
\newlabel{def:heavy-nrqcd@cref}{{[definition][6][20]20.6}{[1][71][]71}}
\newlabel{sec:heavy-inclusive}{{20.3}{71}{第 3 节（测试）}{section.20.3}{}}
\newlabel{sec:heavy-inclusive@cref}{{[section][3][20]20.3}{[1][71][]71}}
\newlabel{thm:heavy-nrqcd-decay}{{20.7}{71}{衰变的 NRQCD 因子化 \tagA }{theorem.20.7}{}}
\newlabel{thm:heavy-nrqcd-decay@cref}{{[theorem][7][20]20.7}{[1][71][]71}}
\newlabel{prop:heavy-nrqcd-matching}{{20.8}{71}{匹配系数 \tagR }{theorem.20.8}{}}
\newlabel{prop:heavy-nrqcd-matching@cref}{{[proposition][8][20]20.8}{[1][71][]71}}
\newlabel{ex:heavy-cv}{{20.9}{71}{阈值处的单圈 $c_v$}{theorem.20.9}{}}
\newlabel{ex:heavy-cv@cref}{{[example][9][20]20.9}{[1][71][]71}}
\newlabel{sec:heavy-nrqcd}{{20.4}{71}{第 4 节（测试）}{section.20.4}{}}
\newlabel{sec:heavy-nrqcd@cref}{{[section][4][20]20.4}{[1][71][]71}}
\newlabel{conj:heavy-nrqcd-production}{{20.10}{71}{产生的 NRQCD 因子化 \tagC }{theorem.20.10}{}}
\newlabel{conj:heavy-nrqcd-production@cref}{{[conjecture][10][20]20.10}{[1][71][]71}}
\newlabel{prop:heavy-quarkonium-fragmentation}{{20.11}{71}{大 $p_T$ 重夸克偶素产生 \tagA }{theorem.20.11}{}}
\newlabel{prop:heavy-quarkonium-fragmentation@cref}{{[proposition][11][20]20.11}{[1][71][]71}}
\newlabel{ch:power}{{21}{73}{幂次修正、次领头幂次、重整化子与非微扰地位}{chapter.21}{}}
\newlabel{ch:power@cref}{{[chapter][21][]21}{[1][73][]73}}
\newlabel{sec:power-twist4}{{21.1}{73}{第 1 节（测试）}{section.21.1}{}}
\newlabel{sec:power-twist4@cref}{{[section][1][21]21.1}{[1][73][]73}}
\newlabel{prop:power-dis-twist4-x}{{21.1}{73}{DIS 扭度 4 修正的 $x$ 空间形式 \tagA }{theorem.21.1}{}}
\newlabel{prop:power-dis-twist4-x@cref}{{[proposition][1][21]21.1}{[1][73][]73}}
\newlabel{prop:power-twist3-sidis}{{21.2}{73}{SIDIS 型过程中的扭度 3 共线因子化 \tagP }{theorem.21.2}{}}
\newlabel{prop:power-twist3-sidis@cref}{{[proposition][2][21]21.2}{[1][73][]73}}
\newlabel{conj:power-twist3-pp}{{21.3}{73}{强子—强子碰撞中单举强子单自旋不对称的扭度 3 因子化 \tagC }{theorem.21.3}{}}
\newlabel{conj:power-twist3-pp@cref}{{[conjecture][3][21]21.3}{[1][73][]73}}
\newlabel{sec:power-twist3}{{21.2}{73}{第 2 节（测试）}{section.21.2}{}}
\newlabel{sec:power-twist3@cref}{{[section][2][21]21.2}{[1][73][]73}}
\newlabel{ex:power-sgp}{{21.4}{73}{最低阶的软胶子极点机制}{theorem.21.4}{}}
\newlabel{ex:power-sgp@cref}{{[example][4][21]21.4}{[1][73][]73}}
\newlabel{prop:power-next-to-soft}{{21.5}{73}{次软定理 \tagR }{theorem.21.5}{}}
\newlabel{prop:power-next-to-soft@cref}{{[proposition][5][21]21.5}{[1][73][]73}}
\newlabel{prop:power-nlp-endpoint}{{21.6}{73}{次领头幂次因子化与端点发散 \tagP }{theorem.21.6}{}}
\newlabel{prop:power-nlp-endpoint@cref}{{[proposition][6][21]21.6}{[1][73][]73}}
\newlabel{sec:power-nlp}{{21.3}{74}{第 3 节（测试）}{section.21.3}{}}
\newlabel{sec:power-nlp@cref}{{[section][3][21]21.3}{[1][74][]74}}
\newlabel{prop:power-rearrangement}{{21.7}{74}{OPE 基的重排自由 \tagR }{theorem.21.7}{}}
\newlabel{prop:power-rearrangement@cref}{{[proposition][7][21]21.7}{[1][74][]74}}
\newlabel{rem:power-practical-ope}{{21.8}{74}{"实用 OPE" 与重整化子吸收的两个猜想性成分}{theorem.21.8}{}}
\newlabel{rem:power-practical-ope@cref}{{[remark][8][21]21.8}{[1][74][]74}}
\newlabel{def:power-borel}{{21.9}{74}{Borel 变换与重整化子}{theorem.21.9}{}}
\newlabel{def:power-borel@cref}{{[definition][9][21]21.9}{[1][74][]74}}
\newlabel{sec:power-ope-basis}{{21.4}{74}{第 4 节（测试）}{section.21.4}{}}
\newlabel{sec:power-ope-basis@cref}{{[section][4][21]21.4}{[1][74][]74}}
\newlabel{prop:power-adler-large-beta0}{{21.10}{74}{大 $\beta _0$ 极限下的 Adler 函数 \tagP }{theorem.21.10}{}}
\newlabel{prop:power-adler-large-beta0@cref}{{[proposition][10][21]21.10}{[1][74][]74}}
\newlabel{ex:power-renormalon-chain}{{21.11}{74}{带 $n$ 个费米子泡的单胶子}{theorem.21.11}{}}
\newlabel{ex:power-renormalon-chain@cref}{{[example][11][21]21.11}{[1][74][]74}}
\newlabel{prop:power-renormalon-examples}{{21.12}{74}{若干过程中的重整化子与幂次修正 \tagP }{theorem.21.12}{}}
\newlabel{prop:power-renormalon-examples@cref}{{[proposition][12][21]21.12}{[1][74][]74}}
\newlabel{sec:power-renormalons}{{21.5}{74}{第 5 节（测试）}{section.21.5}{}}
\newlabel{sec:power-renormalons@cref}{{[section][5][21]21.5}{[1][74][]74}}
\newlabel{conj:power-renormalon-ht}{{21.13}{74}{红外重整化子与高扭度的一般对应 \tagC }{theorem.21.13}{}}
\newlabel{conj:power-renormalon-ht@cref}{{[conjecture][13][21]21.13}{[1][74][]74}}
\newlabel{rem:power-np-status}{{21.14}{74}{非微扰地位}{theorem.21.14}{}}
\newlabel{rem:power-np-status@cref}{{[remark][14][21]21.14}{[1][74][]74}}
\newlabel{thm:power-ope-convergence}{{21.15}{74}{微扰 $\varphi ^4_4$ 中 OPE 的收敛性 \tagR }{theorem.21.15}{}}
\newlabel{thm:power-ope-convergence@cref}{{[theorem][15][21]21.15}{[1][74][]74}}
\newlabel{sec:power-np}{{21.6}{74}{第 6 节（测试）}{section.21.6}{}}
\newlabel{sec:power-np@cref}{{[section][6][21]21.6}{[1][74][]74}}
\newlabel{prop:power-thooft-model}{{21.16}{74}{'t Hooft 模型 \tagP }{theorem.21.16}{}}
\newlabel{prop:power-thooft-model@cref}{{[proposition][16][21]21.16}{[1][74][]74}}
\newlabel{ch:lamet}{{22}{77}{LaMET 与格点因子化：季向东等人的因子化证明}{chapter.22}{}}
\newlabel{ch:lamet@cref}{{[chapter][22][]22}{[1][77][]77}}
\newlabel{sec:lamet-intro}{{22.1}{77}{第 1 节（测试）}{section.22.1}{}}
\newlabel{sec:lamet-intro@cref}{{[section][1][22]22.1}{[1][77][]77}}
\newlabel{conj:lamet-principle}{{22.1}{77}{LaMET 原理（非微扰形式） \tagC }{theorem.22.1}{}}
\newlabel{conj:lamet-principle@cref}{{[conjecture][1][22]22.1}{[1][77][]77}}
\newlabel{rem:lamet-formulations}{{22.2}{77}{LaMET 原理的几种表述及其可证明的内容}{theorem.22.2}{}}
\newlabel{rem:lamet-formulations@cref}{{[remark][2][22]22.2}{[1][77][]77}}
\newlabel{prop:lattice-moments}{{22.3}{77}{格点可计算的矩 \tagR }{theorem.22.3}{}}
\newlabel{prop:lattice-moments@cref}{{[proposition][3][22]22.3}{[1][77][]77}}
\newlabel{def:lattice-quasi}{{22.4}{77}{准分布、赝分布与格点截面}{theorem.22.4}{}}
\newlabel{def:lattice-quasi@cref}{{[definition][4][22]22.4}{[1][77][]77}}
\newlabel{prop:lamet-support}{{22.5}{77}{准关联函数的支撑性质（逐图） \tagR }{theorem.22.5}{}}
\newlabel{prop:lamet-support@cref}{{[proposition][5][22]22.5}{[1][77][]77}}
\newlabel{sec:lamet-principle}{{22.2}{77}{第 2 节（测试）}{section.22.2}{}}
\newlabel{sec:lamet-principle@cref}{{[section][2][22]22.2}{[1][77][]77}}
\newlabel{cor:lamet-entire}{{22.6}{77}{Ioffe 时间分布是指数型整函数 \tagR }{theorem.22.6}{}}
\newlabel{cor:lamet-entire@cref}{{[corollary][6][22]22.6}{[1][77][]77}}
\newlabel{ex:lattice-quasi-tree}{{22.7}{77}{自由夸克的准 PDF：树图与单圈非物理区}{theorem.22.7}{}}
\newlabel{ex:lattice-quasi-tree@cref}{{[example][7][22]22.7}{[1][77][]77}}
\newlabel{asm:lamet-aux-renorm}{{22.8}{78}{含类空辅助静态色源的 QCD 的可重整性 \tagA }{theorem.22.8}{}}
\newlabel{asm:lamet-aux-renorm@cref}{{[assumption][8][22]22.8}{[1][77][]78}}
\newlabel{prop:lattice-renormalizability}{{22.9}{78}{Wilson 线夸克双线性算符的乘法可重整性 \tagA }{theorem.22.9}{}}
\newlabel{prop:lattice-renormalizability@cref}{{[proposition][9][22]22.9}{[1][78][]78}}
\newlabel{prop:lamet-ad-spacelike}{{22.10}{78}{类空与类时 Wilson 线算符的反常维度相同 \tagA }{theorem.22.10}{}}
\newlabel{prop:lamet-ad-spacelike@cref}{{[proposition][10][22]22.10}{[1][78][]78}}
\newlabel{sec:lattice-moments}{{22.3}{78}{第 3 节（测试）}{section.22.3}{}}
\newlabel{sec:lattice-moments@cref}{{[section][3][22]22.3}{[1][78][]78}}
\newlabel{def:lamet-schemes}{{22.11}{78}{准关联函数的重整化方案}{theorem.22.11}{}}
\newlabel{def:lamet-schemes@cref}{{[definition][11][22]22.11}{[1][78][]78}}
\newlabel{lem:lamet-lorentzian}{{22.12}{78}{残余线性项 $e^{-m_0|z|}$ 的 Lorentz 抹平}{theorem.22.12}{}}
\newlabel{lem:lamet-lorentzian@cref}{{[lemma][12][22]22.12}{[1][78][]78}}
\newlabel{prop:lamet-lrr}{{22.13}{78}{Wilson 线自能的线性重整子与领头重整子重求和 \tagP }{theorem.22.13}{}}
\newlabel{prop:lamet-lrr@cref}{{[proposition][13][22]22.13}{[1][78][]78}}
\newlabel{prop:lamet-euclid-minkowski}{{22.14}{78}{欧氏格点矩阵元与 Minkowski 矩阵元的等同 \tagA }{theorem.22.14}{}}
\newlabel{prop:lamet-euclid-minkowski@cref}{{[proposition][14][22]22.14}{[1][78][]78}}
\newlabel{prop:lamet-oneloop}{{22.15}{79}{准 PDF 匹配的单圈结果 \tagR }{theorem.22.15}{}}
\newlabel{prop:lamet-oneloop@cref}{{[proposition][15][22]22.15}{[1][78][]79}}
\newlabel{sec:lattice-quasi}{{22.4}{79}{第 4 节（测试）}{section.22.4}{}}
\newlabel{sec:lattice-quasi@cref}{{[section][4][22]22.4}{[1][79][]79}}
\newlabel{prop:lamet-tails}{{22.16}{79}{单圈准 PDF 的尾部、矩与卷积定义域 \tagR }{theorem.22.16}{}}
\newlabel{prop:lamet-tails@cref}{{[proposition][16][22]22.16}{[1][79][]79}}
\newlabel{prop:lamet-twoloop}{{22.17}{79}{非单态准 PDF 匹配核的两圈结果 \tagR }{theorem.22.17}{}}
\newlabel{prop:lamet-twoloop@cref}{{[proposition][17][22]22.17}{[1][79][]79}}
\newlabel{prop:lamet-polarized}{{22.18}{79}{极化、横度与扭度 3 准分布的匹配 \tagP }{theorem.22.18}{}}
\newlabel{prop:lamet-polarized@cref}{{[proposition][18][22]22.18}{[1][79][]79}}
\newlabel{asm:lamet-ope}{{22.19}{79}{类空 Wilson 线算符的短距离 OPE（重整化微扰论） \tagA }{theorem.22.19}{}}
\newlabel{asm:lamet-ope@cref}{{[assumption][19][22]22.19}{[1][79][]79}}
\newlabel{thm:lamet-coordinate}{{22.20}{79}{坐标空间因子化（赝分布形式） \tagR }{theorem.22.20}{}}
\newlabel{thm:lamet-coordinate@cref}{{[theorem][20][22]22.20}{[1][79][]79}}
\newlabel{sec:lamet-renorm}{{22.5}{79}{第 5 节（测试）}{section.22.5}{}}
\newlabel{sec:lamet-renorm@cref}{{[section][5][22]22.5}{[1][79][]79}}
\newlabel{thm:lamet-parton-momentum}{{22.21}{79}{动量空间形式与部分子动量 \tagR }{theorem.22.21}{}}
\newlabel{thm:lamet-parton-momentum@cref}{{[theorem][21][22]22.21}{[1][79][]79}}
\newlabel{rem:lamet-early-form}{{22.22}{79}{以强子动量为变量的早期匹配形式不成立 \tagV }{theorem.22.22}{}}
\newlabel{rem:lamet-early-form@cref}{{[remark][22][22]22.22}{[1][79][]79}}
\newlabel{prop:lamet-cfp}{{22.23}{80}{光锥规范下的共线减除路线（Ma–Qiu） \tagA }{theorem.22.23}{}}
\newlabel{prop:lamet-cfp@cref}{{[proposition][23][22]22.23}{[1][79][]80}}
\newlabel{thm:lamet-lcs}{{22.24}{80}{“好的格点截面”的因子化（Ma–Qiu） \tagA }{theorem.22.24}{}}
\newlabel{thm:lamet-lcs@cref}{{[theorem][24][22]22.24}{[1][80][]80}}
\newlabel{thm:lattice-lamet}{{22.25}{80}{LaMET 匹配定理（部分子态、领头幂次、全阶） \tagA }{theorem.22.25}{}}
\newlabel{thm:lattice-lamet@cref}{{[theorem][25][22]22.25}{[1][80][]80}}
\newlabel{sec:lamet-oneloop}{{22.6}{80}{第 6 节（测试）}{section.22.6}{}}
\newlabel{sec:lamet-oneloop@cref}{{[section][6][22]22.6}{[1][80][]80}}
\newlabel{prop:lamet-gluon-singlet}{{22.26}{80}{单态夸克与胶子准 PDF：带混合的匹配 \tagA }{theorem.22.26}{}}
\newlabel{prop:lamet-gluon-singlet@cref}{{[proposition][26][22]22.26}{[1][80][]80}}
\newlabel{prop:lamet-hybrid-kernel}{{22.27}{80}{混合方案的单圈匹配核 \tagR }{theorem.22.27}{}}
\newlabel{prop:lamet-hybrid-kernel@cref}{{[proposition][27][22]22.27}{[1][80][]80}}
\newlabel{ex:lamet-hybrid-normalization}{{22.28}{80}{混合核的归一化}{theorem.22.28}{}}
\newlabel{ex:lamet-hybrid-normalization@cref}{{[example][28][22]22.28}{[1][80][]80}}
\newlabel{prop:lamet-scheme-validity}{{22.29}{80}{动量空间匹配在何种重整化方案中成立 \tagA }{theorem.22.29}{}}
\newlabel{prop:lamet-scheme-validity@cref}{{[proposition][29][22]22.29}{[1][80][]80}}
\newlabel{rem:lamet-ratio-large-z}{{22.30}{81}{比值与 RI/MOM 方案在大 $|z|$ 处的动量空间匹配 \tagC }{theorem.22.30}{}}
\newlabel{rem:lamet-ratio-large-z@cref}{{[remark][30][22]22.30}{[1][81][]81}}
\newlabel{sec:lattice-matching}{{22.7}{81}{第 7 节（测试）}{section.22.7}{}}
\newlabel{sec:lattice-matching@cref}{{[section][7][22]22.7}{[1][81][]81}}
\newlabel{asm:lamet-gap}{{22.31}{81}{谱隙与边界值正则性 \tagC }{theorem.22.31}{}}
\newlabel{asm:lamet-gap@cref}{{[assumption][31][22]22.31}{[1][81][]81}}
\newlabel{thm:lamet-exp-decay}{{22.32}{81}{局域性蕴含类空关联函数的指数衰减（Ji–Liu–Su） \tagR }{theorem.22.32}{}}
\newlabel{thm:lamet-exp-decay@cref}{{[theorem][32][22]22.32}{[1][81][]81}}
\newlabel{conj:lamet-decay-qcd}{{22.33}{81}{QCD 准分布矩阵元的长程渐近形式 \tagC }{theorem.22.33}{}}
\newlabel{conj:lamet-decay-qcd@cref}{{[conjecture][33][22]22.33}{[1][81][]81}}
\newlabel{lem:lamet-strip}{{22.34}{81}{准 PDF 的解析带}{theorem.22.34}{}}
\newlabel{lem:lamet-strip@cref}{{[lemma][34][22]22.34}{[1][81][]81}}
\newlabel{prop:lamet-truncation}{{22.35}{81}{Fourier 截断误差界}{theorem.22.35}{}}
\newlabel{prop:lamet-truncation@cref}{{[proposition][35][22]22.35}{[1][81][]81}}
\newlabel{sec:lamet-schemes}{{22.8}{81}{第 8 节（测试）}{section.22.8}{}}
\newlabel{sec:lamet-schemes@cref}{{[section][8][22]22.8}{[1][81][]81}}
\newlabel{prop:lamet-slepian}{{22.36}{81}{有限 Ioffe 时间窗口内可稳定恢复的自由度}{theorem.22.36}{}}
\newlabel{prop:lamet-slepian@cref}{{[proposition][36][22]22.36}{[1][81][]81}}
\newlabel{ex:lamet-fourier-numbers}{{22.37}{81}{截断误差与可恢复模式数的数值}{theorem.22.37}{}}
\newlabel{ex:lamet-fourier-numbers@cref}{{[example][37][22]22.37}{[1][81][]81}}
\newlabel{prop:lamet-tmc}{{22.38}{81}{靶质量修正的精确形式 \tagR }{theorem.22.38}{}}
\newlabel{prop:lamet-tmc@cref}{{[proposition][38][22]22.38}{[1][81][]81}}
\newlabel{ex:lamet-tmc-moments}{{22.39}{82}{$K_N$ 与 $K^\perp _N$（$N\le 4$）}{theorem.22.39}{}}
\newlabel{ex:lamet-tmc-moments@cref}{{[example][39][22]22.39}{[1][81][]82}}
\newlabel{lem:lamet-fourier}{{22.40}{82}{Fourier 引理：$z^2\leftrightarrow \partial _x^2$}{theorem.22.40}{}}
\newlabel{lem:lamet-fourier@cref}{{[lemma][40][22]22.40}{[1][81][]82}}
\newlabel{sec:lamet-fourier}{{22.9}{82}{第 9 节（测试）}{section.22.9}{}}
\newlabel{sec:lamet-fourier@cref}{{[section][9][22]22.9}{[1][82][]82}}
\newlabel{cor:lamet-endpoint}{{22.41}{82}{端点二分}{theorem.22.41}{}}
\newlabel{cor:lamet-endpoint@cref}{{[corollary][41][22]22.41}{[1][82][]82}}
\newlabel{prop:lamet-bvz}{{22.42}{82}{扭度 4 修正的重整子模型（大 $\beta _0$） \tagP }{theorem.22.42}{}}
\newlabel{prop:lamet-bvz@cref}{{[proposition][42][22]22.42}{[1][82][]82}}
\newlabel{asm:lamet-uniform}{{22.43}{82}{扭度展开对 Ioffe 时间的一致性 \tagC }{theorem.22.43}{}}
\newlabel{asm:lamet-uniform@cref}{{[assumption][43][22]22.43}{[1][82][]82}}
\newlabel{asm:lamet-twist4}{{22.44}{82}{领头幂次修正的结构（扭度 4） \tagC }{theorem.22.44}{}}
\newlabel{asm:lamet-twist4@cref}{{[assumption][44][22]22.44}{[1][82][]82}}
\newlabel{asm:lamet-large-z}{{22.45}{82}{大距离区域的控制 \tagC }{theorem.22.45}{}}
\newlabel{asm:lamet-large-z@cref}{{[assumption][45][22]22.45}{[1][82][]82}}
\newlabel{sec:lamet-power}{{22.10}{82}{第 10 节（测试）}{section.22.10}{}}
\newlabel{sec:lamet-power@cref}{{[section][10][22]22.10}{[1][82][]82}}
\newlabel{thm:lamet-hadronic}{{22.46}{82}{强子层次的 LaMET 匹配与幂次修正 \tagR }{theorem.22.46}{}}
\newlabel{thm:lamet-hadronic@cref}{{[theorem][46][22]22.46}{[1][82][]82}}
\newlabel{rem:lamet-domain}{{22.47}{82}{LaMET 预言的适用域}{theorem.22.47}{}}
\newlabel{rem:lamet-domain@cref}{{[remark][47][22]22.47}{[1][82][]82}}
\newlabel{prop:lamet-sdf}{{22.48}{82}{LaMET 与短距离因子化：同一函数的两种展开 \tagR }{theorem.22.48}{}}
\newlabel{prop:lamet-sdf@cref}{{[proposition][48][22]22.48}{[1][82][]82}}
\newlabel{rem:lamet-inverse-problem}{{22.49}{83}{逆问题之争：数学陈述与经验问题}{theorem.22.49}{}}
\newlabel{rem:lamet-inverse-problem@cref}{{[remark][49][22]22.49}{[1][82][]83}}
\newlabel{def:lamet-quasi-gpd}{{22.50}{83}{准 GPD 与准 DA}{theorem.22.50}{}}
\newlabel{def:lamet-quasi-gpd@cref}{{[definition][50][22]22.50}{[1][82][]83}}
\newlabel{sec:lamet-sdf}{{22.11}{83}{第 11 节（测试）}{section.22.11}{}}
\newlabel{sec:lamet-sdf@cref}{{[section][11][22]22.11}{[1][83][]83}}
\newlabel{thm:lamet-quasi-gpd}{{22.51}{83}{准 GPD 与准 DA 的匹配 \tagA }{theorem.22.51}{}}
\newlabel{thm:lamet-quasi-gpd@cref}{{[theorem][51][22]22.51}{[1][83][]83}}
\newlabel{prop:lamet-other-objects}{{22.52}{83}{其他准分布：重子与重介子 LCDA、$\Delta G$、双部分子分布 \tagP }{theorem.22.52}{}}
\newlabel{prop:lamet-other-objects@cref}{{[proposition][52][22]22.52}{[1][83][]83}}
\newlabel{asm:lamet-staple}{{22.53}{83}{订书钉形类空 Wilson 线算符的重整化与 $L\to \infty $ 极限 \tagA }{theorem.22.53}{}}
\newlabel{asm:lamet-staple@cref}{{[assumption][53][22]22.53}{[1][83][]83}}
\newlabel{def:lamet-quasi-tmd}{{22.54}{83}{准 TMD}{theorem.22.54}{}}
\newlabel{def:lamet-quasi-tmd@cref}{{[definition][54][22]22.54}{[1][83][]83}}
\newlabel{def:lamet-soft-offlc}{{22.55}{83}{离光锥软函数（Ji–Liu–Liu 归一化）}{theorem.22.55}{}}
\newlabel{def:lamet-soft-offlc@cref}{{[definition][55][22]22.55}{[1][83][]83}}
\newlabel{sec:lamet-offforward}{{22.12}{83}{第 12 节（测试）}{section.22.12}{}}
\newlabel{sec:lamet-offforward@cref}{{[section][12][22]22.12}{[1][83][]83}}
\newlabel{asm:lamet-soft-rapidity}{{22.56}{83}{离光锥软函数的大快度结构 \tagA }{theorem.22.56}{}}
\newlabel{asm:lamet-soft-rapidity@cref}{{[assumption][56][22]22.56}{[1][83][]83}}
\newlabel{def:lamet-reduced-soft}{{22.57}{83}{约化软函数 $S_r$}{theorem.22.57}{}}
\newlabel{def:lamet-reduced-soft@cref}{{[definition][57][22]22.57}{[1][83][]83}}
\newlabel{asm:lamet-order-limits}{{22.58}{84}{极限次序 \tagA }{theorem.22.58}{}}
\newlabel{asm:lamet-order-limits@cref}{{[assumption][58][22]22.58}{[1][83][]84}}
\newlabel{thm:lamet-quasi-tmd}{{22.59}{84}{准 TMD 因子化 \tagA }{theorem.22.59}{}}
\newlabel{thm:lamet-quasi-tmd@cref}{{[theorem][59][22]22.59}{[1][84][]84}}
\newlabel{ex:lamet-qtmd-oneloop}{{22.60}{84}{单圈检验：指数中的因子 $\frac 12$}{theorem.22.60}{}}
\newlabel{ex:lamet-qtmd-oneloop@cref}{{[example][60][22]22.60}{[1][84][]84}}
\newlabel{sec:lattice-tmd}{{22.13}{84}{第 13 节（测试）}{section.22.13}{}}
\newlabel{sec:lattice-tmd@cref}{{[section][13][22]22.13}{[1][84][]84}}
\newlabel{prop:lattice-cs-kernel}{{22.61}{84}{格点 Collins–Soper 核 \tagA }{theorem.22.61}{}}
\newlabel{prop:lattice-cs-kernel@cref}{{[proposition][61][22]22.61}{[1][84][]84}}
\newlabel{prop:lamet-tmd-ratios}{{22.62}{84}{准 TMD 之比与 TMD 之比 \tagA }{theorem.22.62}{}}
\newlabel{prop:lamet-tmd-ratios@cref}{{[proposition][62][22]22.62}{[1][84][]84}}
\newlabel{cor:lamet-principle-pt}{{22.63}{84}{LaMET 原理的微扰内容 \tagA }{theorem.22.63}{}}
\newlabel{cor:lamet-principle-pt@cref}{{[corollary][63][22]22.63}{[1][84][]84}}
\newlabel{rem:lattice-critiques}{{22.64}{84}{批评与回应}{theorem.22.64}{}}
\newlabel{rem:lattice-critiques@cref}{{[remark][64][22]22.64}{[1][84][]84}}
\newlabel{rem:lamet-open}{{22.65}{84}{开放问题}{theorem.22.65}{}}
\newlabel{rem:lamet-open@cref}{{[remark][65][22]22.65}{[1][84][]84}}
\newlabel{ch:yzliu}{{23}{85}{Yizhuang Liu 的严格因子化研究}{chapter.23}{}}
\newlabel{ch:yzliu@cref}{{[chapter][23][]23}{[1][85][]85}}
\newlabel{sec:yzliu-intro}{{23.1}{85}{第 1 节（测试）}{section.23.1}{}}
\newlabel{sec:yzliu-intro@cref}{{[section][1][23]23.1}{[1][85][]85}}
\newlabel{lem:yzliu-partial-exp}{{23.1}{85}{部分指数化引理 \tagA }{theorem.23.1}{}}
\newlabel{lem:yzliu-partial-exp@cref}{{[lemma][1][23]23.1}{[1][85][]85}}
\newlabel{prop:yzliu-pinch-ultrasoft}{{23.2}{85}{捏点奇点、超软模式与 $L\to \infty $ 极限 \tagA }{theorem.23.2}{}}
\newlabel{prop:yzliu-pinch-ultrasoft@cref}{{[proposition][2][23]23.2}{[1][85][]85}}
\newlabel{asm:yzliu-heavy-pair}{{23.3}{85}{静态色单态 $Q\bar Q$ 对的谱假设 \tagC }{theorem.23.3}{}}
\newlabel{asm:yzliu-heavy-pair@cref}{{[assumption][3][23]23.3}{[1][85][]85}}
\newlabel{thm:yzliu-soft-hqet}{{23.4}{85}{离光锥软函数是快速色源对的形状因子 \tagA }{theorem.23.4}{}}
\newlabel{thm:yzliu-soft-hqet@cref}{{[theorem][4][23]23.4}{[1][85][]85}}
\newlabel{sec:yzliu-soft}{{23.2}{85}{第 2 节（测试）}{section.23.2}{}}
\newlabel{sec:yzliu-soft@cref}{{[section][2][23]23.2}{[1][85][]85}}
\newlabel{ex:yzliu-soft-oneloop}{{23.5}{85}{单圈离光锥软函数与 $S_r$}{theorem.23.5}{}}
\newlabel{ex:yzliu-soft-oneloop@cref}{{[example][5][23]23.5}{[1][85][]85}}
\newlabel{cor:yzliu-soft-form-factor}{{23.6}{86}{软函数是静态色源对的形状因子（非微扰形式） \tagR }{theorem.23.6}{}}
\newlabel{cor:yzliu-soft-form-factor@cref}{{[corollary][6][23]23.6}{[1][85][]86}}
\newlabel{def:yzliu-hyperbolic-angle}{{23.7}{86}{双曲角与六种构型}{theorem.23.7}{}}
\newlabel{def:yzliu-hyperbolic-angle@cref}{{[definition][7][23]23.7}{[1][86][]86}}
\newlabel{asm:yzliu-wightman-pt}{{23.8}{86}{维数正规化微扰论中的 Wightman 结构 \tagA }{theorem.23.8}{}}
\newlabel{asm:yzliu-wightman-pt@cref}{{[assumption][8][23]23.8}{[1][86][]86}}
\newlabel{sec:yzliu-analyticity}{{23.3}{86}{第 3 节（测试）}{section.23.3}{}}
\newlabel{sec:yzliu-analyticity@cref}{{[section][3][23]23.3}{[1][86][]86}}
\newlabel{thm:yzliu-analyticity}{{23.9}{86}{软函数对复双曲角的解析性 \tagA }{theorem.23.9}{}}
\newlabel{thm:yzliu-analyticity@cref}{{[theorem][9][23]23.9}{[1][86][]86}}
\newlabel{prop:yzliu-rapidity-linear}{{23.10}{86}{离光锥软函数中快度发散的因子化（线性性） \tagA }{theorem.23.10}{}}
\newlabel{prop:yzliu-rapidity-linear@cref}{{[proposition][10][23]23.10}{[1][86][]86}}
\newlabel{def:yzliu-intrinsic-soft-lc}{{23.11}{86}{本征软函数的光锥定义}{theorem.23.11}{}}
\newlabel{def:yzliu-intrinsic-soft-lc@cref}{{[definition][11][23]23.11}{[1][86][]86}}
\newlabel{thm:yzliu-intrinsic-soft}{{23.12}{86}{本征软函数：定义的等价、方案无关性与反常维度 \tagA }{theorem.23.12}{}}
\newlabel{thm:yzliu-intrinsic-soft@cref}{{[theorem][12][23]23.12}{[1][86][]86}}
\newlabel{sec:yzliu-rapidity}{{23.4}{87}{第 4 节（测试）}{section.23.4}{}}
\newlabel{sec:yzliu-rapidity@cref}{{[section][4][23]23.4}{[1][86][]87}}
\newlabel{prop:yzliu-sI-twoloop}{{23.13}{87}{$S_I$ 的两圈结构 \tagR }{theorem.23.13}{}}
\newlabel{prop:yzliu-sI-twoloop@cref}{{[proposition][13][23]23.13}{[1][86][]87}}
\newlabel{ex:yzliu-sI-rge}{{23.14}{87}{由 RGE 确定两圈 $L_b^2$ 系数}{theorem.23.14}{}}
\newlabel{ex:yzliu-sI-rge@cref}{{[example][14][23]23.14}{[1][87][]87}}
\newlabel{asm:yzliu-lightcone-uniform}{{23.15}{87}{光锥极限对 ${\rm Im}\,\theta $ 的一致性 \tagP }{theorem.23.15}{}}
\newlabel{asm:yzliu-lightcone-uniform@cref}{{[assumption][15][23]23.15}{[1][87][]87}}
\newlabel{prop:yzliu-universality-classes}{{23.16}{87}{软函数的普适类 \tagP }{theorem.23.16}{}}
\newlabel{prop:yzliu-universality-classes@cref}{{[proposition][16][23]23.16}{[1][87][]87}}
\newlabel{sec:yzliu-time-ordering}{{23.5}{87}{第 5 节（测试）}{section.23.5}{}}
\newlabel{sec:yzliu-time-ordering@cref}{{[section][5][23]23.5}{[1][87][]87}}
\newlabel{def:yzliu-exp-soft}{{23.17}{87}{指数正规子下的软因子}{theorem.23.17}{}}
\newlabel{def:yzliu-exp-soft@cref}{{[definition][17][23]23.17}{[1][87][]87}}
\newlabel{lem:yzliu-matrix-tree}{{23.18}{87}{坐标空间参数函数的正性（全子式矩阵树公式）}{theorem.23.18}{}}
\newlabel{lem:yzliu-matrix-tree@cref}{{[lemma][18][23]23.18}{[1][87][]87}}
\newlabel{lem:yzliu-final-state}{{23.19}{87}{末态求和恒等式}{theorem.23.19}{}}
\newlabel{lem:yzliu-final-state@cref}{{[lemma][19][23]23.19}{[1][87][]87}}
\newlabel{thm:yzliu-time-ordering}{{23.20}{87}{DY 形 TMD 软因子的时序无关性 \tagA }{theorem.23.20}{}}
\newlabel{thm:yzliu-time-ordering@cref}{{[theorem][20][23]23.20}{[1][87][]87}}
\newlabel{sec:yzliu-quasi-tmd}{{23.6}{87}{第 6 节（测试）}{section.23.6}{}}
\newlabel{sec:yzliu-quasi-tmd@cref}{{[section][6][23]23.6}{[1][87][]87}}
\newlabel{prop:yzliu-dy-sidis-exp}{{23.21}{87}{指数正规子下 DY 与 SIDIS 软因子相等 \tagP }{theorem.23.21}{}}
\newlabel{prop:yzliu-dy-sidis-exp@cref}{{[proposition][21][23]23.21}{[1][87][]87}}
\newlabel{conj:yzliu-complex-wilson}{{23.22}{87}{复时空中的 Wilson 圈 \tagC }{theorem.23.22}{}}
\newlabel{conj:yzliu-complex-wilson@cref}{{[conjecture][22][23]23.22}{[1][87][]87}}
\newlabel{prop:yzliu-qtmd-powercount}{{23.23}{88}{准 TMD 因子化的 Collins 式证明：单线记号下的幂次计数 \tagA }{theorem.23.23}{}}
\newlabel{prop:yzliu-qtmd-powercount@cref}{{[proposition][23][23]23.23}{[1][88][]88}}
\newlabel{prop:yzliu-hard-kernel}{{23.24}{88}{准 TMD 匹配核：RGE、重—轻分解与两圈结果 \tagA }{theorem.23.24}{}}
\newlabel{prop:yzliu-hard-kernel@cref}{{[proposition][24][23]23.24}{[1][88][]88}}
\newlabel{sec:yzliu-exclusive}{{23.7}{88}{第 7 节（测试）}{section.23.7}{}}
\newlabel{sec:yzliu-exclusive@cref}{{[section][7][23]23.7}{[1][88][]88}}
\newlabel{ex:yzliu-hl-decomposition}{{23.25}{88}{单圈 $H=C^+_{HL}C^-_{HL}$ 与相位}{theorem.23.25}{}}
\newlabel{ex:yzliu-hl-decomposition@cref}{{[example][25][23]23.25}{[1][88][]88}}
\newlabel{prop:yzliu-lfwf}{{23.26}{88}{准光前波函数振幅的因子化 \tagA }{theorem.23.26}{}}
\newlabel{prop:yzliu-lfwf@cref}{{[proposition][26][23]23.26}{[1][88][]88}}
\newlabel{prop:yzliu-meson-ff}{{23.27}{88}{大动量下的介子形状因子与 $S_I$ 的提取 \tagA }{theorem.23.27}{}}
\newlabel{prop:yzliu-meson-ff@cref}{{[proposition][27][23]23.27}{[1][88][]88}}
\newlabel{prop:yzliu-dy-quasi}{{23.28}{88}{以准 TMD 表示的小 $q_T$ Drell–Yan 截面 \tagA }{theorem.23.28}{}}
\newlabel{prop:yzliu-dy-quasi@cref}{{[proposition][28][23]23.28}{[1][88][]88}}
\newlabel{sec:yzliu-threshold}{{23.8}{89}{第 8 节（测试）}{section.23.8}{}}
\newlabel{sec:yzliu-threshold@cref}{{[section][8][23]23.8}{[1][88][]89}}
\newlabel{prop:yzliu-sivers}{{23.29}{89}{由准 TMD 之比得到 Sivers 不对称 \tagA }{theorem.23.29}{}}
\newlabel{prop:yzliu-sivers@cref}{{[proposition][29][23]23.29}{[1][88][]89}}
\newlabel{def:yzliu-hl-jet}{{23.30}{89}{重—轻 Sudakov 因子与类空喷注函数}{theorem.23.30}{}}
\newlabel{def:yzliu-hl-jet@cref}{{[definition][30][23]23.30}{[1][89][]89}}
\newlabel{thm:yzliu-jet-continuation}{{23.31}{89}{类空与类时重夸克喷注函数的解析延拓 \tagA }{theorem.23.31}{}}
\newlabel{thm:yzliu-jet-continuation@cref}{{[theorem][31][23]23.31}{[1][89][]89}}
\newlabel{thm:yzliu-threshold}{{23.32}{89}{准 PDF 匹配核的阈值因子化（Ji–Liu–Su） \tagA }{theorem.23.32}{}}
\newlabel{thm:yzliu-threshold@cref}{{[theorem][32][23]23.32}{[1][89][]89}}
\newlabel{sec:yzliu-renormalon}{{23.9}{89}{第 9 节（测试）}{section.23.9}{}}
\newlabel{sec:yzliu-renormalon@cref}{{[section][9][23]23.9}{[1][89][]89}}
\newlabel{ex:yzliu-threshold-oneloop}{{23.33}{89}{阈值因子化的单圈检验}{theorem.23.33}{}}
\newlabel{ex:yzliu-threshold-oneloop@cref}{{[example][33][23]23.33}{[1][89][]89}}
\newlabel{prop:yzliu-ad-algebra}{{23.34}{89}{反常维度代数 \tagA }{theorem.23.34}{}}
\newlabel{prop:yzliu-ad-algebra@cref}{{[proposition][34][23]23.34}{[1][89][]89}}
\newlabel{prop:yzliu-threshold-nnlo}{{23.35}{89}{阈值因子化的 NNLO 检验 \tagP }{theorem.23.35}{}}
\newlabel{prop:yzliu-threshold-nnlo@cref}{{[proposition][35][23]23.35}{[1][89][]89}}
\newlabel{rem:yzliu-threshold-domain}{{23.36}{89}{阈值重求和与大 $x$ 适用域}{theorem.23.36}{}}
\newlabel{rem:yzliu-threshold-domain@cref}{{[remark][36][23]23.36}{[1][89][]89}}
\newlabel{sec:yzliu-gpd}{{23.10}{90}{第 10 节（测试）}{section.23.10}{}}
\newlabel{sec:yzliu-gpd@cref}{{[section][10][23]23.10}{[1][89][]90}}
\newlabel{prop:yzliu-linear-renormalon}{{23.37}{90}{重—轻 Sudakov 核的线性红外重整子 \tagP }{theorem.23.37}{}}
\newlabel{prop:yzliu-linear-renormalon@cref}{{[proposition][37][23]23.37}{[1][89][]90}}
\newlabel{conj:yzliu-renormalon-norm}{{23.38}{90}{线性重整子归一化的普适性 \tagC }{theorem.23.38}{}}
\newlabel{conj:yzliu-renormalon-norm@cref}{{[conjecture][38][23]23.38}{[1][90][]90}}
\newlabel{conj:yzliu-threshold-hadronic}{{23.39}{90}{强子夸克双线性矩阵元的阈值渐近 \tagC }{theorem.23.39}{}}
\newlabel{conj:yzliu-threshold-hadronic@cref}{{[conjecture][39][23]23.39}{[1][90][]90}}
\newlabel{prop:yzliu-gpd-oneloop}{{23.40}{90}{准 GPD 单圈匹配核在 $x=\pm \xi $ 附近的结构 \tagR }{theorem.23.40}{}}
\newlabel{prop:yzliu-gpd-oneloop@cref}{{[proposition][40][23]23.40}{[1][90][]90}}
\newlabel{sec:yzliu-models}{{23.11}{90}{第 11 节（测试）}{section.23.11}{}}
\newlabel{sec:yzliu-models@cref}{{[section][11][23]23.11}{[1][90][]90}}
\newlabel{conj:yzliu-gpd-xi}{{23.41}{90}{GPD 在 $x=\pm \xi $ 附近的 LaMET 可计算性 \tagC }{theorem.23.41}{}}
\newlabel{conj:yzliu-gpd-xi@cref}{{[conjecture][41][23]23.41}{[1][90][]90}}
\newlabel{thm:yzliu-gn-convergence}{{23.42}{90}{二维 $U(N)$ Gross–Neveu 模型中扭度展开的绝对收敛 \tagP }{theorem.23.42}{}}
\newlabel{thm:yzliu-gn-convergence@cref}{{[theorem][42][23]23.42}{[1][90][]90}}
\newlabel{prop:yzliu-gn-structure}{{23.43}{90}{Gross–Neveu 模型中扭度展开的结构：重整子相消与真空凝聚 \tagP }{theorem.23.43}{}}
\newlabel{prop:yzliu-gn-structure@cref}{{[proposition][43][23]23.43}{[1][90][]90}}
\newlabel{prop:yzliu-on-factorial}{{23.44}{90}{超可重整二维 $O(N)$ 模型：动量空间 OPE 的阶乘增强 \tagP }{theorem.23.44}{}}
\newlabel{prop:yzliu-on-factorial@cref}{{[proposition][44][23]23.44}{[1][90][]90}}
\newlabel{sec:yzliu-decay}{{23.12}{91}{第 12 节（测试）}{section.23.12}{}}
\newlabel{sec:yzliu-decay@cref}{{[section][12][23]23.12}{[1][91][]91}}
\newlabel{prop:yzliu-nlsm}{{23.45}{91}{二维非线性 $\sigma $ 模型：凝聚与跨级数 \tagP }{theorem.23.45}{}}
\newlabel{prop:yzliu-nlsm@cref}{{[proposition][45][23]23.45}{[1][91][]91}}
\newlabel{prop:yzliu-2dqcd}{{23.46}{91}{二维大 $N_c$ QCD 中的准 PDF \tagP }{theorem.23.46}{}}
\newlabel{prop:yzliu-2dqcd@cref}{{[proposition][46][23]23.46}{[1][91][]91}}
\newlabel{rem:yzliu-transforms}{{23.47}{91}{坐标空间与动量空间之间变换的教训}{theorem.23.47}{}}
\newlabel{rem:yzliu-transforms@cref}{{[remark][47][23]23.47}{[1][91][]91}}
\newlabel{rem:yzliu-instantons}{{23.48}{91}{瞬子与模型研究}{theorem.23.48}{}}
\newlabel{rem:yzliu-instantons@cref}{{[remark][48][23]23.48}{[1][91][]91}}
\newlabel{sec:yzliu-status}{{23.13}{91}{第 13 节（测试）}{section.23.13}{}}
\newlabel{sec:yzliu-status@cref}{{[section][13][23]23.13}{[1][91][]91}}
\newlabel{prop:yzliu-decay-extensions}{{23.49}{91}{标量场论中的谱表示与指数衰减（微扰论各阶） \tagP }{theorem.23.49}{}}
\newlabel{prop:yzliu-decay-extensions@cref}{{[proposition][49][23]23.49}{[1][91][]91}}
\newlabel{rem:yzliu-relation-css}{{23.50}{91}{与 CSS、Collins 与 SCET 的关系}{theorem.23.50}{}}
\newlabel{rem:yzliu-relation-css@cref}{{[remark][50][23]23.50}{[1][91][]91}}
\newlabel{rem:yzliu-open}{{23.51}{91}{开放问题}{theorem.23.51}{}}
\newlabel{rem:yzliu-open@cref}{{[remark][51][23]23.51}{[1][91][]91}}
\newlabel{ch:front}{{24}{92}{其他前沿严格结果}{chapter.24}{}}
\newlabel{ch:front@cref}{{[chapter][24][]24}{[1][92][]92}}
\newlabel{sec:front-intro}{{24.1}{92}{第 1 节（测试）}{section.24.1}{}}
\newlabel{sec:front-intro@cref}{{[section][1][24]24.1}{[1][92][]92}}
\newlabel{def:front-trap}{{24.1}{92}{围道变形、钳夹、Landau 点与好方向}{theorem.24.1}{}}
\newlabel{def:front-trap@cref}{{[definition][1][24]24.1}{[1][92][]92}}
\newlabel{lem:front-min-norm}{{24.2}{92}{Gordan 二择一的构造性形式：最小范数方向}{theorem.24.2}{}}
\newlabel{lem:front-min-norm@cref}{{[lemma][2][24]24.2}{[1][92][]92}}
\newlabel{rem:front-gordan-collins}{{24.3}{92}{Collins 的几何定理即 Gordan 二择一}{theorem.24.3}{}}
\newlabel{rem:front-gordan-collins@cref}{{[remark][3][24]24.3}{[1][92][]92}}
\newlabel{ex:front-gordan-2d}{{24.4}{92}{二维中的两种情形}{theorem.24.4}{}}
\newlabel{ex:front-gordan-2d@cref}{{[example][4][24]24.4}{[1][92][]92}}
\newlabel{prop:front-anomalous}{{24.5}{92}{"规避奇点需要严格正的一阶位移"不成立 \tagV }{theorem.24.5}{}}
\newlabel{prop:front-anomalous@cref}{{[proposition][5][24]24.5}{[1][92][]92}}
\newlabel{sec:front-landau}{{24.2}{92}{第 2 节（测试）}{section.24.2}{}}
\newlabel{sec:front-landau@cref}{{[section][2][24]24.2}{[1][92][]92}}
\newlabel{thm:front-collins-pinch}{{24.6}{92}{物理区钳夹判据（Collins 2020） \tagA }{theorem.24.6}{}}
\newlabel{thm:front-collins-pinch@cref}{{[theorem][6][24]24.6}{[1][92][]92}}
\newlabel{prop:front-landau-not-singular}{{24.7}{92}{Landau 条件只是奇点的必要条件 \tagV }{theorem.24.7}{}}
\newlabel{prop:front-landau-not-singular@cref}{{[proposition][7][24]24.7}{[1][92][]92}}
\newlabel{cor:front-coleman-norton}{{24.8}{93}{动量空间中的 Coleman–Norton 图像（含无质量线与 Wilson 线） \tagR }{theorem.24.8}{}}
\newlabel{cor:front-coleman-norton@cref}{{[corollary][8][24]24.8}{[1][92][]93}}
\newlabel{ex:front-collinear-pinch}{{24.9}{93}{无质量自能的共线钳夹}{theorem.24.9}{}}
\newlabel{ex:front-collinear-pinch@cref}{{[example][9][24]24.9}{[1][92][]93}}
\newlabel{rem:front-cn-history}{{24.10}{93}{Coleman–Norton 与 ELOP 原证明的缺口}{theorem.24.10}{}}
\newlabel{rem:front-cn-history@cref}{{[remark][10][24]24.10}{[1][93][]93}}
\newlabel{sec:front-ope}{{24.3}{93}{第 3 节（测试）}{section.24.3}{}}
\newlabel{sec:front-ope@cref}{{[section][3][24]24.3}{[1][93][]93}}
\newlabel{prop:front-ir-forest}{{24.11}{93}{宽角振幅的红外森林公式（Ma 2020） \tagA }{theorem.24.11}{}}
\newlabel{prop:front-ir-forest@cref}{{[proposition][11][24]24.11}{[1][93][]93}}
\newlabel{prop:front-coordinate-space}{{24.12}{93}{振幅的全阶硬—软—共线因子化：坐标空间路线与算符路线 \tagA }{theorem.24.12}{}}
\newlabel{prop:front-coordinate-space@cref}{{[proposition][12][24]24.12}{[1][93][]93}}
\newlabel{prop:front-local-ir}{{24.13}{93}{被积函数层次的局域红外因子化 \tagP }{theorem.24.13}{}}
\newlabel{prop:front-local-ir@cref}{{[proposition][13][24]24.13}{[1][93][]93}}
\newlabel{rem:front-landau-algebraic}{{24.14}{93}{代数几何的 Landau 分析与偶极子公式}{theorem.24.14}{}}
\newlabel{rem:front-landau-algebraic@cref}{{[remark][14][24]24.14}{[1][93][]93}}
\newlabel{prop:front-ope-scalar}{{24.15}{93}{标量场论中 OPE 的收敛、结合律与弯曲时空 \tagP }{theorem.24.15}{}}
\newlabel{prop:front-ope-scalar@cref}{{[proposition][15][24]24.15}{[1][93][]93}}
\newlabel{sec:front-regions}{{24.4}{94}{第 4 节（测试）}{section.24.4}{}}
\newlabel{sec:front-regions@cref}{{[section][4][24]24.4}{[1][93][]94}}
\newlabel{prop:front-ope-ym}{{24.16}{94}{纯 Yang–Mills 理论中 OPE 在规范不变算符上闭合 \tagP }{theorem.24.16}{}}
\newlabel{prop:front-ope-ym@cref}{{[proposition][16][24]24.16}{[1][93][]94}}
\newlabel{prop:front-flow-renorm}{{24.17}{94}{流方程（Wilson–Polchinski）重整化证明 \tagP }{theorem.24.17}{}}
\newlabel{prop:front-flow-renorm@cref}{{[proposition][17][24]24.17}{[1][94][]94}}
\newlabel{conj:front-ope-convergence}{{24.18}{94}{Yang–Mills 与 QCD 中 OPE 的收敛 \tagC }{theorem.24.18}{}}
\newlabel{conj:front-ope-convergence@cref}{{[conjecture][18][24]24.18}{[1][94][]94}}
\newlabel{rem:front-ope-qcd-status}{{24.19}{94}{对含夸克的 QCD，本书能断言什么}{theorem.24.19}{}}
\newlabel{rem:front-ope-qcd-status@cref}{{[remark][19][24]24.19}{[1][94][]94}}
\newlabel{lem:front-jantzen}{{24.20}{94}{区域展开的主恒等式（Jantzen）}{theorem.24.20}{}}
\newlabel{lem:front-jantzen@cref}{{[lemma][20][24]24.20}{[1][94][]94}}
\newlabel{sec:front-glauber}{{24.5}{94}{第 5 节（测试）}{section.24.5}{}}
\newlabel{sec:front-glauber@cref}{{[section][5][24]24.5}{[1][94][]94}}
\newlabel{def:front-facets}{{24.21}{94}{Lee–Pomeransky 表示、Newton 多面体与本质面}{theorem.24.21}{}}
\newlabel{def:front-facets@cref}{{[definition][21][24]24.21}{[1][94][]94}}
\newlabel{prop:front-lp-facets}{{24.22}{94}{系数全正时的面展开（Semenova–Smirnov–Smirnov） \tagR }{theorem.24.22}{}}
\newlabel{prop:front-lp-facets@cref}{{[proposition][22][24]24.22}{[1][94][]94}}
\newlabel{prop:front-onshell-facets}{{24.23}{94}{宽角 Green 函数在壳展开：区域 = Landau 解 = 面 \tagR }{theorem.24.23}{}}
\newlabel{prop:front-onshell-facets@cref}{{[proposition][23][24]24.23}{[1][94][]94}}
\newlabel{asm:front-mor-complete}{{24.24}{94}{区域集合的完备性 \tagC }{theorem.24.24}{}}
\newlabel{asm:front-mor-complete@cref}{{[assumption][24][24]24.24}{[1][94][]94}}
\newlabel{prop:front-commuting}{{24.25}{94}{多重在壳展开的可交换性 \tagR }{theorem.24.25}{}}
\newlabel{prop:front-commuting@cref}{{[proposition][25][24]24.25}{[1][94][]94}}
\newlabel{sec:front-sll}{{24.6}{95}{第 6 节（测试）}{section.24.6}{}}
\newlabel{sec:front-sll@cref}{{[section][6][24]24.6}{[1][94][]95}}
\newlabel{prop:front-hidden-regions}{{24.26}{95}{"区域 = Newton 多面体的下表面"在 Minkowski 运动学中不成立 \tagV }{theorem.24.26}{}}
\newlabel{prop:front-hidden-regions@cref}{{[proposition][26][24]24.26}{[1][94][]95}}
\newlabel{conj:front-mor-minkowski}{{24.27}{95}{Minkowski 极限下区域方法的有效性 \tagC }{theorem.24.27}{}}
\newlabel{conj:front-mor-minkowski@cref}{{[conjecture][27][24]24.27}{[1][95][]95}}
\newlabel{prop:front-glauber-bfkl}{{24.28}{95}{Glauber 有效理论：Regge 化、BFKL 与零箱减除 \tagA }{theorem.24.28}{}}
\newlabel{prop:front-glauber-bfkl@cref}{{[proposition][28][24]24.28}{[1][95][]95}}
\newlabel{prop:front-syz}{{24.29}{95}{纯 Glauber 梯形图在单标度红外安全观测量中的相消 \tagA }{theorem.24.29}{}}
\newlabel{prop:front-syz@cref}{{[proposition][29][24]24.29}{[1][95][]95}}
\newlabel{conj:front-mpi-violation}{{24.30}{95}{对多重部分子相互作用敏感的观测量：标准因子化的破坏 \tagC }{theorem.24.30}{}}
\newlabel{conj:front-mpi-violation@cref}{{[conjecture][30][24]24.30}{[1][95][]95}}
\newlabel{sec:front-nlp}{{24.7}{95}{第 7 节（测试）}{section.24.7}{}}
\newlabel{sec:front-nlp@cref}{{[section][7][24]24.7}{[1][95][]95}}
\newlabel{prop:front-heavy-fsi}{{24.31}{95}{重末态与旁观者之间的软末态相互作用 \tagA }{theorem.24.31}{}}
\newlabel{prop:front-heavy-fsi@cref}{{[proposition][31][24]24.31}{[1][95][]95}}
\newlabel{rem:front-glauber-status}{{24.32}{96}{Glauber 现状与本书分级的后果}{theorem.24.32}{}}
\newlabel{rem:front-glauber-status@cref}{{[remark][32][24]24.32}{[1][95][]96}}
\newlabel{prop:front-sll-all-orders}{{24.33}{96}{非全局 LHC 观测量的因子化与超领头对数的全阶重求和 \tagA }{theorem.24.33}{}}
\newlabel{prop:front-sll-all-orders@cref}{{[proposition][33][24]24.33}{[1][95][]96}}
\newlabel{prop:front-sll-closed-form}{{24.34}{96}{超领头对数的塔与 $q\bar q$ 单态交换的闭式 \tagA }{theorem.24.34}{}}
\newlabel{prop:front-sll-closed-form@cref}{{[proposition][34][24]24.34}{[1][96][]96}}
\newlabel{ex:front-sll-first}{{24.35}{96}{由闭式读出第一个超领头对数}{theorem.24.35}{}}
\newlabel{ex:front-sll-first@cref}{{[example][35][24]24.35}{[1][96][]96}}
\newlabel{sec:front-nlp-tmd}{{24.8}{96}{第 8 节（测试）}{section.24.8}{}}
\newlabel{sec:front-nlp-tmd@cref}{{[section][8][24]24.8}{[1][96][]96}}
\newlabel{prop:front-restoration}{{24.36}{96}{Glauber 胶子引起的"因子化恢复" \tagP }{theorem.24.36}{}}
\newlabel{prop:front-restoration@cref}{{[proposition][36][24]24.36}{[1][96][]96}}
\newlabel{prop:front-jettiness-cvl}{{24.37}{96}{1-jettiness 中的相干性破坏对数 \tagV }{theorem.24.37}{}}
\newlabel{prop:front-jettiness-cvl@cref}{{[proposition][37][24]24.37}{[1][96][]96}}
\newlabel{asm:front-scet-modes}{{24.38}{96}{SCET 的模式完备性与 Glauber 模式的缺席 \tagA }{theorem.24.38}{}}
\newlabel{asm:front-scet-modes@cref}{{[assumption][38][24]24.38}{[1][96][]96}}
\newlabel{def:front-endpoint}{{24.39}{96}{端点发散与重因子化条件}{theorem.24.39}{}}
\newlabel{def:front-endpoint@cref}{{[definition][39][24]24.39}{[1][96][]96}}
\newlabel{thm:front-hgg}{{24.40}{97}{$h\to \gamma \gamma $（$b$ 夸克圈）的重整化端点因子化定理 \tagA }{theorem.24.40}{}}
\newlabel{thm:front-hgg@cref}{{[theorem][40][24]24.40}{[1][96][]97}}
\newlabel{sec:front-positivity}{{24.9}{97}{第 9 节（测试）}{section.24.9}{}}
\newlabel{sec:front-positivity@cref}{{[section][9][24]24.9}{[1][97][]97}}
\newlabel{ex:front-hgg-ldl}{{24.41}{97}{领头双对数与 ${}_2F_2$}{theorem.24.41}{}}
\newlabel{ex:front-hgg-ldl@cref}{{[example][41][24]24.41}{[1][97][]97}}
\newlabel{prop:front-ggh}{{24.42}{97}{轻夸克诱导的 $gg\to h$ 与重因子化减除方案 \tagA }{theorem.24.42}{}}
\newlabel{prop:front-ggh@cref}{{[proposition][42][24]24.42}{[1][97][]97}}
\newlabel{prop:front-nlp-threshold}{{24.43}{97}{次领头幂次阈值 Drell–Yan 的 $d$ 维因子化 \tagP }{theorem.24.43}{}}
\newlabel{prop:front-nlp-threshold@cref}{{[proposition][43][24]24.43}{[1][97][]97}}
\newlabel{prop:front-nlp-ll}{{24.44}{97}{次领头幂次的领头对数重求和与端点因子化条件 \tagA }{theorem.24.44}{}}
\newlabel{prop:front-nlp-ll@cref}{{[proposition][44][24]24.44}{[1][97][]97}}
\newlabel{prop:front-nlp-dis}{{24.45}{97}{次领头幂次 DIS 与 PDF 因子化：阈值 PDF 的双标度结构 \tagA }{theorem.24.45}{}}
\newlabel{prop:front-nlp-dis@cref}{{[proposition][45][24]24.45}{[1][97][]97}}
\newlabel{sec:front-brst}{{24.10}{98}{第 10 节（测试）}{section.24.10}{}}
\newlabel{sec:front-brst@cref}{{[section][10][24]24.10}{[1][97][]98}}
\newlabel{conj:front-nlp-beyond-ll}{{24.46}{98}{超出领头对数的次领头幂次端点因子化 \tagC }{theorem.24.46}{}}
\newlabel{conj:front-nlp-beyond-ll@cref}{{[conjecture][46][24]24.46}{[1][97][]98}}
\newlabel{conj:front-glauber-nlp}{{24.47}{98}{次领头幂次 TMD 过程中 Glauber 交换的相消 \tagC }{theorem.24.47}{}}
\newlabel{conj:front-glauber-nlp@cref}{{[conjecture][47][24]24.47}{[1][98][]98}}
\newlabel{prop:front-tmd-nlp-operator}{{24.48}{98}{TMD 算符展开与运动学幂次修正 \tagA }{theorem.24.48}{}}
\newlabel{prop:front-tmd-nlp-operator@cref}{{[proposition][48][24]24.48}{[1][98][]98}}
\newlabel{prop:front-sidis-nlp}{{24.49}{98}{SIDIS 方位角不对称的次领头幂次 SCET 因子化 \tagA }{theorem.24.49}{}}
\newlabel{prop:front-sidis-nlp@cref}{{[proposition][49][24]24.49}{[1][98][]98}}
\newlabel{prop:front-nlp-tmd-oneloop}{{24.50}{98}{单圈次领头幂次系数函数与扭度 3 TMD 演化 \tagP }{theorem.24.50}{}}
\newlabel{prop:front-nlp-tmd-oneloop@cref}{{[proposition][50][24]24.50}{[1][98][]98}}
\newlabel{sec:front-smallx}{{24.11}{98}{第 11 节（测试）}{section.24.11}{}}
\newlabel{sec:front-smallx@cref}{{[section][11][24]24.11}{[1][98][]98}}
\newlabel{rem:front-power-rapidity}{{24.51}{98}{次领头幂次的幂律快度发散}{theorem.24.51}{}}
\newlabel{rem:front-power-rapidity@cref}{{[remark][51][24]24.51}{[1][98][]98}}
\newlabel{prop:front-msbar-negative}{{24.52}{98}{$\overline {\rm MS}$ 部分子分布可以为负 \tagV }{theorem.24.52}{}}
\newlabel{prop:front-msbar-negative@cref}{{[proposition][52][24]24.52}{[1][98][]98}}
\newlabel{ex:front-yukawa}{{24.53}{99}{Yukawa 模型分布的端点行为与符号变化}{theorem.24.53}{}}
\newlabel{ex:front-yukawa@cref}{{[example][53][24]24.53}{[1][98][]99}}
\newlabel{prop:front-positivity-large-scale}{{24.54}{99}{足够大标度上的正定性 \tagP }{theorem.24.54}{}}
\newlabel{prop:front-positivity-large-scale@cref}{{[proposition][54][24]24.54}{[1][99][]99}}
\newlabel{rem:front-positivity-tracks}{{24.55}{99}{对本书与拟合的后果}{theorem.24.55}{}}
\newlabel{rem:front-positivity-tracks@cref}{{[remark][55][24]24.55}{[1][99][]99}}
\newlabel{sec:front-dps}{{24.12}{99}{第 12 节（测试）}{section.24.12}{}}
\newlabel{sec:front-dps@cref}{{[section][12][24]24.12}{[1][99][]99}}
\newlabel{rem:front-k-gluon}{{24.56}{99}{为什么需要 BRST：K 胶子与非定域缺口}{theorem.24.56}{}}
\newlabel{rem:front-k-gluon@cref}{{[remark][56][24]24.56}{[1][99][]99}}
\newlabel{prop:front-brst-bilocal}{{24.57}{99}{双局域 PDF 算符 BRST 分类的矩约化 \tagA }{theorem.24.57}{}}
\newlabel{prop:front-brst-bilocal@cref}{{[proposition][57][24]24.57}{[1][99][]99}}
\newlabel{conj:front-brst-multilocal}{{24.58}{99}{多局域光线算符的 BRST 分类 \tagC }{theorem.24.58}{}}
\newlabel{conj:front-brst-multilocal@cref}{{[conjecture][58][24]24.58}{[1][99][]99}}
\newlabel{prop:front-nlo-bk}{{24.59}{99}{NLO BK/JIMWLK 与三圈 BK 核 \tagA }{theorem.24.59}{}}
\newlabel{prop:front-nlo-bk@cref}{{[proposition][59][24]24.59}{[1][99][]99}}
\newlabel{prop:front-bk-collinear}{{24.60}{99}{NLO BK 的不稳定性与共线重求和 \tagP }{theorem.24.60}{}}
\newlabel{prop:front-bk-collinear@cref}{{[proposition][60][24]24.60}{[1][99][]99}}
\newlabel{sec:front-jets}{{24.13}{100}{第 13 节（测试）}{section.24.13}{}}
\newlabel{sec:front-jets@cref}{{[section][13][24]24.13}{[1][99][]100}}
\newlabel{prop:front-regge}{{24.61}{100}{高能极限中的 Wilson 线方法、胶子 Regge 化与 Regge 极限的红外奇异部分 \tagA }{theorem.24.61}{}}
\newlabel{prop:front-regge@cref}{{[proposition][61][24]24.61}{[1][99][]100}}
\newlabel{prop:front-regge-pole}{{24.62}{100}{简单 Regge 极点因子化在 NNLL 失效 \tagV }{theorem.24.62}{}}
\newlabel{prop:front-regge-pole@cref}{{[proposition][62][24]24.62}{[1][100][]100}}
\newlabel{prop:front-smallx-tmd}{{24.63}{100}{小 $x$ 的 TMD 与 Sudakov 对数 \tagA }{theorem.24.63}{}}
\newlabel{prop:front-smallx-tmd@cref}{{[proposition][63][24]24.63}{[1][100][]100}}
\newlabel{prop:front-smallx-helicity}{{24.64}{100}{小 $x$ 螺旋度演化：截距与 BER 的分歧 \tagA }{theorem.24.64}{}}
\newlabel{prop:front-smallx-helicity@cref}{{[proposition][64][24]24.64}{[1][100][]100}}
\newlabel{prop:front-dps-scheme}{{24.65}{100}{无重复计数的双部分子散射 \tagA }{theorem.24.65}{}}
\newlabel{prop:front-dps-scheme@cref}{{[proposition][65][24]24.65}{[1][100][]100}}
\newlabel{sec:front-quarkonium}{{24.14}{100}{第 14 节（测试）}{section.24.14}{}}
\newlabel{sec:front-quarkonium@cref}{{[section][14][24]24.14}{[1][100][]100}}
\newlabel{prop:front-rapidity-renorm}{{24.66}{100}{多部分子软因子快度发散的重整化定理 \tagA }{theorem.24.66}{}}
\newlabel{prop:front-rapidity-renorm@cref}{{[proposition][66][24]24.66}{[1][100][]100}}
\newlabel{prop:front-dpd-evolution}{{24.67}{101}{色关联 DPD 的演化与匹配 \tagP }{theorem.24.67}{}}
\newlabel{prop:front-dpd-evolution@cref}{{[proposition][67][24]24.67}{[1][100][]101}}
\newlabel{prop:front-ngl}{{24.68}{101}{非全局对数：超出领头对数与测量驱动的因子化 \tagA }{theorem.24.68}{}}
\newlabel{prop:front-ngl@cref}{{[proposition][68][24]24.68}{[1][101][]101}}
\newlabel{prop:front-jet-substructure}{{24.69}{101}{喷注子结构、修饰喷注与半单举喷注函数 \tagA }{theorem.24.69}{}}
\newlabel{prop:front-jet-substructure@cref}{{[proposition][69][24]24.69}{[1][101][]101}}
\newlabel{prop:front-lightray-ope}{{24.70}{101}{能量关联子的共线因子化与共形场论中的光线 OPE \tagA }{theorem.24.70}{}}
\newlabel{prop:front-lightray-ope@cref}{{[proposition][70][24]24.70}{[1][101][]101}}
\newlabel{sec:front-qed}{{24.15}{101}{第 15 节（测试）}{section.24.15}{}}
\newlabel{sec:front-qed@cref}{{[section][15][24]24.15}{[1][101][]101}}
\newlabel{prop:front-nrqcd-ir}{{24.71}{101}{NNLO 的红外极点与规范完备的八重态算符 \tagP }{theorem.24.71}{}}
\newlabel{prop:front-nrqcd-ir@cref}{{[proposition][71][24]24.71}{[1][101][]101}}
\newlabel{rem:front-quarkonium-alternatives}{{24.72}{101}{唯象张力与替代方案}{theorem.24.72}{}}
\newlabel{rem:front-quarkonium-alternatives@cref}{{[remark][72][24]24.72}{[1][101][]101}}
\newlabel{thm:front-luxqed}{{24.73}{101}{由结构函数确定质子的光子 PDF（LUXqed） \tagA }{theorem.24.73}{}}
\newlabel{thm:front-luxqed@cref}{{[theorem][73][24]24.73}{[1][101][]101}}
\newlabel{ex:front-luxqed-pointlike}{{24.74}{102}{点状靶极限}{theorem.24.74}{}}
\newlabel{ex:front-luxqed-pointlike@cref}{{[example][74][24]24.74}{[1][101][]102}}
\newlabel{prop:front-lepton-pdf}{{24.75}{102}{轻子与电子的 PDF、QED 与混合演化核 \tagA }{theorem.24.75}{}}
\newlabel{prop:front-lepton-pdf@cref}{{[proposition][75][24]24.75}{[1][102][]102}}
\newlabel{sec:front-status}{{24.16}{102}{第 16 节（测试）}{section.24.16}{}}
\newlabel{sec:front-status@cref}{{[section][16][24]24.16}{[1][102][]102}}
\newlabel{prop:front-ew-pdf}{{24.76}{102}{电弱 PDF 与 $Q\gg M_W$ 的电弱对数 \tagA }{theorem.24.76}{}}
\newlabel{prop:front-ew-pdf@cref}{{[proposition][76][24]24.76}{[1][102][]102}}
\newlabel{prop:front-ew-violation}{{24.77}{102}{电弱 Bloch–Nordsieck 抵消的失效与测量弱荷的观测量 \tagV }{theorem.24.77}{}}
\newlabel{prop:front-ew-violation@cref}{{[proposition][77][24]24.77}{[1][102][]102}}
\newlabel{rem:front-open}{{24.78}{102}{开放问题}{theorem.24.78}{}}
\newlabel{rem:front-open@cref}{{[remark][78][24]24.78}{[1][102][]102}}
\newlabel{ch:summary}{{25}{104}{总结：主表、依赖图与开放问题}{chapter.25}{}}
\newlabel{ch:summary@cref}{{[chapter][25][]25}{[1][104][]104}}
\newlabel{tab:summary-master}{{25.1}{105}{全书主要结论总表\relax }{table.caption.239}{}}
\newlabel{tab:summary-master@cref}{{[table][1][25]25.1}{[1][104][]105}}
\newlabel{tab:summary-master-b}{{25.2}{106}{全书主要结论总表（续）\relax }{table.caption.240}{}}
\newlabel{tab:summary-master-b@cref}{{[table][2][25]25.2}{[1][104][]106}}
\newlabel{tab:summary-master-c}{{25.3}{107}{全书主要结论总表（续）\relax }{table.caption.241}{}}
\newlabel{tab:summary-master-c@cref}{{[table][3][25]25.3}{[1][104][]107}}
\newlabel{tab:summary-master-d}{{25.4}{108}{全书主要结论总表（续）\relax }{table.caption.242}{}}
\newlabel{tab:summary-master-d@cref}{{[table][4][25]25.4}{[1][104][]108}}
\newlabel{tab:summary-master-e}{{25.5}{109}{全书主要结论总表（续）\relax }{table.caption.243}{}}
\newlabel{tab:summary-master-e@cref}{{[table][5][25]25.5}{[1][104][]109}}
\newlabel{tab:summary-master-f}{{25.6}{109}{全书主要结论总表（续）\relax }{table.caption.244}{}}
\newlabel{tab:summary-master-f@cref}{{[table][6][25]25.6}{[1][104][]109}}
\newlabel{app:appA}{{A}{111}{光锥坐标、Dirac 代数、加号分布与 Mellin 变换}{appendix.A}{}}
\newlabel{app:appA@cref}{{[appendix][1][2147483647]A}{[1][111][]111}}
\newlabel{sec:appA-lightfront}{{A.1}{111}{第 1 节（测试）}{section.A.1}{}}
\newlabel{sec:appA-lightfront@cref}{{[subappendix][1][2147483647,1]A.1}{[1][111][]111}}
\newlabel{lem:appA-fierz}{{A.1}{111}{共线 Fierz 分解}{theorem.A.1}{}}
\newlabel{lem:appA-fierz@cref}{{[lemma][1][2147483647,1]A.1}{[1][111][]111}}
\newlabel{def:appA-test-functions}{{A.2}{111}{测试函数类与弱收敛}{theorem.A.2}{}}
\newlabel{def:appA-test-functions@cref}{{[definition][2][2147483647,1]A.2}{[1][111][]111}}
\newlabel{sec:appA-dirac}{{A.2}{111}{第 2 节（测试）}{section.A.2}{}}
\newlabel{sec:appA-dirac@cref}{{[subappendix][2][2147483647,1]A.2}{[1][111][]111}}
\newlabel{lem:appA-plus-product}{{A.3}{111}{乘积约定下的恒等式}{theorem.A.3}{}}
\newlabel{lem:appA-plus-product@cref}{{[lemma][3][2147483647,1]A.3}{[1][111][]111}}
\newlabel{lem:appA-weak-convergence}{{A.4}{111}{由偶数矩得到弱收敛}{theorem.A.4}{}}
\newlabel{lem:appA-weak-convergence@cref}{{[lemma][4][2147483647,1]A.4}{[1][111][]111}}
\newlabel{sec:appA-plus}{{A.3}{111}{第 3 节（测试）}{section.A.3}{}}
\newlabel{sec:appA-plus@cref}{{[subappendix][3][2147483647,1]A.3}{[1][111][]111}}
\newlabel{lem:appA-mellin-table}{{A.5}{111}{Mellin 变换表}{theorem.A.5}{}}
\newlabel{lem:appA-mellin-table@cref}{{[lemma][5][2147483647,1]A.5}{[1][111][]111}}
\newlabel{app:appB}{{B}{112}{Wilson 线：Feynman 规则、程函恒等式与尖点}{appendix.B}{}}
\newlabel{app:appB@cref}{{[appendix][2][2147483647]B}{[1][112][]112}}
\newlabel{sec:appB-definition}{{B.1}{112}{第 1 节（测试）}{section.B.1}{}}
\newlabel{sec:appB-definition@cref}{{[subappendix][1][2147483647,2]B.1}{[1][112][]112}}
\newlabel{prop:appB-rules}{{B.1}{112}{Wilson 线的 Feynman 规则 \tagR }{theorem.B.1}{}}
\newlabel{prop:appB-rules@cref}{{[proposition][1][2147483647,2]B.1}{[1][112][]112}}
\newlabel{sec:appB-rules}{{B.2}{112}{第 2 节（测试）}{section.B.2}{}}
\newlabel{sec:appB-rules@cref}{{[subappendix][2][2147483647,2]B.2}{[1][112][]112}}
\newlabel{lem:appB-shuffle}{{B.2}{112}{程函恒等式（洗牌积）}{theorem.B.2}{}}
\newlabel{lem:appB-shuffle@cref}{{[lemma][2][2147483647,2]B.2}{[1][112][]112}}
\newlabel{sec:appB-eikonal}{{B.3}{112}{第 3 节（测试）}{section.B.3}{}}
\newlabel{sec:appB-eikonal@cref}{{[subappendix][3][2147483647,2]B.3}{[1][112][]112}}
\newlabel{prop:appB-webs}{{B.3}{112}{非阿贝尔指数化 \tagA }{theorem.B.3}{}}
\newlabel{prop:appB-webs@cref}{{[proposition][3][2147483647,2]B.3}{[1][112][]112}}
\newlabel{sec:appB-cusp}{{B.4}{112}{第 4 节（测试）}{section.B.4}{}}
\newlabel{sec:appB-cusp@cref}{{[subappendix][4][2147483647,2]B.4}{[1][112][]112}}
\newlabel{prop:appB-cusp-oneloop}{{B.4}{112}{单圈尖点反常维度 \tagR }{theorem.B.4}{}}
\newlabel{prop:appB-cusp-oneloop@cref}{{[proposition][4][2147483647,2]B.4}{[1][112][]112}}
\newlabel{app:appC}{{C}{113}{维数正规化积分与相空间}{appendix.C}{}}
\newlabel{app:appC@cref}{{[appendix][3][2147483647]C}{[1][113][]113}}
\newlabel{sec:appC-msbar}{{C.1}{113}{第 1 节（测试）}{section.C.1}{}}
\newlabel{sec:appC-msbar@cref}{{[subappendix][1][2147483647,3]C.1}{[1][113][]113}}
\newlabel{rem:appC-Seps}{{C.1}{113}{$S_\varepsilon $ 约定的影响}{theorem.C.1}{}}
\newlabel{rem:appC-Seps@cref}{{[remark][1][2147483647,3]C.1}{[1][113][]113}}
\newlabel{rem:appC-d-kernels}{{C.2}{113}{$D$ 维劈裂核与平均约定}{theorem.C.2}{}}
\newlabel{rem:appC-d-kernels@cref}{{[remark][2][2147483647,3]C.2}{[1][113][]113}}
\newlabel{sec:appC-integrals}{{C.2}{113}{第 2 节（测试）}{section.C.2}{}}
\newlabel{sec:appC-integrals@cref}{{[subappendix][2][2147483647,3]C.2}{[1][113][]113}}
\newlabel{lem:appC-bubble}{{C.3}{113}{单圈两点函数}{theorem.C.3}{}}
\newlabel{lem:appC-bubble@cref}{{[lemma][3][2147483647,3]C.3}{[1][113][]113}}
\newlabel{lem:appC-ibp}{{C.4}{113}{分部积分（三角规则）}{theorem.C.4}{}}
\newlabel{lem:appC-ibp@cref}{{[lemma][4][2147483647,3]C.4}{[1][113][]113}}
\newlabel{sec:appC-phase-space}{{C.3}{113}{第 3 节（测试）}{section.C.3}{}}
\newlabel{sec:appC-phase-space@cref}{{[subappendix][3][2147483647,3]C.3}{[1][113][]113}}
\newlabel{lem:appC-phase-space}{{C.5}{113}{$D$ 维相空间}{theorem.C.5}{}}
\newlabel{lem:appC-phase-space@cref}{{[lemma][5][2147483647,3]C.5}{[1][113][]113}}
\newlabel{app:appD}{{D}{114}{Landau 方程与 Coleman–Norton 定理}{appendix.D}{}}
\newlabel{app:appD@cref}{{[appendix][4][2147483647]D}{[1][114][]114}}
\newlabel{sec:appD-gordan}{{D.1}{114}{第 1 节（测试）}{section.D.1}{}}
\newlabel{sec:appD-gordan@cref}{{[subappendix][1][2147483647,4]D.1}{[1][114][]114}}
\newlabel{lem:appD-gordan}{{D.1}{114}{Gordan 择一定理}{theorem.D.1}{}}
\newlabel{lem:appD-gordan@cref}{{[lemma][1][2147483647,4]D.1}{[1][114][]114}}
\newlabel{sec:appD-momentum}{{D.2}{114}{第 2 节（测试）}{section.D.2}{}}
\newlabel{sec:appD-momentum@cref}{{[subappendix][2][2147483647,4]D.2}{[1][114][]114}}
\newlabel{rem:appD-parameter-proof}{{D.2}{114}{Landau 条件的 Feynman 参数证明及其注意事项}{theorem.D.2}{}}
\newlabel{rem:appD-parameter-proof@cref}{{[remark][2][2147483647,4]D.2}{[1][114][]114}}
\newlabel{app:appE}{{E}{115}{幂次计数公式的推导}{appendix.E}{}}
\newlabel{app:appE@cref}{{[appendix][5][2147483647]E}{[1][115][]115}}
\newlabel{sec:appE-units}{{E.1}{115}{第 1 节（测试）}{section.E.1}{}}
\newlabel{sec:appE-units@cref}{{[subappendix][1][2147483647,5]E.1}{[1][115][]115}}
\newlabel{prop:appE-degree}{{E.1}{115}{约化图的超表观红外发散度 \tagR }{theorem.E.1}{}}
\newlabel{prop:appE-degree@cref}{{[proposition][1][2147483647,5]E.1}{[1][115][]115}}
\newlabel{sec:appE-formula}{{E.2}{115}{第 2 节（测试）}{section.E.2}{}}
\newlabel{sec:appE-formula@cref}{{[subappendix][2][2147483647,5]E.2}{[1][115][]115}}
\newlabel{rem:appE-actual-bound}{{E.2}{115}{从超表观度到实际界 \tagA }{theorem.E.2}{}}
\newlabel{rem:appE-actual-bound@cref}{{[remark][2][2147483647,5]E.2}{[1][115][]115}}
\newlabel{app:appF}{{F}{116}{图形 Ward 恒等式}{appendix.F}{}}
\newlabel{app:appF@cref}{{[appendix][6][2147483647]F}{[1][116][]116}}
\newlabel{sec:appF-one}{{F.1}{116}{第 1 节（测试）}{section.F.1}{}}
\newlabel{sec:appF-one@cref}{{[subappendix][1][2147483647,6]F.1}{[1][116][]116}}
\newlabel{ex:appF-telescoping}{{F.1}{116}{有限 Wilson 线的伸缩恒等式}{theorem.F.1}{}}
\newlabel{ex:appF-telescoping@cref}{{[example][1][2147483647,6]F.1}{[1][116][]116}}
\newlabel{sec:appF-two}{{F.2}{116}{第 2 节（测试）}{section.F.2}{}}
\newlabel{sec:appF-two@cref}{{[subappendix][2][2147483647,6]F.2}{[1][116][]116}}
\newlabel{prop:appF-two-gluon}{{F.2}{116}{两个纵向胶子：树图层次的完整证明 \tagR }{theorem.F.2}{}}
\newlabel{prop:appF-two-gluon@cref}{{[proposition][2][2147483647,6]F.2}{[1][116][]116}}
\newlabel{sec:appF-induction}{{F.3}{116}{第 3 节（测试）}{section.F.3}{}}
\newlabel{sec:appF-induction@cref}{{[subappendix][3][2147483647,6]F.3}{[1][116][]116}}
\newlabel{prop:appF-induction}{{F.3}{116}{多胶子情形的归纳论证 \tagA }{theorem.F.3}{}}
\newlabel{prop:appF-induction@cref}{{[proposition][3][2147483647,6]F.3}{[1][116][]116}}
\newlabel{app:appG}{{G}{117}{单圈与两圈计算细节}{appendix.G}{}}
\newlabel{app:appG@cref}{{[appendix][7][2147483647]G}{[1][117][]117}}
\newlabel{app:appH}{{H}{118}{文献导读}{appendix.H}{}}
\newlabel{app:appH@cref}{{[appendix][8][2147483647]H}{[1][118][]118}}
\newlabel{app:appI}{{I}{119}{维数正规化中的 $\gamma _5$}{appendix.I}{}}
\newlabel{app:appI@cref}{{[appendix][9][2147483647]I}{[1][119][]119}}
\newlabel{sec:appI-hvbm}{{I.1}{119}{第 1 节（测试）}{section.I.1}{}}
\newlabel{sec:appI-hvbm@cref}{{[subappendix][1][2147483647,9]I.1}{[1][119][]119}}
\newlabel{def:appI-hvbm}{{I.1}{119}{HVBM 方案}{theorem.I.1}{}}
\newlabel{def:appI-hvbm@cref}{{[definition][1][2147483647,9]I.1}{[1][119][]119}}
\newlabel{sec:appI-larin}{{I.2}{119}{第 2 节（测试）}{section.I.2}{}}
\newlabel{sec:appI-larin@cref}{{[subappendix][2][2147483647,9]I.2}{[1][119][]119}}
\newlabel{prop:appI-larin}{{I.2}{119}{Larin 方案与非单态轴矢流的有限重整化 \tagR }{theorem.I.2}{}}
\newlabel{prop:appI-larin@cref}{{[proposition][2][2147483647,9]I.2}{[1][119][]119}}
\newlabel{sec:appI-anomaly}{{I.3}{119}{第 3 节（测试）}{section.I.3}{}}
\newlabel{sec:appI-anomaly@cref}{{[subappendix][3][2147483647,9]I.3}{[1][119][]119}}
\newlabel{rem:appI-anomaly}{{I.3}{119}{轴矢反常}{theorem.I.3}{}}
\newlabel{rem:appI-anomaly@cref}{{[remark][3][2147483647,9]I.3}{[1][119][]119}}
\newlabel{app:appJ}{{J}{120}{BRST 对称性与割图中的鬼场}{appendix.J}{}}
\newlabel{app:appJ@cref}{{[appendix][10][2147483647]J}{[1][120][]120}}
\newlabel{sec:appJ-brst}{{J.1}{120}{第 1 节（测试）}{section.J.1}{}}
\newlabel{sec:appJ-brst@cref}{{[subappendix][1][2147483647,10]J.1}{[1][120][]120}}
\newlabel{prop:appJ-ghost-cuts}{{J.1}{120}{割图中的鬼场贡献 \tagR }{theorem.J.1}{}}
\newlabel{prop:appJ-ghost-cuts@cref}{{[proposition][1][2147483647,10]J.1}{[1][120][]120}}
\newlabel{lem:appJ-cut-hermiticity}{{J.2}{120}{割图规则的厄米性 \tagR }{theorem.J.2}{}}
\newlabel{lem:appJ-cut-hermiticity@cref}{{[lemma][2][2147483647,10]J.2}{[1][120][]120}}
\newlabel{sec:appJ-cuts}{{J.2}{120}{第 2 节（测试）}{section.J.2}{}}
\newlabel{sec:appJ-cuts@cref}{{[subappendix][2][2147483647,10]J.2}{[1][120][]120}}
\newlabel{prop:appJ-local-operators}{{J.3}{120}{规范不变定域算符的重整化与混合 \tagR }{theorem.J.3}{}}
\newlabel{prop:appJ-local-operators@cref}{{[proposition][3][2147483647,10]J.3}{[1][120][]120}}
\newlabel{rem:appJ-nonlocal}{{J.4}{120}{非定域算符}{theorem.J.4}{}}
\newlabel{rem:appJ-nonlocal@cref}{{[remark][4][2147483647,10]J.4}{[1][120][]120}}
\newlabel{app:appK}{{K}{121}{矩问题与解析延拓引理}{appendix.K}{}}
\newlabel{app:appK@cref}{{[appendix][11][2147483647]K}{[1][121][]121}}
\newlabel{sec:appK-hausdorff}{{K.1}{121}{第 1 节（测试）}{section.K.1}{}}
\newlabel{sec:appK-hausdorff@cref}{{[subappendix][1][2147483647,11]K.1}{[1][121][]121}}
\newlabel{lem:appK-hausdorff}{{K.1}{121}{Hausdorff 唯一性}{theorem.K.1}{}}
\newlabel{lem:appK-hausdorff@cref}{{[lemma][1][2147483647,11]K.1}{[1][121][]121}}
\newlabel{thm:appK-even-moments}{{K.2}{121}{偶数矩决定函数}{theorem.K.2}{}}
\newlabel{thm:appK-even-moments@cref}{{[theorem][2][2147483647,11]K.2}{[1][121][]121}}
\newlabel{cor:appK-pdf}{{K.3}{121}{PDF 的唯一性}{theorem.K.3}{}}
\newlabel{cor:appK-pdf@cref}{{[corollary][3][2147483647,11]K.3}{[1][121][]121}}
\newlabel{sec:appK-continuation}{{K.2}{121}{第 2 节（测试）}{section.K.2}{}}
\newlabel{sec:appK-continuation@cref}{{[subappendix][2][2147483647,11]K.2}{[1][121][]121}}
\newlabel{prop:appK-kernels}{{K.4}{121}{领头阶核由偶数矩唯一确定}{theorem.K.4}{}}
\newlabel{prop:appK-kernels@cref}{{[proposition][4][2147483647,11]K.4}{[1][121][]121}}
\newlabel{thm:appK-carlson}{{K.5}{121}{Carlson 定理}{theorem.K.5}{}}
\newlabel{thm:appK-carlson@cref}{{[theorem][5][2147483647,11]K.5}{[1][121][]121}}
\newlabel{prop:appK-continuation}{{K.6}{121}{多项式有界条件下延拓的唯一性}{theorem.K.6}{}}
\newlabel{prop:appK-continuation@cref}{{[proposition][6][2147483647,11]K.6}{[1][121][]121}}
\newlabel{sec:appK-entire}{{K.3}{121}{第 3 节（测试）}{section.K.3}{}}
\newlabel{sec:appK-entire@cref}{{[subappendix][3][2147483647,11]K.3}{[1][121][]121}}
\newlabel{thm:appK-muntz}{{K.7}{121}{Müntz–Szász 定理}{theorem.K.7}{}}
\newlabel{thm:appK-muntz@cref}{{[theorem][7][2147483647,11]K.7}{[1][121][]121}}
\newlabel{lem:appK-entire}{{K.8}{121}{受控增长的整函数}{theorem.K.8}{}}
\newlabel{lem:appK-entire@cref}{{[lemma][8][2147483647,11]K.8}{[1][121][]121}}
\newlabel{lem:appK-markov}{{K.9}{121}{V. A. Markov 系数不等式}{theorem.K.9}{}}
\newlabel{lem:appK-markov@cref}{{[lemma][9][2147483647,11]K.9}{[1][121][]121}}
\newlabel{sec:appK-counterexamples}{{K.4}{121}{第 4 节（测试）}{section.K.4}{}}
\newlabel{sec:appK-counterexamples@cref}{{[subappendix][4][2147483647,11]K.4}{[1][121][]121}}
\newlabel{rem:appK-counterexamples}{{K.10}{121}{反例与稳定性}{theorem.K.10}{}}
\newlabel{rem:appK-counterexamples@cref}{{[remark][10][2147483647,11]K.10}{[1][121][]121}}
\newlabel{app:appL}{{L}{122}{光线算符与定域算符：矩关系的全阶证明}{appendix.L}{}}
\newlabel{app:appL@cref}{{[appendix][12][2147483647]L}{[1][122][]122}}
\newlabel{sec:appL-vertices}{{L.1}{122}{第 1 节（测试）}{section.L.1}{}}
\newlabel{sec:appL-vertices@cref}{{[subappendix][1][2147483647,12]L.1}{[1][122][]122}}
\newlabel{lem:appL-divided-differences}{{L.1}{122}{均差恒等式}{theorem.L.1}{}}
\newlabel{lem:appL-divided-differences@cref}{{[lemma][1][2147483647,12]L.1}{[1][122][]122}}
\newlabel{sec:appL-bare}{{L.2}{122}{第 2 节（测试）}{section.L.2}{}}
\newlabel{sec:appL-bare@cref}{{[subappendix][2][2147483647,12]L.2}{[1][122][]122}}
\newlabel{prop:appL-bare}{{L.2}{122}{裸算符恒等式 \tagR }{theorem.L.2}{}}
\newlabel{prop:appL-bare@cref}{{[proposition][2][2147483647,12]L.2}{[1][122][]122}}
\newlabel{sec:appL-renormalized}{{L.3}{122}{第 3 节（测试）}{section.L.3}{}}
\newlabel{sec:appL-renormalized@cref}{{[subappendix][3][2147483647,12]L.3}{[1][122][]122}}
\newlabel{prop:appL-renormalized}{{L.3}{122}{重整化后的矩关系 \tagR }{theorem.L.3}{}}
\newlabel{prop:appL-renormalized@cref}{{[proposition][3][2147483647,12]L.3}{[1][122][]122}}
\newlabel{app:appM}{{M}{123}{OPE：微扰证明、无质量 QCD、非微扰条件定理}{appendix.M}{}}
\newlabel{app:appM@cref}{{[appendix][13][2147483647]M}{[1][123][]123}}
\newlabel{sec:appM-overview}{{M.1}{123}{第 1 节（测试）}{section.M.1}{}}
\newlabel{sec:appM-overview@cref}{{[subappendix][1][2147483647,13]M.1}{[1][123][]123}}
\newlabel{def:appM-bphz}{{M.1}{123}{图、表观发散度、森林与 $R$ 运算}{theorem.M.1}{}}
\newlabel{def:appM-bphz@cref}{{[definition][1][2147483647,13]M.1}{[1][123][]123}}
\newlabel{lem:appM-degree}{{M.2}{123}{表观发散度的外线公式}{theorem.M.2}{}}
\newlabel{lem:appM-degree@cref}{{[lemma][2][2147483647,13]M.2}{[1][123][]123}}
\newlabel{thm:appM-zimmermann}{{M.3}{123}{Zimmermann 收敛定理（顶点可加的超减除） \tagR }{theorem.M.3}{}}
\newlabel{thm:appM-zimmermann@cref}{{[theorem][3][2147483647,13]M.3}{[1][123][]123}}
\newlabel{ex:appM-nonadditive}{{M.4}{123}{非顶点可加的超减除会失败}{theorem.M.4}{}}
\newlabel{ex:appM-nonadditive@cref}{{[example][4][2147483647,13]M.4}{[1][123][]123}}
\newlabel{thm:appM-weinberg}{{M.5}{123}{Weinberg 定理（Hahn–Zimmermann 形式）}{theorem.M.5}{}}
\newlabel{thm:appM-weinberg@cref}{{[theorem][5][2147483647,13]M.5}{[1][123][]123}}
\newlabel{cor:appM-weinberg-degree}{{M.6}{123}{渐近指数与子图发散度}{theorem.M.6}{}}
\newlabel{cor:appM-weinberg-degree@cref}{{[corollary][6][2147483647,13]M.6}{[1][123][]123}}
\newlabel{rem:appM-forest-asymptotics}{{M.7}{123}{森林公式不能逐项用 Weinberg 定理}{theorem.M.7}{}}
\newlabel{rem:appM-forest-asymptotics@cref}{{[remark][7][2147483647,13]M.7}{[1][123][]123}}
\newlabel{sec:appM-bphz}{{M.2}{123}{第 2 节（测试）}{section.M.2}{}}
\newlabel{sec:appM-bphz@cref}{{[subappendix][2][2147483647,13]M.2}{[1][123][]123}}
\newlabel{def:appM-normal-product}{{M.8}{123}{正规乘积与冻结 $q$ 的双局域正规乘积}{theorem.M.8}{}}
\newlabel{def:appM-normal-product@cref}{{[definition][8][2147483647,13]M.8}{[1][123][]123}}
\newlabel{lem:appM-contraction}{{M.9}{123}{收缩与桥}{theorem.M.9}{}}
\newlabel{lem:appM-contraction@cref}{{[lemma][9][2147483647,13]M.9}{[1][123][]123}}
\newlabel{lem:appM-chain-expansion}{{M.10}{123}{链展开恒等式}{theorem.M.10}{}}
\newlabel{lem:appM-chain-expansion@cref}{{[lemma][10][2147483647,13]M.10}{[1][123][]123}}
\newlabel{thm:appM-zimmermann-identity}{{M.11}{124}{Zimmermann 恒等式 \tagR }{theorem.M.11}{}}
\newlabel{thm:appM-zimmermann-identity@cref}{{[theorem][11][2147483647,13]M.11}{[1][123][]124}}
\newlabel{cor:appM-reduction}{{M.12}{124}{化到最小正规乘积 \tagR }{theorem.M.12}{}}
\newlabel{cor:appM-reduction@cref}{{[corollary][12][2147483647,13]M.12}{[1][124][]124}}
\newlabel{ex:appM-callan-symanzik}{{M.13}{124}{$\varphi ^4$ 中的 Zimmermann 恒等式}{theorem.M.13}{}}
\newlabel{ex:appM-callan-symanzik@cref}{{[example][13][2147483647,13]M.13}{[1][124][]124}}
\newlabel{lem:appM-routing}{{M.14}{124}{冻结 $q$ 时外动量分配的相容性}{theorem.M.14}{}}
\newlabel{lem:appM-routing@cref}{{[lemma][14][2147483647,13]M.14}{[1][124][]124}}
\newlabel{sec:appM-chain}{{M.3}{124}{第 3 节（测试）}{section.M.3}{}}
\newlabel{sec:appM-chain@cref}{{[subappendix][3][2147483647,13]M.3}{[1][124][]124}}
\newlabel{lem:appM-integrability}{{M.15}{124}{冻结 $q$ 正规乘积的良定义与链分解 \tagR }{theorem.M.15}{}}
\newlabel{lem:appM-integrability@cref}{{[lemma][15][2147483647,13]M.15}{[1][124][]124}}
\newlabel{lem:appM-contact}{{M.16}{124}{接触项 \tagR }{theorem.M.16}{}}
\newlabel{lem:appM-contact@cref}{{[lemma][16][2147483647,13]M.16}{[1][124][]124}}
\newlabel{thm:appM-ope-identity}{{M.17}{124}{微扰论中的 OPE 恒等式（有质量模型） \tagR }{theorem.M.17}{}}
\newlabel{thm:appM-ope-identity@cref}{{[theorem][17][2147483647,13]M.17}{[1][124][]124}}
\newlabel{ex:appM-tree}{{M.18}{124}{树图 OPE 就是 Taylor 展开}{theorem.M.18}{}}
\newlabel{ex:appM-tree@cref}{{[example][18][2147483647,13]M.18}{[1][124][]124}}
\newlabel{rem:appM-mechanism}{{M.19}{124}{指数 $D-1-a$ 的启发式来源}{theorem.M.19}{}}
\newlabel{rem:appM-mechanism@cref}{{[remark][19][2147483647,13]M.19}{[1][124][]124}}
\newlabel{asm:appM-remainder}{{M.20}{124}{冻结 $q$ 双局域正规乘积的动量空间余项估计 \tagA }{theorem.M.20}{}}
\newlabel{asm:appM-remainder@cref}{{[assumption][20][2147483647,13]M.20}{[1][124][]124}}
\newlabel{lem:appM-derivative-bounds}{{M.21}{124}{由齐次性、O(4) 与 Lagrange 恒等式得到的导数界}{theorem.M.21}{}}
\newlabel{lem:appM-derivative-bounds@cref}{{[lemma][21][2147483647,13]M.21}{[1][124][]124}}
\newlabel{sec:appM-remainder}{{M.4}{125}{第 4 节（测试）}{section.M.4}{}}
\newlabel{sec:appM-remainder@cref}{{[subappendix][4][2147483647,13]M.4}{[1][124][]125}}
\newlabel{def:appM-blocks}{{M.22}{125}{精确投影、块与条件 (NE)}{theorem.M.22}{}}
\newlabel{def:appM-blocks@cref}{{[definition][22][2147483647,13]M.22}{[1][124][]125}}
\newlabel{lem:appM-ne}{{M.23}{125}{条件 (NE) 的适用范围}{theorem.M.23}{}}
\newlabel{lem:appM-ne@cref}{{[lemma][23][2147483647,13]M.23}{[1][125][]125}}
\newlabel{lem:appM-uniqueness}{{M.24}{125}{OPE 系数的唯一性 \tagR }{theorem.M.24}{}}
\newlabel{lem:appM-uniqueness@cref}{{[lemma][24][2147483647,13]M.24}{[1][125][]125}}
\newlabel{thm:appM-scaling}{{M.25}{125}{约化系数的标度与余项（有质量模型） \tagR }{theorem.M.25}{}}
\newlabel{thm:appM-scaling@cref}{{[theorem][25][2147483647,13]M.25}{[1][125][]125}}
\newlabel{ex:appM-unreduced}{{M.26}{125}{未约化系数违反标度（锐化值）}{theorem.M.26}{}}
\newlabel{ex:appM-unreduced@cref}{{[example][26][2147483647,13]M.26}{[1][125][]125}}
\newlabel{lem:appM-symbol}{{M.27}{125}{动量空间余项到坐标空间余项}{theorem.M.27}{}}
\newlabel{lem:appM-symbol@cref}{{[lemma][27][2147483647,13]M.27}{[1][125][]125}}
\newlabel{cor:appM-coordinate}{{M.28}{125}{坐标空间的 OPE 余项 \tagR }{theorem.M.28}{}}
\newlabel{cor:appM-coordinate@cref}{{[corollary][28][2147483647,13]M.28}{[1][125][]125}}
\newlabel{sec:appM-massless}{{M.5}{125}{第 5 节（测试）}{section.M.5}{}}
\newlabel{sec:appM-massless@cref}{{[subappendix][5][2147483647,13]M.5}{[1][125][]125}}
\newlabel{rem:appM-polynomial}{{M.29}{125}{未识别的多项式项}{theorem.M.29}{}}
\newlabel{rem:appM-polynomial@cref}{{[remark][29][2147483647,13]M.29}{[1][125][]125}}
\newlabel{rem:appM-remainder-reduction}{{M.30}{125}{余项估计化归为零动量匹配估计；唯一性论证的循环性}{theorem.M.30}{}}
\newlabel{rem:appM-remainder-reduction@cref}{{[remark][30][2147483647,13]M.30}{[1][125][]125}}
\newlabel{rem:appM-exceptional}{{M.31}{125}{例外动量与前向矩阵元}{theorem.M.31}{}}
\newlabel{rem:appM-exceptional@cref}{{[remark][31][2147483647,13]M.31}{[1][125][]125}}
\newlabel{prop:appM-zero-momentum-insertion}{{M.32}{125}{零动量扭度 2 算符插入在维数正规化中的存在性 \tagA }{theorem.M.32}{}}
\newlabel{prop:appM-zero-momentum-insertion@cref}{{[proposition][32][2147483647,13]M.32}{[1][125][]125}}
\newlabel{asm:appM-forward}{{M.33}{125}{前向在壳无质量部分子矩阵元的 OPE 余项 \tagA }{theorem.M.33}{}}
\newlabel{asm:appM-forward@cref}{{[assumption][33][2147483647,13]M.33}{[1][125][]125}}
\newlabel{def:appM-alien}{{M.34}{126}{规范不变算符、BRST 恰当算符与运动方程算符}{theorem.M.34}{}}
\newlabel{def:appM-alien@cref}{{[definition][34][2147483647,13]M.34}{[1][126][]126}}
\newlabel{thm:appM-local-cohomology}{{M.35}{126}{定域 BRST 上同调}{theorem.M.35}{}}
\newlabel{thm:appM-local-cohomology@cref}{{[theorem][35][2147483647,13]M.35}{[1][126][]126}}
\newlabel{sec:appM-brst}{{M.6}{126}{第 6 节（测试）}{section.M.6}{}}
\newlabel{sec:appM-brst@cref}{{[subappendix][6][2147483647,13]M.6}{[1][126][]126}}
\newlabel{lem:appM-lbc-formal}{{M.36}{126}{形式维数下逐圈的定域 BRST 上同调}{theorem.M.36}{}}
\newlabel{lem:appM-lbc-formal@cref}{{[lemma][36][2147483647,13]M.36}{[1][126][]126}}
\newlabel{rem:appM-joglekar-lee}{{M.37}{126}{Joglekar–Lee 混合定理的结构}{theorem.M.37}{}}
\newlabel{rem:appM-joglekar-lee@cref}{{[remark][37][2147483647,13]M.37}{[1][126][]126}}
\newlabel{lem:appM-alien-vanish}{{M.38}{126}{外来算符的物理矩阵元为零（非前向或红外正规化） \tagR }{theorem.M.38}{}}
\newlabel{lem:appM-alien-vanish@cref}{{[lemma][38][2147483647,13]M.38}{[1][126][]126}}
\newlabel{rem:appM-forward-aliens}{{M.39}{126}{前向在壳胶子矩阵元的红外微妙之处}{theorem.M.39}{}}
\newlabel{rem:appM-forward-aliens@cref}{{[remark][39][2147483647,13]M.39}{[1][126][]126}}
\newlabel{asm:appM-brst-ope}{{M.40}{126}{OPE 的 BRST 结构：带一个复合旁观者的欧氏展开 (Eε) \tagA }{theorem.M.40}{}}
\newlabel{asm:appM-brst-ope@cref}{{[assumption][40][2147483647,13]M.40}{[1][126][]126}}
\newlabel{thm:appM-brst-reduction}{{M.41}{126}{OPE 系数的 BRST 约化 \tagR }{theorem.M.41}{}}
\newlabel{thm:appM-brst-reduction@cref}{{[theorem][41][2147483647,13]M.41}{[1][126][]126}}
\newlabel{lem:appM-forward-forests}{{M.42}{126}{前向 2-森林与 Symanzik 多项式 \tagR }{theorem.M.42}{}}
\newlabel{lem:appM-forward-forests@cref}{{[lemma][42][2147483647,13]M.42}{[1][126][]126}}
\newlabel{sec:appM-forward}{{M.7}{127}{第 7 节（测试）}{section.M.7}{}}
\newlabel{sec:appM-forward@cref}{{[subappendix][7][2147483647,13]M.7}{[1][126][]127}}
\newlabel{lem:appM-newton}{{M.43}{127}{Newton 多面体控制与 Taylor 级数}{theorem.M.43}{}}
\newlabel{lem:appM-newton@cref}{{[lemma][43][2147483647,13]M.43}{[1][126][]127}}
\newlabel{lem:appM-matroid}{{M.44}{127}{拟阵形式与维数正规化的可容许性}{theorem.M.44}{}}
\newlabel{lem:appM-matroid@cref}{{[lemma][44][2147483647,13]M.44}{[1][127][]127}}
\newlabel{thm:appM-forward-expansion}{{M.45}{127}{前向在壳 DR 振幅的逐图展开定理 \tagR }{theorem.M.45}{}}
\newlabel{thm:appM-forward-expansion@cref}{{[theorem][45][2147483647,13]M.45}{[1][127][]127}}
\newlabel{ex:appM-forward-triangle}{{M.46}{127}{前向三角图：精确的 Taylor 系数}{theorem.M.46}{}}
\newlabel{ex:appM-forward-triangle@cref}{{[example][46][2147483647,13]M.46}{[1][127][]127}}
\newlabel{ex:appM-forward-nonuniform}{{M.47}{127}{离壳极限与在壳 DR 值的非一致性}{theorem.M.47}{}}
\newlabel{ex:appM-forward-nonuniform@cref}{{[example][47][2147483647,13]M.47}{[1][127][]127}}
\newlabel{lem:appM-t2-projector}{{M.48}{127}{D 维中 $T_2$ 的投影与前向振幅的全纯性 \tagR }{theorem.M.48}{}}
\newlabel{lem:appM-t2-projector@cref}{{[lemma][48][2147483647,13]M.48}{[1][127][]127}}
\newlabel{cor:appM-forward-ope}{{M.49}{127}{每阶的前向在壳部分子 OPE \tagA }{theorem.M.49}{}}
\newlabel{cor:appM-forward-ope@cref}{{[corollary][49][2147483647,13]M.49}{[1][127][]127}}
\newlabel{sec:appM-np}{{M.8}{127}{第 8 节（测试）}{section.M.8}{}}
\newlabel{sec:appM-np@cref}{{[subappendix][8][2147483647,13]M.8}{[1][127][]127}}
\newlabel{asm:appM-oswi}{{M.50}{127}{前向在壳胶子振幅的 Ward 恒等式 \tagA }{theorem.M.50}{}}
\newlabel{asm:appM-oswi@cref}{{[assumption][50][2147483647,13]M.50}{[1][127][]127}}
\newlabel{lem:appM-gluon-regularity}{{M.51}{127}{由在壳 Ward 恒等式得到胶子靶在 $\omega =0$ 的正则性 \tagR }{theorem.M.51}{}}
\newlabel{lem:appM-gluon-regularity@cref}{{[lemma][51][2147483647,13]M.51}{[1][127][]127}}
\newlabel{cor:appM-cut-plane}{{M.52}{128}{部分子前向振幅在割平面上的全纯性、不连续性与矩 \tagA }{theorem.M.52}{}}
\newlabel{cor:appM-cut-plane@cref}{{[corollary][52][2147483647,13]M.52}{[1][127][]128}}
\newlabel{ex:appM-forward-oneloop}{{M.53}{128}{单圈精确到 $\varepsilon $ 的检验}{theorem.M.53}{}}
\newlabel{ex:appM-forward-oneloop@cref}{{[example][53][2147483647,13]M.53}{[1][128][]128}}
\newlabel{lem:appM-forward-tube}{{M.54}{128}{前向管全纯性与与 $Q$ 无关的增长指数 \tagR }{theorem.M.54}{}}
\newlabel{lem:appM-forward-tube@cref}{{[lemma][54][2147483647,13]M.54}{[1][128][]128}}
\newlabel{lem:appM-hatP-zero}{{M.55}{128}{$\hat P=0$ 点的不变量论证 \tagR }{theorem.M.55}{}}
\newlabel{lem:appM-hatP-zero@cref}{{[lemma][55][2147483647,13]M.55}{[1][128][]128}}
\newlabel{lem:appM-curved-approach}{{M.56}{128}{沿弯曲路径逼近的边界值}{theorem.M.56}{}}
\newlabel{lem:appM-curved-approach@cref}{{[lemma][56][2147483647,13]M.56}{[1][128][]128}}
\newlabel{sec:appM-status}{{M.9}{128}{第 9 节（测试）}{section.M.9}{}}
\newlabel{sec:appM-status@cref}{{[subappendix][9][2147483647,13]M.9}{[1][128][]128}}
\newlabel{lem:appM-positivity}{{M.57}{128}{不连续性的正性（$Q^2$ 抹平） \tagR }{theorem.M.57}{}}
\newlabel{lem:appM-positivity@cref}{{[lemma][57][2147483647,13]M.57}{[1][128][]128}}
\newlabel{lem:appM-green}{{M.58}{128}{几乎解析延拓与 Green 公式：不连续性测度的增长}{theorem.M.58}{}}
\newlabel{lem:appM-green@cref}{{[lemma][58][2147483647,13]M.58}{[1][128][]128}}
\newlabel{lem:appM-norm-bound}{{M.59}{128}{RG 演化矩阵的范数界（质量混合的重标度） \tagR }{theorem.M.59}{}}
\newlabel{lem:appM-norm-bound@cref}{{[lemma][59][2147483647,13]M.59}{[1][128][]128}}
\newlabel{rem:appM-svz}{{M.60}{128}{实用 OPE、基的重排与重整子吸收}{theorem.M.60}{}}
\newlabel{rem:appM-svz@cref}{{[remark][60][2147483647,13]M.60}{[1][128][]128}}
\newlabel{rem:appM-supporting}{{M.61}{128}{支持性的严格结果}{theorem.M.61}{}}
\newlabel{rem:appM-supporting@cref}{{[remark][61][2147483647,13]M.61}{[1][128][]128}}
\newlabel{rem:appM-instantons}{{M.62}{128}{瞬子与渐近性}{theorem.M.62}{}}
\newlabel{rem:appM-instantons@cref}{{[remark][62][2147483647,13]M.62}{[1][128][]128}}
\newlabel{rem:appM-vacuum}{{M.63}{128}{真空 OPE}{theorem.M.63}{}}
\newlabel{rem:appM-vacuum@cref}{{[remark][63][2147483647,13]M.63}{[1][128][]128}}
\newlabel{app:appN}{{N}{129}{Collins 减除法的精确形式}{appendix.N}{}}
\newlabel{app:appN@cref}{{[appendix][14][2147483647]N}{[1][129][]129}}
\newlabel{sec:appN-setting}{{N.1}{129}{第 1 节（测试）}{section.N.1}{}}
\newlabel{sec:appN-setting@cref}{{[subappendix][1][2147483647,14]N.1}{[1][129][]129}}
\newlabel{def:appN-integrand}{{N.1}{129}{被积函数、幂次参数、围道与测度}{theorem.N.1}{}}
\newlabel{def:appN-integrand@cref}{{[definition][1][2147483647,14]N.1}{[1][129][]129}}
\newlabel{lem:appN-suppressed-sets}{{N.2}{129}{幂次压低类的基本性质}{theorem.N.2}{}}
\newlabel{lem:appN-suppressed-sets@cref}{{[lemma][2][2147483647,14]N.2}{[1][129][]129}}
\newlabel{rem:appN-hard-soft-lines}{{N.3}{129}{区域集合的两条约定}{theorem.N.3}{}}
\newlabel{rem:appN-hard-soft-lines@cref}{{[remark][3][2147483647,14]N.3}{[1][129][]129}}
\newlabel{lem:appN-beta}{{N.4}{129}{邻域指数的可行选择}{theorem.N.4}{}}
\newlabel{lem:appN-beta@cref}{{[lemma][4][2147483647,14]N.4}{[1][129][]129}}
\newlabel{ex:appN-vertex-regions}{{N.5}{129}{两圈平面顶点的九个领头区域}{theorem.N.5}{}}
\newlabel{ex:appN-vertex-regions@cref}{{[example][5][2147483647,14]N.5}{[1][129][]129}}
\newlabel{asm:appN-contour}{{N.6}{129}{(GA1) 围道 \tagA }{theorem.N.6}{}}
\newlabel{asm:appN-contour@cref}{{[assumption][6][2147483647,14]N.6}{[1][129][]129}}
\newlabel{def:appN-meet-closed}{{N.7}{129}{(GA2) 对逐线交封闭的区域集合}{theorem.N.7}{}}
\newlabel{def:appN-meet-closed@cref}{{[definition][7][2147483647,14]N.7}{[1][129][]129}}
\newlabel{sec:appN-inputs}{{N.2}{129}{第 2 节（测试）}{section.N.2}{}}
\newlabel{sec:appN-inputs@cref}{{[subappendix][2][2147483647,14]N.2}{[1][129][]129}}
\newlabel{asm:appN-envelope}{{N.8}{129}{(GA3) 幂次计数包络 \tagA }{theorem.N.8}{}}
\newlabel{asm:appN-envelope@cref}{{[assumption][8][2147483647,14]N.8}{[1][129][]129}}
\newlabel{asm:appN-separation}{{N.9}{130}{(GA4) 分离性 \tagA }{theorem.N.9}{}}
\newlabel{asm:appN-separation@cref}{{[assumption][9][2147483647,14]N.9}{[1][130][]130}}
\newlabel{rem:appN-inputs-status}{{N.10}{130}{各输入的地位}{theorem.N.10}{}}
\newlabel{rem:appN-inputs-status@cref}{{[remark][10][2147483647,14]N.10}{[1][130][]130}}
\newlabel{thm:appN-subtraction}{{N.11}{130}{减除定理（条件形式） \tagR }{theorem.N.11}{}}
\newlabel{thm:appN-subtraction@cref}{{[theorem][11][2147483647,14]N.11}{[1][130][]130}}
\newlabel{def:appN-fibered}{{N.12}{130}{纤维化密度、可检验条件 (FR)(FI) 与误差类型 E-N$^\varphi $}{theorem.N.12}{}}
\newlabel{def:appN-fibered@cref}{{[definition][12][2147483647,14]N.12}{[1][130][]130}}
\newlabel{thm:appN-fibered-subtraction}{{N.13}{130}{纤维化减除定理}{theorem.N.13}{}}
\newlabel{thm:appN-fibered-subtraction@cref}{{[theorem][13][2147483647,14]N.13}{[1][130][]130}}
\newlabel{lem:appN-identification}{{N.14}{130}{纯分母乘积的纤维化–协变等同}{theorem.N.14}{}}
\newlabel{lem:appN-identification@cref}{{[lemma][14][2147483647,14]N.14}{[1][130][]130}}
\newlabel{sec:appN-theorem}{{N.3}{130}{第 3 节（测试）}{section.N.3}{}}
\newlabel{sec:appN-theorem@cref}{{[subappendix][3][2147483647,14]N.3}{[1][130][]130}}
\newlabel{lem:appN-lp-certificate}{{N.15}{130}{热带/线性规划证书}{theorem.N.15}{}}
\newlabel{lem:appN-lp-certificate@cref}{{[lemma][15][2147483647,14]N.15}{[1][130][]130}}
\newlabel{lem:appN-lc-positivity}{{N.16}{130}{光锥能量分母的定号性（标量三次模型、无质量内线） \tagR }{theorem.N.16}{}}
\newlabel{lem:appN-lc-positivity@cref}{{[lemma][16][2147483647,14]N.16}{[1][130][]130}}
\newlabel{asm:appN-lce}{{N.17}{130}{光锥序表示与协变值的等同（LCE） \tagA }{theorem.N.17}{}}
\newlabel{asm:appN-lce@cref}{{[assumption][17][2147483647,14]N.17}{[1][130][]130}}
\newlabel{cor:appN-cs3-fibered}{{N.18}{131}{红外有限图的 (CS3$^\varphi $) \tagR }{theorem.N.18}{}}
\newlabel{cor:appN-cs3-fibered@cref}{{[corollary][18][2147483647,14]N.18}{[1][130][]131}}
\newlabel{rem:appN-forest}{{N.19}{131}{与 BPHZ 森林公式的正确类比}{theorem.N.19}{}}
\newlabel{rem:appN-forest@cref}{{[remark][19][2147483647,14]N.19}{[1][131][]131}}
\newlabel{asm:appN-values}{{N.20}{131}{发散链项的取值规则}{theorem.N.20}{}}
\newlabel{asm:appN-values@cref}{{[assumption][20][2147483647,14]N.20}{[1][131][]131}}
\newlabel{asm:appN-collins}{{N.21}{131}{一般图的 Collins 条件 \tagA }{theorem.N.21}{}}
\newlabel{asm:appN-collins@cref}{{[assumption][21][2147483647,14]N.21}{[1][131][]131}}
\newlabel{sec:appN-approximators}{{N.4}{131}{第 4 节（测试）}{section.N.4}{}}
\newlabel{sec:appN-approximators@cref}{{[subappendix][4][2147483647,14]N.4}{[1][131][]131}}
\newlabel{def:appN-tmd-approx}{{N.22}{131}{TMD 近似子与快度正规化}{theorem.N.22}{}}
\newlabel{def:appN-tmd-approx@cref}{{[definition][22][2147483647,14]N.22}{[1][131][]131}}
\newlabel{def:appN-excl-approx}{{N.23}{131}{遍举近似子与分布振幅投影}{theorem.N.23}{}}
\newlabel{def:appN-excl-approx@cref}{{[definition][23][2147483647,14]N.23}{[1][131][]131}}
\newlabel{lem:appN-algebra}{{N.24}{131}{近似子的代数}{theorem.N.24}{}}
\newlabel{lem:appN-algebra@cref}{{[lemma][24][2147483647,14]N.24}{[1][131][]131}}
\newlabel{rem:appN-abs-edge}{{N.25}{131}{(Abs) 的一个未被覆盖的情形}{theorem.N.25}{}}
\newlabel{rem:appN-abs-edge@cref}{{[remark][25][2147483647,14]N.25}{[1][131][]131}}
\newlabel{rem:appN-routing}{{N.26}{131}{没有适配动量分配时}{theorem.N.26}{}}
\newlabel{rem:appN-routing@cref}{{[remark][26][2147483647,14]N.26}{[1][131][]131}}
\newlabel{lem:appN-error-table}{{N.27}{131}{误差表各行的证明 \tagR }{theorem.N.27}{}}
\newlabel{lem:appN-error-table@cref}{{[lemma][27][2147483647,14]N.27}{[1][131][]131}}
\newlabel{asm:appN-ward-product}{{N.28}{131}{Ward 恒等式的乘积结构与减除的重组 \tagA }{theorem.N.28}{}}
\newlabel{asm:appN-ward-product@cref}{{[assumption][28][2147483647,14]N.28}{[1][131][]131}}
\newlabel{sec:appN-processes}{{N.5}{132}{第 5 节（测试）}{section.N.5}{}}
\newlabel{sec:appN-processes@cref}{{[subappendix][5][2147483647,14]N.5}{[1][132][]132}}
\newlabel{asm:appN-brst-nonlocal}{{N.29}{132}{多定域光线算符的 BRST 分类 \tagC }{theorem.N.29}{}}
\newlabel{asm:appN-brst-nonlocal@cref}{{[assumption][29][2147483647,14]N.29}{[1][132][]132}}
\newlabel{asm:appN-cutsum}{{N.30}{132}{割图求和的公共表示 \tagA }{theorem.N.30}{}}
\newlabel{asm:appN-cutsum@cref}{{[assumption][30][2147483647,14]N.30}{[1][132][]132}}
\newlabel{asm:appN-uv}{{N.31}{132}{被积函数层次的紫外重整化 \tagA }{theorem.N.31}{}}
\newlabel{asm:appN-uv@cref}{{[assumption][31][2147483647,14]N.31}{[1][132][]132}}
\newlabel{asm:appN-multren}{{N.32}{132}{共线因子的乘法重整化（维数正规化路线） \tagA }{theorem.N.32}{}}
\newlabel{asm:appN-multren@cref}{{[assumption][32][2147483647,14]N.32}{[1][132][]132}}
\newlabel{asm:appN-ir-regulator}{{N.33}{132}{保持 Ward 恒等式的部分子靶红外正规化 \tagC }{theorem.N.33}{}}
\newlabel{asm:appN-ir-regulator@cref}{{[assumption][33][2147483647,14]N.33}{[1][132][]132}}
\newlabel{prop:appN-sia}{{N.34}{132}{$e^+e^-\to hX$ 的减除结构 \tagR }{theorem.N.34}{}}
\newlabel{prop:appN-sia@cref}{{[proposition][34][2147483647,14]N.34}{[1][132][]132}}
\newlabel{prop:appN-thrust-zerobin}{{N.35}{132}{推力的零箱积分无标度 \tagR }{theorem.N.35}{}}
\newlabel{prop:appN-thrust-zerobin@cref}{{[proposition][35][2147483647,14]N.35}{[1][132][]132}}
\newlabel{sec:appN-exclusive}{{N.6}{133}{第 6 节（测试）}{section.N.6}{}}
\newlabel{sec:appN-exclusive@cref}{{[subappendix][6][2147483647,14]N.6}{[1][132][]133}}
\newlabel{prop:appN-thrust}{{N.36}{133}{双喷注事件形状（推力）的减除结构 \tagR }{theorem.N.36}{}}
\newlabel{prop:appN-thrust@cref}{{[proposition][36][2147483647,14]N.36}{[1][132][]133}}
\newlabel{prop:appN-dy}{{N.37}{133}{DY 共线因子化的条件结构 \tagR }{theorem.N.37}{}}
\newlabel{prop:appN-dy@cref}{{[proposition][37][2147483647,14]N.37}{[1][133][]133}}
\newlabel{asm:appN-tmd}{{N.38}{133}{TMD 输入 \tagA }{theorem.N.38}{}}
\newlabel{asm:appN-tmd@cref}{{[assumption][38][2147483647,14]N.38}{[1][133][]133}}
\newlabel{rem:appN-spacelike-selfenergy}{{N.39}{133}{类空正规子线的自能}{theorem.N.39}{}}
\newlabel{rem:appN-spacelike-selfenergy@cref}{{[remark][39][2147483647,14]N.39}{[1][133][]133}}
\newlabel{lem:appN-absorption}{{N.40}{133}{一阶零箱吸收关系 \tagR }{theorem.N.40}{}}
\newlabel{lem:appN-absorption@cref}{{[lemma][40][2147483647,14]N.40}{[1][133][]133}}
\newlabel{prop:appN-tmd-window}{{N.41}{133}{TMD 快度窗口中的减除定理 \tagR }{theorem.N.41}{}}
\newlabel{prop:appN-tmd-window@cref}{{[proposition][41][2147483647,14]N.41}{[1][133][]133}}
\newlabel{asm:appN-endpoint}{{N.42}{133}{端点条件 \tagA }{theorem.N.42}{}}
\newlabel{asm:appN-endpoint@cref}{{[assumption][42][2147483647,14]N.42}{[1][133][]133}}
\newlabel{sec:appN-examples}{{N.7}{133}{第 7 节（测试）}{section.N.7}{}}
\newlabel{sec:appN-examples@cref}{{[subappendix][7][2147483647,14]N.7}{[1][133][]133}}
\newlabel{thm:appN-structure}{{N.43}{133}{由减除法得到的因子化结构（条件形式） \tagR }{theorem.N.43}{}}
\newlabel{thm:appN-structure@cref}{{[theorem][43][2147483647,14]N.43}{[1][133][]133}}
\newlabel{prop:appN-triangle-exact}{{N.44}{134}{单圈离壳三角图：Collins 减除的完整计算与精确余项 \tagR }{theorem.N.44}{}}
\newlabel{prop:appN-triangle-exact@cref}{{[proposition][44][2147483647,14]N.44}{[1][134][]134}}
\newlabel{prop:appN-ladder}{{N.45}{134}{两级梯形图：嵌套区域与余项的乘积结构 \tagR }{theorem.N.45}{}}
\newlabel{prop:appN-ladder@cref}{{[proposition][45][2147483647,14]N.45}{[1][134][]134}}
\newlabel{prop:appN-triangle-cs}{{N.46}{134}{单圈离壳三角图：纤维化 Collins 条件的完整验证与积分减除定理 \tagR }{theorem.N.46}{}}
\newlabel{prop:appN-triangle-cs@cref}{{[proposition][46][2147483647,14]N.46}{[1][134][]134}}
\newlabel{prop:appN-phi36-dis}{{N.47}{134}{$\varphi ^3_6$ 中 DIS 前向振幅：全部单圈图与单梯级梯形图 \tagR }{theorem.N.47}{}}
\newlabel{prop:appN-phi36-dis@cref}{{[proposition][47][2147483647,14]N.47}{[1][134][]134}}
\newlabel{prop:appN-massive-photon}{{N.48}{134}{阿贝尔有质量光子：类光近似子的精确余项与快度正规子 \tagR }{theorem.N.48}{}}
\newlabel{prop:appN-massive-photon@cref}{{[proposition][48][2147483647,14]N.48}{[1][134][]134}}
\newlabel{prop:appN-twoloop-algebra}{{N.49}{134}{两圈平面顶点：区域代数与字面约定下 (Abs) 的失效 \tagR }{theorem.N.49}{}}
\newlabel{prop:appN-twoloop-algebra@cref}{{[proposition][49][2147483647,14]N.49}{[1][134][]134}}
\newlabel{sec:appN-status}{{N.8}{134}{第 8 节（测试）}{section.N.8}{}}
\newlabel{sec:appN-status@cref}{{[subappendix][8][2147483647,14]N.8}{[1][134][]134}}
\newlabel{asm:appN-analytic-reg}{{N.50}{134}{两圈图样积分的辅助解析正规化 \tagA }{theorem.N.50}{}}
\newlabel{asm:appN-analytic-reg@cref}{{[assumption][50][2147483647,14]N.50}{[1][134][]134}}
\newlabel{prop:appN-twoloop-remainder}{{N.51}{134}{两圈平面顶点：Collins 项与区域积分的等同及积分余项 \tagR }{theorem.N.51}{}}
\newlabel{prop:appN-twoloop-remainder@cref}{{[proposition][51][2147483647,14]N.51}{[1][134][]134}}
\newlabel{rem:appN-m1-infrared}{{N.52}{135}{超可重整模型 (M1) 的红外结构：全阶例子需要的模型}{theorem.N.52}{}}
\newlabel{rem:appN-m1-infrared@cref}{{[remark][52][2147483647,14]N.52}{[1][134][]135}}
\newlabel{rem:appN-fibered-allorders}{{N.53}{135}{全阶的纤维化减除：条件形式、范围与定位}{theorem.N.53}{}}
\newlabel{rem:appN-fibered-allorders@cref}{{[remark][53][2147483647,14]N.53}{[1][135][]135}}
\newlabel{app:appO}{{O}{136}{光锥排序微扰论与 Drell–Yan 的 Glauber 相消}{appendix.O}{}}
\newlabel{app:appO@cref}{{[appendix][15][2147483647]O}{[1][136][]136}}
\newlabel{sec:appO-overview}{{O.1}{136}{第 1 节（测试）}{section.O.1}{}}
\newlabel{sec:appO-overview@cref}{{[subappendix][1][2147483647,15]O.1}{[1][136][]136}}
\newlabel{def:appO-leading-data}{{O.1}{136}{领头区域的数据}{theorem.O.1}{}}
\newlabel{def:appO-leading-data@cref}{{[definition][1][2147483647,15]O.1}{[1][136][]136}}
\newlabel{def:appO-lcopt}{{O.2}{136}{光锥时间、中间态与能量分母}{theorem.O.2}{}}
\newlabel{def:appO-lcopt@cref}{{[definition][2][2147483647,15]O.2}{[1][136][]136}}
\newlabel{lem:appO-lcopt}{{O.3}{136}{由 Feynman 积分导出光锥序微扰论（有限 $\eta _0$） \tagR }{theorem.O.3}{}}
\newlabel{lem:appO-lcopt@cref}{{[lemma][3][2147483647,15]O.3}{[1][136][]136}}
\newlabel{lem:appO-numerators}{{O.4}{136}{分子、瞬时项与瞬时圈 \tagR }{theorem.O.4}{}}
\newlabel{lem:appO-numerators@cref}{{[lemma][4][2147483647,15]O.4}{[1][136][]136}}
\newlabel{rem:appO-contact}{{O.5}{136}{与 $V_H$ 重合的接触顶点在领头幂次不出现}{theorem.O.5}{}}
\newlabel{rem:appO-contact@cref}{{[remark][5][2147483647,15]O.5}{[1][136][]136}}
\newlabel{asm:appO-lcopt}{{O.6}{136}{光锥序微扰论的适用性 \tagA }{theorem.O.6}{}}
\newlabel{asm:appO-lcopt@cref}{{[assumption][6][2147483647,15]O.6}{[1][136][]136}}
\newlabel{def:appO-orientation}{{O.7}{136}{初态与末态中间态、取向与标准流}{theorem.O.7}{}}
\newlabel{def:appO-orientation@cref}{{[definition][7][2147483647,15]O.7}{[1][136][]136}}
\newlabel{sec:appO-lcopt}{{O.2}{137}{第 2 节（测试）}{section.O.2}{}}
\newlabel{sec:appO-lcopt@cref}{{[subappendix][2][2147483647,15]O.2}{[1][136][]137}}
\newlabel{lem:appO-orientation}{{O.8}{137}{光锥序中间态的取向与 $k^-$ 极点（精确形式） \tagR }{theorem.O.8}{}}
\newlabel{lem:appO-orientation@cref}{{[lemma][8][2147483647,15]O.8}{[1][136][]137}}
\newlabel{def:appO-pinch}{{O.9}{137}{Glauber 区与 Glauber 夹挤}{theorem.O.9}{}}
\newlabel{def:appO-pinch@cref}{{[definition][9][2147483647,15]O.9}{[1][137][]137}}
\newlabel{lem:appO-glauber-pc}{{O.10}{137}{Glauber 幂次计数 \tagR }{theorem.O.10}{}}
\newlabel{lem:appO-glauber-pc@cref}{{[lemma][10][2147483647,15]O.10}{[1][137][]137}}
\newlabel{lem:appO-hard-locality}{{O.11}{137}{硬顶点的光锥定域性与无权变量 $\mathcal E$ \tagR }{theorem.O.11}{}}
\newlabel{lem:appO-hard-locality@cref}{{[lemma][11][2147483647,15]O.11}{[1][137][]137}}
\newlabel{def:appO-chain}{{O.12}{137}{末态链、割位置与约化权重}{theorem.O.12}{}}
\newlabel{def:appO-chain@cref}{{[definition][12][2147483647,15]O.12}{[1][137][]137}}
\newlabel{lem:appO-cut-position}{{O.13}{137}{割位置恒等式（含插入）}{theorem.O.13}{}}
\newlabel{lem:appO-cut-position@cref}{{[lemma][13][2147483647,15]O.13}{[1][137][]137}}
\newlabel{lem:appO-unitarity}{{O.14}{137}{光锥幺正性引理（闭时路径形式） \tagR }{theorem.O.14}{}}
\newlabel{lem:appO-unitarity@cref}{{[lemma][14][2147483647,15]O.14}{[1][137][]137}}
\newlabel{sec:appO-orientation}{{O.3}{137}{第 3 节（测试）}{section.O.3}{}}
\newlabel{sec:appO-orientation@cref}{{[subappendix][3][2147483647,15]O.3}{[1][137][]137}}
\newlabel{rem:appO-cut-sum-content}{{O.15}{137}{幺正性引理用到了什么、没有消去什么}{theorem.O.15}{}}
\newlabel{rem:appO-cut-sum-content@cref}{{[remark][15][2147483647,15]O.15}{[1][137][]137}}
\newlabel{def:appO-windows}{{O.16}{137}{可局域端点、尺度适配窗口与软角区}{theorem.O.16}{}}
\newlabel{def:appO-windows@cref}{{[definition][16][2147483647,15]O.16}{[1][137][]137}}
\newlabel{lem:appO-equal-time}{{O.17}{137}{等光锥时刻插入}{theorem.O.17}{}}
\newlabel{lem:appO-equal-time@cref}{{[lemma][17][2147483647,15]O.17}{[1][137][]137}}
\newlabel{lem:appO-insertion-values}{{O.18}{137}{光锥序中的插入：(ET1)(ET2) 与等时值 \tagR }{theorem.O.18}{}}
\newlabel{lem:appO-insertion-values@cref}{{[lemma][18][2147483647,15]O.18}{[1][137][]137}}
\newlabel{rem:appO-constant-c}{{O.19}{137}{常数 $c$ 与正规化依赖性}{theorem.O.19}{}}
\newlabel{rem:appO-constant-c@cref}{{[remark][19][2147483647,15]O.19}{[1][137][]137}}
\newlabel{prop:appO-ss-localize}{{O.20}{137}{一条旁观者–旁观者线的局域化与 Glauber 算符 \tagR }{theorem.O.20}{}}
\newlabel{prop:appO-ss-localize@cref}{{[proposition][20][2147483647,15]O.20}{[1][137][]137}}
\newlabel{def:appO-sectors}{{O.21}{138}{插入扇区、类 (SEC) 与结构条件 (EQ)(ADJ)(REAL)}{theorem.O.21}{}}
\newlabel{def:appO-sectors@cref}{{[definition][21][2147483647,15]O.21}{[1][137][]138}}
\newlabel{sec:appO-unitarity}{{O.4}{138}{第 4 节（测试）}{section.O.4}{}}
\newlabel{sec:appO-unitarity@cref}{{[subappendix][4][2147483647,15]O.4}{[1][138][]138}}
\newlabel{lem:appO-joint-localization}{{O.22}{138}{非对易插入的联合局域化：扇区规则、次序权重与误差界 \tagR }{theorem.O.22}{}}
\newlabel{lem:appO-joint-localization@cref}{{[lemma][22][2147483647,15]O.22}{[1][138][]138}}
\newlabel{lem:appO-multipole-localization}{{O.23}{138}{多极点外部因子的联合局域化 \tagR }{theorem.O.23}{}}
\newlabel{lem:appO-multipole-localization@cref}{{[lemma][23][2147483647,15]O.23}{[1][138][]138}}
\newlabel{lem:appO-quadrant-residues}{{O.24}{138}{两条 SS 线的精确 $k^-$ 留数：象限公式 \tagR }{theorem.O.24}{}}
\newlabel{lem:appO-quadrant-residues@cref}{{[lemma][24][2147483647,15]O.24}{[1][138][]138}}
\newlabel{lem:appO-preintegration-localization}{{O.25}{138}{远离壳的软线：在 $S$ 圈积分之前局域化 \tagR }{theorem.O.25}{}}
\newlabel{lem:appO-preintegration-localization@cref}{{[lemma][25][2147483647,15]O.25}{[1][138][]138}}
\newlabel{lem:appO-glauber-operator}{{O.26}{138}{Glauber 插入的非阿贝尔算符结构 \tagR }{theorem.O.26}{}}
\newlabel{lem:appO-glauber-operator@cref}{{[lemma][26][2147483647,15]O.26}{[1][138][]138}}
\newlabel{def:appO-glob-model}{{O.27}{138}{全局两时模型 $\mathbb M_{\rm glob}$、NF 域、模式与次序权重}{theorem.O.27}{}}
\newlabel{def:appO-glob-model@cref}{{[definition][27][2147483647,15]O.27}{[1][138][]138}}
\newlabel{def:appO-regulator-class}{{O.28}{138}{正规子类 $\mathfrak X$ 与双侧窗口}{theorem.O.28}{}}
\newlabel{def:appO-regulator-class@cref}{{[definition][28][2147483647,15]O.28}{[1][138][]138}}
\newlabel{sec:appO-localization}{{O.5}{138}{第 5 节（测试）}{section.O.5}{}}
\newlabel{sec:appO-localization@cref}{{[subappendix][5][2147483647,15]O.5}{[1][138][]138}}
\newlabel{lem:appO-kplus-rotation}{{O.29}{138}{$k^+$ 围道的转动}{theorem.O.29}{}}
\newlabel{lem:appO-kplus-rotation@cref}{{[lemma][29][2147483647,15]O.29}{[1][138][]138}}
\newlabel{lem:appO-ray-identity}{{O.30}{138}{射线恒等式}{theorem.O.30}{}}
\newlabel{lem:appO-ray-identity@cref}{{[lemma][30][2147483647,15]O.30}{[1][138][]138}}
\newlabel{lem:appO-two-insertion-bounds}{{O.31}{138}{转动围道上双插入喷注函数的界}{theorem.O.31}{}}
\newlabel{lem:appO-two-insertion-bounds@cref}{{[lemma][31][2147483647,15]O.31}{[1][138][]138}}
\newlabel{lem:appO-halfcoinc-colour}{{O.32}{138}{半重合对的色代数 \tagR }{theorem.O.32}{}}
\newlabel{lem:appO-halfcoinc-colour@cref}{{[lemma][32][2147483647,15]O.32}{[1][138][]138}}
\newlabel{def:appO-family}{{O.33}{139}{旁观者交换族与积分映射}{theorem.O.33}{}}
\newlabel{def:appO-family@cref}{{[definition][33][2147483647,15]O.33}{[1][138][]139}}
\newlabel{lem:appO-side-independent-windows}{{O.34}{139}{与侧分配无关的窗口消去窗口边缘项 \tagR }{theorem.O.34}{}}
\newlabel{lem:appO-side-independent-windows@cref}{{[lemma][34][2147483647,15]O.34}{[1][138][]139}}
\newlabel{prop:appO-ss-cancel}{{O.35}{139}{旁观者–旁观者 Glauber 交换在族求和后相消 \tagR }{theorem.O.35}{}}
\newlabel{prop:appO-ss-cancel@cref}{{[proposition][35][2147483647,15]O.35}{[1][139][]139}}
\newlabel{sec:appO-rewriting}{{O.6}{139}{第 6 节（测试）}{section.O.6}{}}
\newlabel{sec:appO-rewriting@cref}{{[subappendix][6][2147483647,15]O.6}{[1][139][]139}}
\newlabel{ex:appO-single-exchange}{{O.36}{139}{单条 Glauber 交换}{theorem.O.36}{}}
\newlabel{ex:appO-single-exchange@cref}{{[example][36][2147483647,15]O.36}{[1][139][]139}}
\newlabel{ex:appO-two-exchange}{{O.37}{139}{两条 Glauber 交换与 $p_{\rm opp}$}{theorem.O.37}{}}
\newlabel{ex:appO-two-exchange@cref}{{[example][37][2147483647,15]O.37}{[1][139][]139}}
\newlabel{asm:appO-adjacent-swap}{{O.38}{139}{相邻交换对称性（ADJ） \tagP }{theorem.O.38}{}}
\newlabel{asm:appO-adjacent-swap@cref}{{[assumption][38][2147483647,15]O.38}{[1][139][]139}}
\newlabel{asm:appO-real-jets}{{O.39}{139}{旁观者喷注函数的实性（REAL） \tagP }{theorem.O.39}{}}
\newlabel{asm:appO-real-jets@cref}{{[assumption][39][2147483647,15]O.39}{[1][139][]139}}
\newlabel{asm:appO-collinear-ends}{{O.40}{139}{共线端点的匹配 \tagA }{theorem.O.40}{}}
\newlabel{asm:appO-collinear-ends@cref}{{[assumption][40][2147483647,15]O.40}{[1][139][]139}}
\newlabel{prop:appO-twoexch-spectator}{{O.41}{139}{全重合 SS 对的 $O(g^4)$ 侧求和（未割线）：全局两时模型中的相消 \tagP }{theorem.O.41}{}}
\newlabel{prop:appO-twoexch-spectator@cref}{{[proposition][41][2147483647,15]O.41}{[1][139][]139}}
\newlabel{prop:appO-twoexch-regulated}{{O.42}{140}{正规化版本、端点极点与双侧正规子 \tagP }{theorem.O.42}{}}
\newlabel{prop:appO-twoexch-regulated@cref}{{[proposition][42][2147483647,15]O.42}{[1][139][]140}}
\newlabel{sec:appO-deformation}{{O.7}{140}{第 7 节（测试）}{section.O.7}{}}
\newlabel{sec:appO-deformation@cref}{{[subappendix][7][2147483647,15]O.7}{[1][140][]140}}
\newlabel{prop:appO-strict-window-shape}{{O.43}{140}{严格窗口及其对正规子形状的依赖 \tagR }{theorem.O.43}{}}
\newlabel{prop:appO-strict-window-shape@cref}{{[proposition][43][2147483647,15]O.43}{[1][140][]140}}
\newlabel{prop:appO-pergraph-weights}{{O.44}{140}{全重合对的逐图次序权重不由窗口数据决定 \tagP }{theorem.O.44}{}}
\newlabel{prop:appO-pergraph-weights@cref}{{[proposition][44][2147483647,15]O.44}{[1][140][]140}}
\newlabel{rem:appO-halfcoinc}{{O.45}{140}{半重合对：无正规子的两时模型}{theorem.O.45}{}}
\newlabel{rem:appO-halfcoinc@cref}{{[remark][45][2147483647,15]O.45}{[1][140][]140}}
\newlabel{prop:appO-halfcoinc-sensitivity}{{O.46}{140}{半重合对：$W_{12}$ 的实部与对 $k^-$ 正规子的敏感性 \tagP }{theorem.O.46}{}}
\newlabel{prop:appO-halfcoinc-sensitivity@cref}{{[proposition][46][2147483647,15]O.46}{[1][140][]140}}
\newlabel{rem:appO-cut-ss}{{O.47}{140}{被割的直接 SS 线（实发射）}{theorem.O.47}{}}
\newlabel{rem:appO-cut-ss@cref}{{[remark][47][2147483647,15]O.47}{[1][140][]140}}
\newlabel{prop:appO-g6-ladder}{{O.48}{140}{$O(g^6)$：直接线梯形图类在幺正性下不封闭 \tagA }{theorem.O.48}{}}
\newlabel{prop:appO-g6-ladder@cref}{{[proposition][48][2147483647,15]O.48}{[1][140][]140}}
\newlabel{lem:appO-carrier}{{O.49}{140}{整函数载子}{theorem.O.49}{}}
\newlabel{lem:appO-carrier@cref}{{[lemma][49][2147483647,15]O.49}{[1][140][]140}}
\newlabel{sec:appO-wilson}{{O.8}{141}{第 8 节（测试）}{section.O.8}{}}
\newlabel{sec:appO-wilson@cref}{{[subappendix][8][2147483647,15]O.8}{[1][140][]141}}
\newlabel{def:appO-weak-equivalence}{{O.50}{141}{测量相容的割图求和与弱等价}{theorem.O.50}{}}
\newlabel{def:appO-weak-equivalence@cref}{{[definition][50][2147483647,15]O.50}{[1][140][]141}}
\newlabel{def:appO-rewriting}{{O.51}{141}{改写映射 $\mathfrak r_R$}{theorem.O.51}{}}
\newlabel{def:appO-rewriting@cref}{{[definition][51][2147483647,15]O.51}{[1][141][]141}}
\newlabel{asm:appO-envelope}{{O.52}{141}{纤维积分误差的可积包络 \tagA }{theorem.O.52}{}}
\newlabel{asm:appO-envelope@cref}{{[assumption][52][2147483647,15]O.52}{[1][141][]141}}
\newlabel{asm:appO-mixing}{{O.53}{141}{Glauber–软混合与重合连接的完整图集相消 \tagA }{theorem.O.53}{}}
\newlabel{asm:appO-mixing@cref}{{[assumption][53][2147483647,15]O.53}{[1][141][]141}}
\newlabel{prop:appO-one-sided}{{O.54}{141}{割图求和后的单侧性（弱等价） \tagR }{theorem.O.54}{}}
\newlabel{prop:appO-one-sided@cref}{{[proposition][54][2147483647,15]O.54}{[1][141][]141}}
\newlabel{rem:appO-weak}{{O.55}{141}{为何只有弱等价；SA/AS 线；窗口外}{theorem.O.55}{}}
\newlabel{rem:appO-weak@cref}{{[remark][55][2147483647,15]O.55}{[1][141][]141}}
\newlabel{rem:appO-h-emission}{{O.56}{141}{H-发射端点的障碍}{theorem.O.56}{}}
\newlabel{rem:appO-h-emission@cref}{{[remark][56][2147483647,15]O.56}{[1][141][]141}}
\newlabel{sec:appO-main}{{O.9}{141}{第 9 节（测试）}{section.O.9}{}}
\newlabel{sec:appO-main@cref}{{[subappendix][9][2147483647,15]O.9}{[1][141][]141}}
\newlabel{def:appO-deformation-map}{{O.57}{141}{形变映射与 Glauber 流}{theorem.O.57}{}}
\newlabel{def:appO-deformation-map@cref}{{[definition][57][2147483647,15]O.57}{[1][141][]141}}
\newlabel{lem:appO-no-crossing}{{O.58}{141}{形变不跨越奇点 \tagR }{theorem.O.58}{}}
\newlabel{lem:appO-no-crossing@cref}{{[lemma][58][2147483647,15]O.58}{[1][141][]141}}
\newlabel{def:appO-graded}{{O.59}{141}{分级构型}{theorem.O.59}{}}
\newlabel{def:appO-graded@cref}{{[definition][59][2147483647,15]O.59}{[1][141][]141}}
\newlabel{def:appO-target}{{O.60}{142}{形变目标与相容性}{theorem.O.60}{}}
\newlabel{def:appO-target@cref}{{[definition][60][2147483647,15]O.60}{[1][141][]142}}
\newlabel{prop:appO-target}{{O.61}{142}{$\mathcal G_0$ 与分级构型的形变目标 \tagR }{theorem.O.61}{}}
\newlabel{prop:appO-target@cref}{{[proposition][61][2147483647,15]O.61}{[1][141][]142}}
\newlabel{asm:appO-deformation}{{O.62}{142}{一般构型、相容性、软角区与 H-发射端点 \tagA }{theorem.O.62}{}}
\newlabel{asm:appO-deformation@cref}{{[assumption][62][2147483647,15]O.62}{[1][142][]142}}
\newlabel{prop:appO-eikonal}{{O.63}{142}{变形围道上的程函近似 \tagR }{theorem.O.63}{}}
\newlabel{prop:appO-eikonal@cref}{{[proposition][63][2147483647,15]O.63}{[1][142][]142}}
\newlabel{sec:appO-checklist}{{O.10}{142}{第 10 节（测试）}{section.O.10}{}}
\newlabel{sec:appO-checklist@cref}{{[subappendix][10][2147483647,15]O.10}{[1][142][]142}}
\newlabel{asm:appO-rewritten-cs}{{O.64}{142}{改写形式与族求和形式的 Collins 条件 \tagA }{theorem.O.64}{}}
\newlabel{asm:appO-rewritten-cs@cref}{{[assumption][64][2147483647,15]O.64}{[1][142][]142}}
\newlabel{prop:appO-past-wilson}{{O.65}{142}{Ward 恒等式与指向过去的 Wilson 线 \tagA }{theorem.O.65}{}}
\newlabel{prop:appO-past-wilson@cref}{{[proposition][65][2147483647,15]O.65}{[1][142][]142}}
\newlabel{rem:appO-timelike}{{O.66}{142}{类时与类空的非光锥 Wilson 线}{theorem.O.66}{}}
\newlabel{rem:appO-timelike@cref}{{[remark][66][2147483647,15]O.66}{[1][142][]142}}
\newlabel{rem:appO-onshell-unitarity}{{O.67}{142}{在壳线的两粒子幺正性（$O(g^4)$）}{theorem.O.67}{}}
\newlabel{rem:appO-onshell-unitarity@cref}{{[remark][67][2147483647,15]O.67}{[1][142][]142}}
\newlabel{thm:appO-main}{{O.68}{142}{色单态单举 DY 中 Glauber 区相消的条件定理 \tagR }{theorem.O.68}{}}
\newlabel{thm:appO-main@cref}{{[theorem][68][2147483647,15]O.68}{[1][142][]142}}
\newlabel{cor:appO-interface}{{O.69}{143}{与附录 N 的接口 \tagR }{theorem.O.69}{}}
\newlabel{cor:appO-interface@cref}{{[corollary][69][2147483647,15]O.69}{[1][143][]143}}
\newlabel{prop:appO-tmd}{{O.70}{143}{TMD 版本 \tagR }{theorem.O.70}{}}
\newlabel{prop:appO-tmd@cref}{{[proposition][70][2147483647,15]O.70}{[1][143][]143}}
\newlabel{sec:appO-neighbours}{{O.11}{143}{第 11 节（测试）}{section.O.11}{}}
\newlabel{sec:appO-neighbours@cref}{{[subappendix][11][2147483647,15]O.11}{[1][143][]143}}
\newlabel{rem:appO-measurement}{{O.71}{143}{测量相容性的统一判据}{theorem.O.71}{}}
\newlabel{rem:appO-measurement@cref}{{[remark][71][2147483647,15]O.71}{[1][143][]143}}
\newlabel{rem:appO-hh}{{O.72}{143}{$pp\to h_1h_2X$ 在 (O-3) 失效}{theorem.O.72}{}}
\newlabel{rem:appO-hh@cref}{{[remark][72][2147483647,15]O.72}{[1][143][]143}}
\newlabel{rem:appO-gaps}{{O.73}{143}{快度间隙与衍射}{theorem.O.73}{}}
\newlabel{rem:appO-gaps@cref}{{[remark][73][2147483647,15]O.73}{[1][143][]143}}
\newlabel{rem:appO-dft}{{O.74}{143}{Doria–Frenkel–Taylor}{theorem.O.74}{}}
\newlabel{rem:appO-dft@cref}{{[remark][74][2147483647,15]O.74}{[1][143][]143}}
\newlabel{rem:appO-double-dy}{{O.75}{143}{双 Drell–Yan}{theorem.O.75}{}}
\newlabel{rem:appO-double-dy@cref}{{[remark][75][2147483647,15]O.75}{[1][143][]143}}
\newlabel{rem:appO-rs-dictionary}{{O.76}{143}{与 Glauber SCET 的对照}{theorem.O.76}{}}
\newlabel{rem:appO-rs-dictionary@cref}{{[remark][76][2147483647,15]O.76}{[1][143][]143}}
\newlabel{rem:appO-literature-comparison}{{O.77}{143}{与其他全阶论证的比较；TMD 版本；有争议的提议}{theorem.O.77}{}}
\newlabel{rem:appO-literature-comparison@cref}{{[remark][77][2147483647,15]O.77}{[1][143][]143}}
\newlabel{app:appQ}{{Q}{215}{经核实的文献勘误}{appendix.Q}{}}
\newlabel{app:appQ@cref}{{[appendix][17][2147483647]Q}{[1][215][]215}}
\newlabel{sec:appQ-criteria}{{Q.1}{215}{第 1 节（测试）}{section.Q.1}{}}
\newlabel{sec:appQ-criteria@cref}{{[subappendix][1][2147483647,17]Q.1}{[1][215][]215}}
\newlabel{rem:appQ-scope}{{Q.1}{215}{附录 Q 不包含什么}{theorem.Q.1}{}}
\newlabel{rem:appQ-scope@cref}{{[remark][1][2147483647,17]Q.1}{[1][215][]215}}
\newlabel{book:bibstart}{{Q.1}{216}{第 1 节（测试）}{section*.267}{}}
\newlabel{book:bibstart@cref}{{[subappendix][1][2147483647,17]Q.1}{[1][216][]216}}
